% Documento da pesquisa: Gestão Hospitalar no SUS
% Partes atuais: Análise Exploratória (seções 1 a 9) e Estimação
% (seções 10 a 13) do painel de hospitais SUS/SP.
% Linguagem adaptada para leitura por gestores (não estatísticos).
% Compilar no Overleaf com pdfLaTeX. Ver LEIAME_overleaf.txt.
% Este arquivo substitui o antigo main.tex.

\documentclass[12pt,a4paper]{article}

\usepackage[T1]{fontenc}
\usepackage[utf8]{inputenc}
\usepackage[brazilian]{babel}
% margens laterais mínimas para as imagens ocuparem quase toda a página
\usepackage[top=2cm, bottom=2cm, left=1cm, right=1cm]{geometry}
\usepackage{graphicx}
\usepackage{subcaption}
\usepackage{booktabs}
\usepackage{float}
\usepackage{amsmath}
\usepackage{caption}
\usepackage{xcolor}
\usepackage{enumitem}
\usepackage[hidelinks]{hyperref}

\captionsetup{font=small, labelfont=bf}
\captionsetup[sub]{font=footnotesize}

\graphicspath{{figuras_analise_exploratoria/}{figuras_estimacao/}{figuras_fase2/}{figuras_selecao_entrevistas/}}

% Paginação: não esticar o espaço vertical para preencher a página. Com
% as figuras fixadas em [H], isto evita os grandes buracos brancos entre
% parágrafos --- o conteúdo assenta na altura natural e a folga vai para
% o rodapé, deixando o documento visualmente mais limpo.
\raggedbottom
% Evita linhas órfãs e viúvas (uma linha solta no topo ou no rodapé).
\widowpenalty=10000
\clubpenalty=10000

% subfigura padrão: largura quase total, uma por linha, centralizada
\newcommand{\subfig}[2]{%
  \begin{subfigure}[t]{0.62\textwidth}
    \centering
    \includegraphics[width=\linewidth]{#1}
    \caption{#2}
  \end{subfigure}\par\medskip}

% bloco "Como ler": explicação curta e acessível de cada figura
\newcommand{\comoler}[1]{%
  \par\medskip
  \noindent\begingroup\small\itshape
  \textcolor{black!70}{\textbf{Como ler:} #1}%
  \endgroup\par\medskip}

% bloco "Como foi feito": explica a construção do gráfico --- o que está
% no eixo horizontal e no vertical, quais variáveis foram escolhidas e
% que transformações (escala log, winsorização) foram aplicadas.
\newcommand{\comofeito}[1]{%
  \par\medskip
  \noindent\begingroup\small\itshape
  \textcolor{black!70}{\textbf{Como foi feito:} #1}%
  \endgroup\par\medskip}

\title{Gestão Hospitalar no SUS}

\begin{document}

% ══════════════════════════════ CAPA ══════════════════════════════
\begin{titlepage}
  \centering
  \vspace*{2.5cm}
  {\LARGE\bfseries Gestão Hospitalar no SUS: o Impacto do Modelo
   Institucional e Organizacional na Prestação de Serviço\par}
  \vspace{1.2cm}
  {\Large Análise Exploratória e Estimação do Painel de Hospitais
   SUS/SP, 2015 a 2025\par}
  \vspace{3.5cm}
  {\large
   João Eduardo Pastori Garcia\\
   {\normalsize Estatístico Responsável}\\[1.4em]
   Priscilla Reinisch Perdicaris\\
   {\normalsize Pesquisadora Principal}\\[1.4em]
   Alberto Tomasi Diniz Tiefensee\\
   {\normalsize Revisor e Responsável pelo Banco de Dados}\par}
  \vfill
  {\large São Paulo\\ Julho de 2026}
\end{titlepage}

\tableofcontents
\clearpage

% Cada seção começa em página nova (\clearpage também descarrega qualquer
% figura pendente antes da quebra), para o documento paginar limpo.
\section{O painel de dados e como ler este documento}
\label{sec:painel}

Esta parte do documento apresenta a Análise Exploratória dos dados da
pesquisa: antes de qualquer conclusão sobre o efeito do modelo de
gestão, é preciso conhecer em detalhe o material com que trabalhamos.
O que se faz aqui é responder perguntas simples e fundamentais: como se
distribuem os indicadores de desempenho entre os hospitais? Quem são os
hospitais de cada categoria de gestão e eles são comparáveis entre si?
O que aconteceu nos anos da pandemia? O que os dados mostram nos cinco
hospitais que trocaram a administração Direta pela gestão por OSS? E
onde estão os problemas de qualidade de registro que precisam de
cuidado?

\subsection{O painel}

A base é o painel consolidado do projeto, com \textbf{289 hospitais
acompanhados ano a ano de 2015 a 2025, totalizando 3.179 observações}
(uma observação corresponde a um hospital em um ano). O painel é
\textbf{balanceado}: todos os 289 hospitais aparecem nos 11 anos, sem
buracos, o que permite comparar a evolução de cada um deles em bases
iguais. Esse conjunto resulta de um funil documentado de seleção que
partiu de 830 estabelecimentos da base bruta do SIH/SUS e reteve
apenas hospitais com porte mínimo (mais de 50 leitos na mediana dos
anos), produção contínua no período e classificação de complexidade
disponível, excluindo ainda os hospitais de campanha da pandemia e,
na exclusão final por incomparabilidade estrutural, 25 hospitais de
populações incomparáveis com o hospital geral agudo (pediátricos
e onco-pediátricos, crônicos e de reabilitação, e psiquiátricos).

Os indicadores analisados são os quatro oficiais do estudo, sempre nas
versões que excluem as internações por COVID: \textbf{mortalidade
geral} (proporção de internações que terminam em óbito);
\textbf{tempo médio de permanência} (TMP, em dias);
\textbf{faturamento por saída} (valor pago por internação encerrada, em
reais); e \textbf{fração de alta complexidade} (proporção das
internações classificadas como de alta complexidade). A eles se somam
as duas \textbf{taxas de ocupação} (internação e UTI), em escala
percentual.

Três regras valem para todo o documento. Primeira: a categoria de
gestão de cada hospital (Direta, OSS, Público Municipal, Filantrópico,
Privado) segue a classificação assistencial do SIH como definição
adotada pela equipe; a categoria \textbf{Privado reúne apenas 3
hospitais} e por isso nunca deve ser lida como retrato de um grupo,
apenas registrada. Segunda: o \textbf{porte} de cada hospital é fixo,
calculado sobre a mediana de leitos de todos os anos com os cortes
oficiais (até 50 leitos, pequeno porte; 51 a 150, médio; 151 a 500,
grande; acima de 500, especial); como o funil já excluiu os pequenos, o
painel tem \textbf{159 hospitais médios, 120 grandes e 10 especiais}, e
o rótulo anual do SIH, que oscila com o cadastro, não é usado.
Terceira: a complexidade do hospital é medida sempre pela
\textbf{complexidade Barcelona}, escore construído a partir da estrutura
física e tecnológica do hospital, em todas as análises deste documento.

\subsection{Guia rápido de leitura dos gráficos}

Alguns tipos de gráfico e de medida se repetem ao longo do documento;
vale fixá-los uma única vez, em linguagem direta.

\newpage
\textbf{Mediana e média.} A mediana é o valor do hospital do meio da
fila: metade dos hospitais está abaixo dela, metade acima. A média é a
soma dividida pelo total e é sensível a valores extremos: um único
hospital muito caro puxa a média para cima, mas quase não move a
mediana. Como nossos indicadores têm valores extremos legítimos, a
mediana é quase sempre o retrato mais fiel do hospital típico, e por
isso é a medida preferida aqui. Os \textbf{percentis} generalizam a
ideia: o percentil 90 (p90) é o valor abaixo do qual estão 90\% dos
hospitais.

\textbf{Histograma.} É a fotografia da distribuição: o eixo horizontal
traz o valor do indicador e a altura de cada barra mostra quantas
observações caem naquele trecho. Onde as barras são altas, os hospitais
se concentram; caudas compridas para a direita indicam poucos casos com
valores muito altos.

\textbf{Gráfico de violino.} É um histograma de pé e espelhado, um por
categoria: a silhueta engorda onde há mais hospitais. As linhas
internas marcam os quartis e o número impresso é a mediana do grupo.

\textbf{Escala logarítmica.} Alguns gráficos usam eixo em escala
logarítmica, em que as marcas avançam por multiplicação (100, 1.000,
10.000) em vez de soma. Isso permite enxergar ao mesmo tempo hospitais
de R\$~500 e de R\$~20.000 por saída sem que os pequenos virem um
borrão; os rótulos dos eixos mostram sempre os valores reais.

\clearpage
\section{Como se distribuem os indicadores}
\label{sec:univariada}

A Tabela~\ref{tab:univariada} resume as distribuições dos sete
indicadores e as Figuras~\ref{fig:univA} a~\ref{fig:univC} trazem os
histogramas. Antes
dos números, um esclarecimento que simplifica todo o restante do
documento: a \textbf{mortalidade ajustada é, na prática, idêntica à
mortalidade geral} neste painel (a exclusão dos óbitos fetais, maternos
e neonatais quase não altera as taxas; a concordância entre as duas é
quase total, com correlação de 1,000 após arredondamento, e as
pequenas diferenças aparecem apenas na fração de zeros, 0,9\% contra
1,0\%). Por isso ela aparece na tabela, mas as
figuras mostram apenas a mortalidade geral, e tudo o que for dito vale
igualmente para a versão ajustada.

Em palavras simples, o que os dados mostram: o hospital típico do
painel tem \textbf{mortalidade em torno de 5,6\%} (metade dos hospitais
está entre 3,7\% e 7,1\%), \textbf{permanência típica de 4,3 dias} e
\textbf{faturamento de cerca de R\$~1.034 por internação encerrada}. Mas os
extremos importam. Na mortalidade há um segundo grupo de hospitais com
taxas próximas de zero (maternidades e especializados) e \textbf{0,9\%
das observações têm mortalidade exatamente zero}. No tempo de
permanência, a distribuição é assimétrica à direita, com máximo de
24,8 dias e \textbf{sem a antiga subpopulação de longa permanência}: os
hospitais de perfil psiquiátrico ou de cuidados prolongados, que na
versão anterior do painel se acumulavam num teto de 30,5 dias, foram
removidos na consolidação e não aparecem mais nos dados. No faturamento, a cauda direita é longa: enquanto a mediana é
R\$~1.034, o máximo chega a R\$~27.888, e é por isso que a média
(R\$~1.431) fica 38\% acima da mediana. A fração de alta complexidade é
o indicador mais concentrado: \textbf{31,3\% dos hospitais ano não
realizam nenhuma internação de alta complexidade}. Na ocupação,
\textbf{15,5\% das observações de internação ultrapassam 100\%}
(hospital operando acima da capacidade cadastrada, informação real de
sobrecarga) e a UTI combina \textbf{21,9\% de zeros} (hospitais sem UTI
ativa) com um máximo impossível de 875\%, examinado na
Seção~\ref{sec:qualidade}.

\begin{table}[H]
\centering
\caption{Resumo estatístico dos sete indicadores (3.179 observações de
hospital e ano; mortalidades e alta complexidade em proporção;
ocupações em percentual). DP é o desvio padrão, medida do espalhamento
dos valores; p90 é o valor abaixo do qual estão 90\% das observações;
assimetria positiva indica cauda comprida à direita. As versões geral e
ajustada da mortalidade coincidem em quase todas as estatísticas,
divergindo de forma visível apenas na fração de zeros.}
\label{tab:univariada}
\small
\begin{tabular}{lrrrrrrr}
\toprule
Indicador & Média & DP & Mediana & p90 & Máximo & Assimetria & Zeros \\
\midrule
Mortalidade geral      & 0,0552 & 0,0301 & 0,0558 & 0,0886 & 0,2602 & 0,80 & 0,9\% \\
Mortalidade ajustada   & 0,0552 & 0,0301 & 0,0558 & 0,0886 & 0,2602 & 0,80 & 1,0\% \\
TMP (dias)             & 4,65   & 1,92   & 4,29   & 6,73   & 24,79  & 2,40 & 0,03\% \\
Faturamento por saída (R\$)  & 1.431,37 & 1.708,50 & 1.033,97 & 2.465,42 & 27.887,60 & 6,71 & 0\% \\
Fração alta complex.   & 0,0759 & 0,1563 & 0,0029 & 0,2234 & 1,0000 & 3,39 & 31,3\% \\
Ocupação internação (\%) & 70,84 & 28,17 & 72,70 & 104,26 & 236,61 & -0,01 & 0\% \\
Ocupação UTI (\%)      & 59,84  & 44,89  & 69,38  & 96,67  & 875,21 & 2,95 & 21,9\% \\
\bottomrule
\end{tabular}
\end{table}

\comofeito{cada linha é um indicador e as colunas são estatísticas
descritivas calculadas sobre as 3.179 observações do painel (cada
observação é um hospital em um ano): \textbf{Média} (soma dividida pelo
total), \textbf{DP} (desvio padrão, o espalhamento em torno da média),
\textbf{Mediana} (valor do hospital do meio da fila), \textbf{p90} (valor
abaixo do qual estão 90\% das observações), \textbf{Máximo},
\textbf{Assimetria} (positiva indica cauda comprida à direita) e
\textbf{Zeros} (percentual de observações exatamente iguais a zero).
Mortalidades e alta complexidade estão em proporção; as ocupações, em
percentual.}

Na mortalidade, excepcionalmente, a média fica \textbf{abaixo} da
mediana: aqui não é a cauda direita que domina, mas o grupo de
hospitais de baixo óbito perto de zero, que puxa a média para baixo, o
mesmo grupo que a Seção~\ref{sec:distribuicoes} identifica. Nos
indicadores de cauda direita longa (TMP, faturamento e alta complexidade)
vale a regra usual, com a média acima da mediana; nas duas ocupações a
média também fica abaixo da mediana, pelo peso dos valores baixos e,
na UTI, dos zeros.

\begin{figure}[H]
\centering
\subfig{fig_ae_01_hist_mort_all.png}{Mortalidade geral}
\subfig{fig_ae_01_hist_tmp.png}{TMP}
\caption{Distribuição da mortalidade geral e do tempo médio de
permanência (289 hospitais, 2015 a 2025). Em cada gráfico, a linha
tracejada vertical é a mediana, a pontilhada é a média e os valores de
pico e máximo estão anotados.}
\label{fig:univA}
\end{figure}

\comofeito{cada painel é um histograma construído sobre as 3.179
observações do painel (cada observação é um hospital em um ano). O eixo
horizontal traz o valor do indicador --- mortalidade geral como
proporção de internações que terminam em óbito (já excluídas as
internações por COVID) e TMP em dias --- e o eixo vertical, a densidade,
que mede o quanto as observações se concentram em cada trecho de valor.
A linha tracejada vertical marca a mediana e a pontilhada, a média.}

\comoler{cada barra mostra quantas observações caem naquele valor do
indicador. Na mortalidade, note as duas corcovas: a da esquerda, perto
de zero, são as maternidades e especializados; a principal fica perto
de 5,5\%. No TMP, quase todos os hospitais estão entre 3 e 7 dias, com uma cauda
curta que se estende até cerca de 25 dias, sem o grupo isolado de longa
permanência que o painel anterior mostrava.}

\begin{figure}[H]
\centering
\subfig{fig_ae_01_hist_custo_saida.png}{Faturamento por saída}
\subfig{fig_ae_01_hist_pct_alta_complex.png}{Fração alta complexidade}
\caption{Distribuição do faturamento por saída e da fração de alta
complexidade.}
\label{fig:univB}
\end{figure}

\comofeito{mesma construção da figura anterior. No eixo horizontal, o
faturamento por saída em reais (valor pago por internação encerrada) e a
fração de alta complexidade (proporção das internações classificadas
como de alta complexidade, entre 0 e 1); no eixo vertical, a densidade
das mesmas 3.179 observações.}

\comoler{no faturamento, a concentração fica perto de R\$~1.000 e a cauda se
estica até quase R\$~28.000, o que explica a média (pontilhada) ficar
bem à direita da mediana (tracejada). Na alta complexidade, a barra
gigante em zero mostra os 31,3\% de hospitais ano sem nenhuma
internação desse tipo.}

\begin{figure}[H]
\centering
\subfig{fig_ae_01_hist_ocupacao_internacao.png}{Ocupação de internação}
\subfig{fig_ae_01_hist_ocupacao_uti.png}{Ocupação de UTI}
\caption{Distribuição das taxas de ocupação de internação e de UTI, em
escala percentual.}
\label{fig:univC}
\end{figure}

\comofeito{histogramas das duas taxas de ocupação em escala percentual:
eixo horizontal em \%, eixo vertical em densidade, sobre as 3.179
observações. A ocupação é calculada como diárias sobre a capacidade
cadastrada de leitos SUS; a linha vertical em 100\% separa a operação
dentro da capacidade cadastrada (à esquerda) da sobrecarga (à direita).}

\comoler{a linha pontilhada vertical marca 100\%: tudo à direita dela é
operação acima da capacidade cadastrada. Na internação, isso ocorre em
15,5\% das observações. Na UTI, a barra em zero são os hospitais sem
UTI ativa, e a cauda que se arrasta até 875\% é o problema de cadastro
examinado na Seção~\ref{sec:qualidade}.}

\clearpage
\section{O formato de cada indicador}
\label{sec:distribuicoes}

Cada indicador tem uma forma própria de variar entre os hospitais, e
conhecer essa forma tem uma consequência prática para a gestão: por trás
de um histograma de aparência estranha costuma haver um fato concreto da
rede --- dois perfis de hospital convivendo lado a lado, um acúmulo de
zeros com significado, uma cauda que exige leitura em escala própria.
Esta seção reúne esses achados de formato. A justificativa estatística
formal --- qual família de distribuição descreve cada indicador e o
quanto ela se ajusta aos dados, com as tabelas de nota de ajuste e os
gráficos QQ correspondentes --- foi reunida no
Apêndice~\ref{ap:diagnosticos}, para não interromper a leitura.

\subsection{Mortalidade: uma proporção com dois grupos de hospitais}

A mortalidade é uma proporção (fica entre 0 e 1), e o molde clássico
para proporções é a distribuição \textbf{Beta}. O ajuste de uma Beta
única descreve bem a escala geral, mas erra o pico de hospitais com
mortalidade quase nula. A Figura~\ref{fig:mistura} mostra a solução e,
ao mesmo tempo, a explicação: somando \textbf{duas Betas} (uma para
cada grupo de hospitais), a curva acompanha o histograma com nitidez. O
ponto central para a gestão é que esses dois grupos
\textbf{existem de verdade no painel}: um grupo de \textbf{35 hospitais
de baixo óbito estrutural} (mortalidade mediana abaixo de 2\%; 16 são
hospitais especializados e os outros 19, hospitais gerais de baixo
volume) e os \textbf{254 demais}, hospitais gerais com mortalidade em
torno de 6\%.
A aparente anomalia do histograma é, portanto, apenas o retrato de dois
perfis assistenciais diferentes convivendo na rede. Como também existem
observações com mortalidade exatamente zero (0,9\%), o molde completo
para a mortalidade é a \textbf{Beta com um tratamento especial para os
zeros}.

\comoler{no painel da esquerda, compare as duas curvas com as barras: a
tracejada passa por cima do vale e erra os dois picos; a escura desenha
os dois picos. No painel da direita, as curvas tracejadas coloridas são
os dois grupos de hospitais (baixo óbito à esquerda, demais à direita)
e a curva preta é a soma dos dois, praticamente idêntica ao histograma:
a forma estranha da mortalidade é só a mistura de dois perfis de
hospital.}

\begin{figure}[H]
\centering
\subfig{fig_ae_12_mistura_mort_all.png}{Um molde só vs a soma de dois}
\subfig{fig_ae_12_decomposicao_mort_all.png}{Os dois grupos reais de
hospitais}
\caption{Mortalidade geral: à esquerda, a curva cinza tracejada é o
molde único (que erra o pico perto de zero) e a curva escura é a soma
de dois moldes, que acompanha o histograma; à direita, a mesma soma
reconstruída a partir dos dois grupos reais de hospitais do painel.}
\label{fig:mistura}
\end{figure}

\comofeito{os dois painéis têm no eixo horizontal a mortalidade geral
(proporção, entre 0 e 1) e no vertical a quantidade das 3.179
observações em cada faixa. À esquerda, sobre o histograma, foram desenhadas duas
curvas ajustadas: uma Beta única e a soma de duas Betas. À direita, a
soma de duas Betas foi reconstruída separando de fato os hospitais nos
dois grupos observados no painel (baixo óbito estrutural e demais), pela
mortalidade mediana de cada hospital.}

\subsection{Alta complexidade: o peso dos zeros}

A fração de alta complexidade também é uma proporção, mas com uma
particularidade dominante: \textbf{um terço das observações vale
exatamente zero} (hospitais que não realizam alta complexidade) e uma
pequena parcela vale exatamente um (hospitais que só realizam alta
complexidade). Esses dois acúmulos nas pontas exigem um tratamento
estatístico próprio --- o modelo conhecido pela sigla \textbf{ZOIB},
cuja construção e ajuste estão no Apêndice~\ref{ap:diagnosticos}. Para a
gestão, o que importa é que \textbf{os zeros não são aleatórios}
(Figura~\ref{fig:zeros}): a proporção de hospitais sem alta complexidade
cai de 80,1\% na faixa 2 de complexidade para zero na faixa 6, e de
37,5\% nos Filantrópicos para 10,7\% nas OSS. Ou seja, saber a
complexidade, o porte e a categoria do hospital já diz muito sobre a
chance de ele estar no grupo do zero.

\comoler{cada barra é a proporção de hospitais-ano daquele grupo que
não realizou nenhuma internação de alta complexidade. A queda das
barras da esquerda para a direita no painel (b) mostra que o zero é
questão de perfil e estrutura: quase todos os hospitais de faixa 2
estão no zero, e nenhum de faixa 6 está.}

\begin{figure}[H]
\centering
\subfig{fig_ae_15_zeros_pct_alta_categoria.png}{Por categoria de
gestão}
\subfig{fig_ae_15_zeros_pct_alta_faixa.png}{Por faixa de complexidade}
\caption{Percentual de observações com alta complexidade exatamente
zero, por grupo de hospitais.}
\label{fig:zeros}
\end{figure}

\comofeito{gráfico de barras. O eixo horizontal traz os grupos de
hospitais --- as categorias de gestão no painel (a) e as faixas
Barcelona de complexidade no painel (b) --- e o eixo vertical, o
percentual de observações daquele grupo em que a fração de alta
complexidade é exatamente zero. Cada barra é, portanto, a parcela de
hospitais-ano do grupo que não realizou nenhuma internação de alta
complexidade.}

\subsection{Faturamento e tempo de permanência: valores com cauda
longa}

Faturamento por saída e TMP são valores positivos com cauda à
direita\footnote{Guia dos tamanhos de amostra que
aparecem no documento: \textbf{3.179} é o painel completo;
\textbf{3.178} exclui a única observação com TMP igual a zero (Hospital
Infante D.\ Henrique, 2021), incompatível com a escala logarítmica;
\textbf{3.177} exclui as duas observações de
denominador frágil da Seção~\ref{sec:qualidade}, regra usada nos
modelos da fase de estimação; e \textbf{2.482} são as observações com
UTI ativa depois de retiradas essas duas frágeis. Cada conta foi
conferida no painel.}, e a leitura prática que interessa à gestão é uma
só: \textbf{os dois indicadores devem ser analisados em escala
logarítmica}. Nessa escala, o faturamento vira um morro quase simétrico
(Figura~\ref{fig:histlog}) e o TMP hoje se ajusta quase perfeitamente,
porque a consolidação do painel (ETAPA F) removeu os hospitais de longa
permanência que, na versão anterior, criavam um teto de 30,5 dias que
distorcia a distribuição. A comparação formal entre as famílias de
distribuição candidatas --- com a nota de ajuste de cada uma e os
gráficos correspondentes --- está reunida no
Apêndice~\ref{ap:diagnosticos} (Tabela~\ref{tab:familias}), e confirma
que faturamento e TMP são bem descritos pelo molde log logístico nessa
escala.

\comoler{na escala logarítmica, o faturamento (a) vira um morro quase
simétrico, sinal clássico de que a análise deve ser feita nessa escala.
No TMP (b), o morro principal fica entre 3 e 7 dias, com uma cauda
curta à direita e sem o segundo grupo isolado de longa permanência que
o painel anterior exibia.}

\begin{figure}[H]
\centering
\subfig{fig_ae_03_hist_log_custo_saida.png}{Faturamento por saída}
\subfig{fig_ae_03_hist_log_tmp.png}{TMP}
\caption{Distribuições de faturamento e TMP com o eixo horizontal em escala
logarítmica (os rótulos mostram os valores reais).}
\label{fig:histlog}
\end{figure}

\comofeito{são os mesmos histogramas de faturamento por saída e TMP, mas
com o eixo horizontal em \textbf{escala logarítmica} (as marcas avançam
por multiplicação, e os rótulos mostram os valores reais em reais e em
dias); o eixo vertical segue acompanhando a quantidade de observações
por faixa. A escala
log foi escolhida justamente porque esses dois indicadores têm cauda
longa à direita, que na escala linear comprime a maioria dos hospitais
num só bloco.}

\subsection{Produção: contagens muito mais variáveis que o usual}

A produção anual (número de saídas) é uma contagem. O molde mais
simples para contagens, a Poisson, supõe que a variação entre hospitais
seja comparável à média; aqui \textbf{a variação é milhares de vezes
maior que a média} (razão de 6.249 no agregado e de 102 mesmo comparando
cada hospital só consigo mesmo ao longo dos anos). O molde adequado é a
\textbf{Binomial Negativa}, que aceita essa variação extra.

\clearpage
\section{O perfil de cada categoria de gestão}
\label{sec:categorias}

O painel contém \textbf{171 hospitais Filantrópicos, 56 Públicos
Municipais, 37 sob gestão de OSS, 27 da administração Direta e 3
Privados}. Os cinco hospitais que trocaram de gestão são contados
\textbf{nas duas categorias} nas contagens por período; por isso a soma
das categorias excede em cinco o total de 289 hospitais. Nas
Figuras~\ref{fig:composicao} e seguintes, cada hospital recebe
\textbf{atribuição única pela categoria em que passou a maior parte do
período} (os quatro conversores tardios como Direta; Sorocaba, com sete
anos como OSS, como OSS). A
Tabela~\ref{tab:categorias} traz as medianas de cada grupo e as
Figuras~\ref{fig:violA} a~\ref{fig:violC} mostram a distribuição
completa.

O retrato operacional é nítido e fala direto à gestão: \textbf{os
hospitais OSS operam com ocupação mediana de 92,6\% e os Públicos
Municipais de 86,9\%, contra 65,5\% dos Filantrópicos e apenas 64,1\%
da administração Direta}. O tempo de permanência conta história
parecida: \textbf{6,01 dias na Direta}, contra 4,82 nos Municipais,
4,74 nas OSS e 3,98 nos Filantrópicos. Essas diferenças brutas refletem,
em parte, o papel que o Estado atribui a cada rede (as OSS receberam
majoritariamente hospitais gerais de médio e grande porte em regiões
pressionadas) e não devem ser lidas como ranking de gestão. Já a
mortalidade mediana bruta inverteu em relação ao painel anterior:
\textbf{a administração Direta passou a ter a mais alta (5,84\%) e as
OSS a mais baixa (4,94\%)}, com Filantrópicos (5,65\%) e Públicos
Municipais (5,74\%) entre as duas. Essa inversão vem da composição de
saída da ETAPA F --- que retirou do painel hospitais de baixa
mortalidade concentrados na Direta e nos Filantrópicos --- e está
discutida nas Seções~\ref{sec:efeitogestao} e~\ref{sec:inferenciafinal}.
É essencial não tirar conclusões de
causa dessas comparações brutas: as categorias atendem perfis
diferentes de complexidade e porte, e essas diferenças precisam ser
descontadas antes de qualquer julgamento. O faturamento real mediano por
saída (R\$ de dez/2025)
varia pouco entre as categorias públicas e filantrópicas (de R\$~1.293
nos Municipais a R\$~1.438 nas OSS). A categoria \textbf{Privado é um
mundo à parte}: faturamento mediano de R\$~16.262 e 98,4\% de alta
complexidade, perfil de centros especializados; com só 3 hospitais,
serve de registro, não de comparação.

\begin{table}[H]
\centering
\caption{Medianas por categoria de gestão, período completo
(mortalidade em proporção; ocupações em percentual; faturamento real,
deflacionado pelo IPCA para R\$ de dez/2025, conforme a regra da
Seção~\ref{sec:tendencia}).}
\label{tab:categorias}
\small
\begin{tabular}{lrrrrrr}
\toprule
Categoria & Mortalidade & TMP & Faturamento real (R\$ dez/2025) & Alta complex. &
Ocup.\ int.\ (\%) & Ocup.\ UTI (\%) \\
\midrule
Direta            & 0,0584 & 6,01 & 1.366,65  & 0,0213 & 64,1 & 58,5 \\
OSS               & 0,0494 & 4,74 & 1.438,03  & 0,0242 & 92,6 & 73,9 \\
Público Municipal & 0,0574 & 4,82 & 1.292,71  & 0,0010 & 86,9 & 69,5 \\
Filantrópico      & 0,0565 & 3,98 & 1.303,22  & 0,0018 & 65,5 & 70,2 \\
Privado           & 0,0256 & 4,24 & 16.262,04 & 0,9836 & 40,5 & 3,5  \\
\bottomrule
\end{tabular}
\end{table}

\comofeito{cada linha é uma categoria de gestão e cada coluna traz a
\textbf{mediana} daquele indicador entre todas as observações
(hospital-ano) da categoria, no período completo. Cada observação hospital-ano
entra na categoria vigente naquele ano; os cinco conversores contribuem
para as duas categorias, como nas contagens por período do texto.
Mortalidade e alta complexidade estão em proporção; as ocupações, em
percentual; o faturamento, em valores reais deflacionados pelo IPCA
para R\$ de dez/2025.}

\comoler{a silhueta engorda onde há mais hospitais; as linhas internas
marcam os quartis e o número é a mediana. Na mortalidade, os violinos
das quatro categorias públicas e filantrópicas ocupam alturas
parecidas; no TMP, o violino da Direta está visivelmente deslocado para
cima.}

\begin{figure}[H]
\centering
\subfig{fig_ae_04_violino_mort_all.png}{Mortalidade geral}
\subfig{fig_ae_04_violino_tmp.png}{TMP (escala log)}
\caption{Distribuição da mortalidade e do TMP por categoria de gestão.
Cada violino é a silhueta da distribuição do grupo; o número impresso é
a mediana. A \emph{largura} do violino é uma densidade: ela mostra onde
os hospitais daquela categoria se concentram, e não quantos por cento
deles estão em cada altura. Violinos de categorias diferentes têm
larguras comparáveis entre si, mas a largura não se lê como quantidade
de hospitais --- para isso, veja as contagens na
Figura~\ref{fig:composicao}.}
\label{fig:violA}
\end{figure}

\comofeito{gráfico de violino, um por categoria de gestão. O eixo
horizontal traz as cinco categorias institucionais e o vertical, o valor
do indicador (mortalidade em proporção; TMP em dias, aqui em escala
logarítmica pela cauda longa). Cada violino é a silhueta da distribuição
daquela categoria (engorda onde há mais hospitais), as linhas internas
marcam os quartis e o número impresso é a mediana do grupo. Cada
observação hospital-ano entra na categoria vigente naquele ano, como na
Tabela~\ref{tab:categorias}.}

\comoler{no faturamento, os quatro grupos públicos e filantrópicos são
parecidos e o Privado flutua sozinho lá em cima, uma ordem de grandeza
acima. Nos Filantrópicos, os violinos mais compridos refletem a enorme
diversidade interna do grupo, das pequenas Santas Casas aos grandes
hospitais de referência.}

\begin{figure}[H]
\centering
\subfig{fig_ae_04_violino_custo_real.png}{Faturamento real por saída,
R\$ de dez/2025 (escala log)}
\subfig{fig_ae_04_violino_pct_alta_complex.png}{Fração alta
complexidade}
\caption{Distribuição do faturamento real por saída (R\$ de dez/2025) e
da fração de alta complexidade por categoria de gestão.}
\label{fig:violB}
\end{figure}

\comofeito{mesma construção da figura anterior, agora para o faturamento
real por saída (em R\$ de dez/2025, escala logarítmica) e para a fração
de alta complexidade (proporção entre 0 e 1), com as categorias no eixo
horizontal e o indicador no vertical.}

\comoler{na ocupação de internação, veja como os violinos de OSS e
Municipal estão inteiros mais altos que os de Direta e Filantrópico: é
a diferença de regime operacional em uma imagem. Na UTI, os corpos dos
violinos são mais parecidos, mas as caudas de zeros (hospitais sem UTI)
puxam as silhuetas para baixo.}

\begin{figure}[H]
\centering
\subfig{fig_ae_04_violino_ocupacao_internacao.png}{Ocupação de
internação}
\subfig{fig_ae_04_violino_ocupacao_uti.png}{Ocupação de UTI}
\caption{Distribuição das taxas de ocupação por categoria de gestão.}
\label{fig:violC}
\end{figure}

\comofeito{mesma construção dos violinos anteriores, para as duas taxas
de ocupação (internação e UTI) em escala percentual no eixo vertical,
por categoria de gestão no eixo horizontal.}

Uma pergunta indispensável antes de comparar categorias:
\textbf{elas têm hospitais parecidos o bastante para a comparação ser
justa?} A Figura~\ref{fig:composicao} responde contando os hospitais de
cada categoria em cada faixa de complexidade e porte. A sobreposição é
boa nas faixas 3 e 4, onde estão 20 dos 26 hospitais da Direta e 29 dos
33 das OSS; mas \textbf{a Direta não tem nenhum hospital na faixa 5 e
as OSS não têm nenhum nas faixas 2 e 6}; o Privado só existe nas faixas
2 e 3; e a faixa 2 é quase toda Filantrópica (32 dos 43 hospitais). No
porte, as OSS e os Públicos Municipais não têm hospital especial. Em consequência,
\textbf{comparações entre categorias são seguras nas faixas 3 e 4 e nos
portes médio e grande}, onde todos os grupos públicos e filantrópicos
estão presentes; fora dessa região faltam hospitais para comparar.

\comoler{cada célula informa o número de hospitais daquela categoria
naquela faixa ou porte. As células com zero são o alerta: onde há zero,
não existe hospital daquele perfil naquela categoria, e portanto não há
como comparar. A mancha escura na linha dos Filantrópicos, faixas 2 e
3, mostra onde está o grosso da rede.}

\begin{figure}[H]
\centering
\subfig{fig_ae_13_composicao_categoria_faixa.png}{Categoria por faixa
de complexidade}
\subfig{fig_ae_13_composicao_categoria_porte.png}{Categoria por porte}
\caption{Quantos hospitais cada categoria tem em cada faixa de
complexidade e em cada porte (contagem de CNES; azul mais escuro
indica mais hospitais). Aqui cada hospital conta uma única vez, na
categoria em que passou a maior parte do período; por isso as células
somam 289, e não a soma por período do texto.}
\label{fig:composicao}
\end{figure}

\comofeito{mapa de calor (contagem). No painel (a), o eixo horizontal
traz as faixas Barcelona de complexidade e, no (b), os portes; o eixo
vertical traz as categorias de gestão. O número em cada célula é a
\textbf{quantidade de hospitais} (CNES) daquela categoria naquele perfil,
e a intensidade do azul acompanha essa contagem. Cada hospital é contado
uma única vez, na categoria em que passou a maior parte do período, de
modo que as células somam 289.}


\clearpage
\section{Complexidade e porte: os grandes divisores da rede}
\label{sec:complexidadeporte}

Antes de qualquer comparação entre modelos de gestão, é preciso
reconhecer o quanto a complexidade e o porte, sozinhos, separam os
hospitais. As Figuras~\ref{fig:faixaA} a~\ref{fig:faixaC} mostram a
mediana de cada indicador ao longo das faixas Barcelona de complexidade
(da faixa 2, menor, à 6, maior), e as Figuras~\ref{fig:porteA}
a~\ref{fig:porteC}, ao longo dos portes.

O padrão é de escada: o faturamento mediano por saída sobe de \textbf{R\$~495
na faixa 2 para R\$~2.772 na faixa 6} (mais de cinco vezes), o tempo de
permanência vai de 3,27 para 6,30 dias, a fração de alta complexidade
sai de zero e chega a 22,5\%, e a ocupação de internação sobe de 37,5\%
para a casa dos 90\%. A mortalidade cresce com a complexidade até a
faixa 5 (de 2,88\% para 6,40\%) e \textbf{recua na faixa 6 (5,43\%)},
um sinal de que os hospitais de ponta convertem estrutura em resultado.
Vale notar que metade dos 49 hospitais da faixa 2 não opera UTI. No
porte, a mesma escada: faturamento mediano de \textbf{R\$~834 nos médios,
R\$~1.470 nos grandes e R\$~2.754 nos especiais}.

A conclusão para todo o restante do documento: \textbf{complexidade e
porte precisam ser sempre levados em conta ao comparar categorias de
gestão}, sob pena de atribuir ao modelo de gestão diferenças que são,
na verdade, de perfil assistencial. E, pela salvaguarda definida na
Seção~\ref{sec:painel}, a complexidade entra nas análises de
mortalidade apenas pela versão estrutural do escore.

\begin{figure}[H]
\centering
\subfig{fig_ae_05_faixa_mort_all.png}{Mortalidade geral}
\subfig{fig_ae_05_faixa_tmp.png}{TMP}
\caption{Mediana da mortalidade e do TMP por faixa Barcelona de
complexidade (2 a 6), com o valor sobre cada barra.}
\label{fig:faixaA}
\end{figure}

\comofeito{gráfico de barras. O eixo horizontal traz as faixas Barcelona
de complexidade (da 2, menor, à 6, maior) e o eixo vertical, a
\textbf{mediana} do indicador entre todos os hospitais ano daquela faixa
(mortalidade em proporção; TMP em dias), com o valor impresso sobre cada
barra. A faixa Barcelona é a métrica de complexidade construída no
projeto; nas análises de mortalidade usa-se sempre a versão estrutural
do escore.}

\comoler{cada barra é a mediana do indicador naquela faixa de
complexidade, com o valor impresso em cima. Na mortalidade, repare no
degrau que sobe até a faixa 5 e desce na 6; no TMP, a escada é sempre
ascendente.}

\begin{figure}[H]
\centering
\subfig{fig_ae_05_faixa_custo_saida.png}{Faturamento por saída}
\subfig{fig_ae_05_faixa_pct_alta_complex.png}{Fração alta complexidade}
\caption{Mediana do faturamento por saída e da fração de alta complexidade
por faixa Barcelona.}
\label{fig:faixaB}
\end{figure}

\comofeito{mesma construção da figura anterior, agora com a mediana do
faturamento por saída (em reais) e da fração de alta complexidade
(proporção) no eixo vertical, ao longo das faixas Barcelona no eixo
horizontal.}

\comoler{a escada do faturamento é a mais íngreme do documento: da faixa 2
para a 6, a mediana mais que quintuplica. A alta complexidade parte do
zero absoluto nas faixas baixas.}

\begin{figure}[H]
\centering
\subfig{fig_ae_05_faixa_ocupacao_internacao.png}{Ocupação de
internação}
\subfig{fig_ae_05_faixa_ocupacao_uti.png}{Ocupação de UTI}
\caption{Mediana das taxas de ocupação por faixa Barcelona.}
\label{fig:faixaC}
\end{figure}

\comofeito{mesma construção, com a mediana das duas taxas de ocupação
(percentual) no eixo vertical, por faixa Barcelona no eixo horizontal.}

\comoler{quanto mais complexo o hospital, mais cheio ele opera. Na
UTI, a barra da faixa 2 vale zero porque metade daqueles hospitais não
tem UTI.}

\begin{figure}[H]
\centering
\subfig{fig_ae_06_porte_mort_all.png}{Mortalidade geral}
\subfig{fig_ae_06_porte_tmp.png}{TMP}
\caption{Mediana da mortalidade e do TMP por porte (classificação fixa
pela mediana de leitos: 173 hospitais médios, 130 grandes e 11
especiais).}
\label{fig:porteA}
\end{figure}

\comofeito{mesma leitura das figuras por faixa, mas agora o eixo
horizontal traz o \textbf{porte} do hospital (classificação fixa pela
mediana de leitos de todos os anos: médio, grande e especial) e o
vertical, a mediana da mortalidade e do TMP.}

\comoler{mesma leitura das figuras por faixa, agora por porte:
hospitais maiores seguram o paciente por mais tempo e têm mortalidade
um pouco maior, coerente com casos mais graves.}

\begin{figure}[H]
\centering
\subfig{fig_ae_06_porte_custo_saida.png}{Faturamento por saída}
\subfig{fig_ae_06_porte_pct_alta_complex.png}{Fração alta complexidade}
\caption{Mediana do faturamento por saída e da fração de alta complexidade
por porte.}
\label{fig:porteB}
\end{figure}

\comofeito{mesma construção, com a mediana do faturamento por saída (em
reais) e da fração de alta complexidade (proporção) no eixo vertical, por
porte no eixo horizontal.}

\comoler{o faturamento mediano dos hospitais especiais é mais de três vezes o
dos médios, e a alta complexidade se concentra quase toda nos portes
grande e especial.}

\begin{figure}[H]
\centering
\subfig{fig_ae_06_porte_ocupacao_internacao.png}{Ocupação de
internação}
\subfig{fig_ae_06_porte_ocupacao_uti.png}{Ocupação de UTI}
\caption{Mediana das taxas de ocupação por porte.}
\label{fig:porteC}
\end{figure}

\comofeito{mesma construção, com a mediana das duas taxas de ocupação
(percentual) no eixo vertical, por porte no eixo horizontal.}

\comoler{hospitais maiores operam mais cheios, tanto na internação
quanto na UTI: mais um retrato de que porte e pressão assistencial
andam juntos.}

\clearpage
\section{A evolução no tempo e a marca da pandemia}
\label{sec:tendencia}

As Figuras~\ref{fig:tendA} a~\ref{fig:tendC} acompanham a mediana anual
de cada indicador de 2015 a 2025, com a faixa cinza destacando o biênio
2020 e 2021. A pandemia deixou marca em quase tudo: a mortalidade (já sem as
internações COVID) sobe de cerca de 5,3\% para \textbf{5,85\% em 2020 e
6,28\% em 2021} e só retorna ao patamar histórico em 2025, sinal de que
o biênio afetou também os pacientes não COVID. A ocupação de internação
despenca para \textbf{64,5\% em 2020} (leitos esvaziados por
cancelamentos e medo) enquanto a de UTI dispara para \textbf{100,8\% em
2021}, seu recorde na série: a mediana da rede passa acima da
capacidade cadastrada, e não apenas a cauda, o que é a assinatura
aritmética da superlotação discutida na Seção~\ref{sec:qualidade}. O tempo de permanência atravessa tudo
estável, em torno de 4,3 dias.

O faturamento pede uma leitura em duas camadas. Em valores correntes
(nominais), a mediana sobe de R\$~866 em 2015 para R\$~1.283 em 2025.
Mas parte dessa alta é só inflação. Corrigindo tudo para o poder de
compra de 2025 com o IPCA (Figura~\ref{fig:custoreal} e
Tabela~\ref{tab:custoreal}), a história muda de sinal: \textbf{em
valores de 2025, o faturamento mediano caiu de R\$~1.427 para R\$~1.283, uma
queda real de cerca de 10\% na década}. Toda a alta aparente é
inflação. Qualquer avaliação da evolução do faturamento deve, portanto, ser
feita em valores reais ou com o nível de preços de cada ano devidamente
considerado.

Por categoria (Figuras~\ref{fig:tendcatA} a~\ref{fig:tendcatC}), o
choque pandêmico é
generalizado, mas a queda de ocupação de 2020 é mais funda na Direta e
nos Filantrópicos; a fração de alta complexidade das OSS cresce de
forma consistente na segunda metade do período; e, no faturamento em
valores correntes, a Direta
parte de nível parecido com as demais categorias e passa a liderar o
faturamento mediano a partir de 2020.

\begin{table}[H]
\centering
\caption{Faturamento por saída mediano por ano: em valores correntes e em
valores de 2025 (corrigidos pelo IPCA, variação de dezembro a dezembro,
IBGE). O fator indica por quanto o valor do ano é multiplicado para
chegar a preços de 2025.}
\label{tab:custoreal}
\small
\begin{tabular}{lrrr}
\toprule
Ano & Nominal (R\$) & Valores de 2025 (R\$) & Fator IPCA \\
\midrule
2015 & 866,15   & 1.427,31 & 1,6479 \\
2016 & 879,45   & 1.363,47 & 1,5504 \\
2017 & 882,36   & 1.328,77 & 1,5059 \\
2018 & 903,01   & 1.310,72 & 1,4515 \\
2019 & 934,26   & 1.300,04 & 1,3915 \\
2020 & 1.023,02 & 1.361,99 & 1,3313 \\
2021 & 1.050,43 & 1.270,66 & 1,2097 \\
2022 & 1.176,37 & 1.345,12 & 1,1435 \\
2023 & 1.197,21 & 1.308,50 & 1,0930 \\
2024 & 1.266,82 & 1.320,79 & 1,0426 \\
2025 & 1.283,46 & 1.283,46 & 1,0000 \\
\bottomrule
\end{tabular}
\end{table}

\comofeito{cada linha é um ano. A coluna \textbf{Nominal} é o faturamento
mediano por saída em reais correntes daquele ano; a coluna \textbf{Valores
de 2025} é o mesmo valor multiplicado pelo \textbf{Fator IPCA} (a
variação acumulada do IPCA de dezembro a dezembro, IBGE, do ano até
2025), trazendo tudo para o poder de compra de 2025. Em 2025 o fator é
1,0000 por definição, então as duas colunas coincidem.}
\newpage
\comoler{a linha é o hospital mediano em cada ano; a área azul em volta
contém a metade central dos hospitais (do quartil inferior ao
superior). Dentro da faixa cinza da pandemia, note a corcova da
mortalidade; o TMP, em contraste, atravessa a década quase imóvel.}

\begin{figure}[H]
\centering
\subfig{fig_ae_07_tendencia_mort_all.png}{Mortalidade geral}
\subfig{fig_ae_07_tendencia_tmp.png}{TMP}
\caption{Mediana anual da mortalidade e do TMP, com a metade central
dos hospitais na área sombreada em azul; a faixa cinza destaca 2020 e
2021 e os pontos de máximo, mínimo e último ano estão anotados.}
\label{fig:tendA}
\end{figure}

\comofeito{série temporal. O eixo horizontal é o ano (2015 a 2025) e o
vertical, o valor do indicador (mortalidade em proporção; TMP em dias). A
linha liga a \textbf{mediana} de todos os hospitais em cada ano; a área
sombreada em azul ao redor é a metade central dos hospitais (do quartil
inferior ao superior) e a faixa cinza destaca o biênio pandêmico de 2020
e 2021. Pontos de máximo, mínimo e último ano estão anotados.}
\newpage
\comoler{o faturamento real recua ao longo da década: descontada a
inflação, o valor pago por internação encerrada perde poder de compra. O
contraste direto com a linha nominal está na
Figura~\ref{fig:custoreal}. A alta complexidade mediana cresce nos anos
finais da série.}

\begin{figure}[H]
\centering
\subfig{fig_ae_07_tendencia_custo_real.png}{Faturamento real por saída
(R\$ de dez/2025)}
\subfig{fig_ae_07_tendencia_pct_alta_complex.png}{Fração alta
complexidade}
\caption{Mediana anual do faturamento real por saída (R\$ de dez/2025) e
da fração de alta complexidade.}
\label{fig:tendB}
\end{figure}

\comofeito{mesma construção da figura anterior (mediana anual e metade
central sombreada), para o faturamento real por saída (deflacionado pelo
IPCA para R\$ de dez/2025) e a fração de alta complexidade; ano no eixo
horizontal, indicador no vertical.}
\newpage
\comoler{dentro da faixa cinza da pandemia, os dois movimentos em
sentidos opostos: a internação esvazia em 2020 (mergulho para 64,5\%) e
a UTI lota em 2021 (pico de 86,1\%, recorde da série).}

\begin{figure}[H]
\centering
\subfig{fig_ae_07_tendencia_ocupacao_internacao.png}{Ocupação de
internação}
\subfig{fig_ae_07_tendencia_ocupacao_uti.png}{Ocupação de UTI}
\caption{Mediana anual das taxas de ocupação.}
\label{fig:tendC}
\end{figure}

\comofeito{mesma construção, para as duas taxas de ocupação (percentual)
no eixo vertical ao longo dos anos no eixo horizontal.}
\newpage
\comoler{as duas linhas contam histórias opostas. A clara (valores
correntes) sobe o tempo todo; a escura (valores de 2025) mostra o que o
dinheiro de fato compra e cai ao longo da década: de R\$~1.427 para
R\$~1.283 por internação. Quando as duas linhas se encontram, em 2025,
é porque ali o valor corrente e o corrigido coincidem por definição.}

\begin{figure}[H]
\centering
\includegraphics[width=.62\textwidth]{fig_ae_07_tendencia_custo_saida_real.png}
\caption{Faturamento por saída mediano: a linha clara tracejada é o valor
corrente e a linha escura é o valor corrigido pelo IPCA para o poder de
compra de 2025.}
\label{fig:custoreal}
\end{figure}

\comofeito{o eixo horizontal é o ano e o vertical, o faturamento mediano
por saída em reais. São desenhadas \textbf{duas linhas}: a clara
tracejada é o valor corrente (nominal) e a escura é o mesmo valor
corrigido pelo IPCA (variação de dezembro a dezembro, IBGE) para o poder
de compra de 2025. As duas coincidem em 2025 por definição, pois é o ano
de referência da correção.}
\newpage
\comoler{cada linha é o hospital mediano de uma categoria em cada ano.
Na mortalidade, todas as categorias fazem a corcova da pandemia. No
TMP, a linha da Direta corre acima das demais durante a década
inteira.}

\begin{figure}[H]
\centering
\subfig{fig_ae_08_categoria_mort_all.png}{Mortalidade geral}
\subfig{fig_ae_08_categoria_tmp.png}{TMP}
\caption{Mediana anual da mortalidade e do TMP por categoria de gestão;
cada categoria tem cor e marcador próprios, e a linha pontilhada cinza
do Privado (3 hospitais) é apenas registro.}
\label{fig:tendcatA}
\end{figure}

\comofeito{série temporal com uma linha por categoria de gestão. O eixo
horizontal é o ano e o vertical, a \textbf{mediana anual} do indicador
(mortalidade em proporção; TMP em dias) dentro de cada categoria. Cada
categoria tem cor e marcador próprios; a linha pontilhada cinza do
Privado (só 3 hospitais) aparece apenas como registro.}
\newpage
\comoler{no faturamento, as quatro categorias sobem juntas até 2019 e se
abrem a partir de 2020, com a Direta na frente. Na alta complexidade, a
linha das OSS se descola das demais na segunda metade do período.}

\begin{figure}[H]
\centering
\subfig{fig_ae_08_categoria_custo_real.png}{Faturamento real por saída,
R\$ de dez/2025 (Privado omitido por escala)}
\subfig{fig_ae_08_categoria_pct_alta_complex.png}{Fração alta
complexidade (Privado omitido por escala)}
\caption{Mediana anual do faturamento real por saída (R\$ de dez/2025) e
da fração de alta
complexidade por categoria. O Privado foi omitido destes dois painéis
porque sua escala (faturamento dez vezes maior) esconderia as demais
categorias; seus valores estão na Tabela~\ref{tab:categorias}.}
\label{fig:tendcatB}
\end{figure}

\comofeito{mesma construção da figura anterior (uma linha de mediana
anual por categoria), para o faturamento real por saída (R\$ de
dez/2025) e a fração de alta complexidade. O Privado foi omitido destes
dois painéis porque sua escala, dez vezes maior, esconderia as demais.}
\newpage
\comoler{na ocupação de internação, as duas linhas de cima (OSS e
Municipal) e as duas de baixo (Filantrópico e Direta) mantêm a
separação de quase 25 pontos de ocupação durante toda a década (as de
cima na casa dos 90\% e as de baixo na dos 65\%), com o
mergulho de 2020 visível em todas. Na UTI, o pico coletivo de 2021 é a
marca da pandemia.}

\begin{figure}[H]
\centering
\subfig{fig_ae_08_categoria_ocupacao_internacao.png}{Ocupação de
internação}
\subfig{fig_ae_08_categoria_ocupacao_uti.png}{Ocupação de UTI}
\caption{Mediana anual das taxas de ocupação por categoria de gestão.}
\label{fig:tendcatC}
\end{figure}

\comofeito{mesma construção, com uma linha de mediana anual por categoria
para as duas taxas de ocupação (percentual) no eixo vertical.}


\clearpage
\section{Os cinco hospitais que mudaram de gestão}
\label{sec:conversores}

Cinco hospitais do painel trocaram a administração Direta pela gestão
por OSS dentro do período de estudo, e eles são o material mais valioso
da pesquisa: permitem observar o mesmo hospital antes e depois da
mudança. São eles o \textbf{Conjunto Hospitalar de Sorocaba (CNES
2081695)}, convertido em 2019; o \textbf{Centro de Referência da Saúde
da Mulher, Pérola Byington (CNES 2078287)}, convertido em 2023; e três
convertidos em 2025: o \textbf{Hospital Kátia de Souza Rodrigues,
Taipas (CNES 2082225)}, o \textbf{Hospital Geral de Vila Penteado (CNES
2091755)} e o \textbf{Hospital Estadual de Presidente Prudente (CNES
2750511)}. As Figuras~\ref{fig:conv2081695}
a~\ref{fig:conv2750511b} mostram, para cada um, a trajetória de quatro
indicadores contra a referência típica das categorias Direta e OSS.

No \textbf{Conjunto Hospitalar de Sorocaba}, único caso com sete anos
de observação após a mudança, a mortalidade vinha subindo antes da
conversão, atingiu 9,9\% em 2018 e caiu de forma sustentada depois,
chegando a 7,4\% em 2022 antes de um repique recente. Comparando as
medianas antes e depois: mortalidade \textbf{14,9\% menor} (de 9,34\%
para 7,95\%), permanência \textbf{14,0\% menor} (de 8,29 para 7,13
dias) e faturamento por saída 19,3\% maior, alta parecida com a do setor como
um todo no mesmo intervalo. O hospital segue com mortalidade bem acima
da OSS típica, coerente com seu papel de grande complexo regional que
concentra casos graves.

No \textbf{Pérola Byington}, a mudança de 2023 coincide com melhora
mais expressiva: mortalidade \textbf{39,5\% menor} (de 7,55\% para
4,57\%, convergindo para o padrão das OSS), permanência 17,5\% menor e
faturamento \textbf{11,4\% menor num período em que o faturamento do setor subia};
vale para o faturamento a mesma ressalva feita adiante para a mortalidade, a
de melhora iniciada visivelmente antes da mudança formal de gestão.

Duas cautelas de leitura. Em Sorocaba, o pico de mortalidade
imediatamente antes da conversão levanta a possibilidade de a queda
posterior ser, em parte, retorno natural ao patamar anterior. No
Pérola, a melhora começa visivelmente em 2021 e 2022, antes da mudança
formal de gestão, o que sugere que outros fatores também atuavam.
Quanto aos três convertidos em 2025, há \textbf{apenas um ano de
observação após a mudança}, insuficiente para qualquer leitura de
tendência; suas figuras ficam como registro do ponto de partida.

\comoler{em cada painel, a linha preta é o hospital; as linhas
tracejadas coloridas são o hospital mediano da Direta e o da OSS, como
referência; e a linha vertical tracejada marca a entrada na gestão OSS.
Veja na mortalidade a subida até 2018, a queda sustentada depois da
linha vertical e a distância que ainda separa o hospital da OSS
típica.}

\begin{figure}[H]
\centering
\subfig{fig_ae_09_conversor_2081695_mort_all.png}{Mortalidade geral}
\subfig{fig_ae_09_conversor_2081695_tmp.png}{TMP}
\caption{Conjunto Hospitalar de Sorocaba (CNES 2081695): conversão para
OSS em 2019. Mortalidade geral e TMP.}
\label{fig:conv2081695}
\end{figure}

\comofeito{série temporal de um único hospital. O eixo horizontal é o ano
(2015 a 2025) e o vertical, o indicador (mortalidade em proporção; TMP em
dias). A linha preta é o próprio hospital; as duas linhas tracejadas
coloridas são o hospital mediano da categoria Direta e o da OSS, usados
como referência; e a linha vertical tracejada marca o ano de entrada na
gestão OSS. Assim se compara o hospital com ele mesmo antes e depois da
mudança e com as duas redes.}

\begin{figure}[H]
\centering
\subfig{fig_ae_09_conversor_2081695_custo_saida.png}{Faturamento por saída}
\subfig{fig_ae_09_conversor_2081695_pct_alta_complex.png}{Fração alta
complexidade}
\caption{Conjunto Hospitalar de Sorocaba (CNES 2081695): conversão para
OSS em 2019. Faturamento por saída e fração de alta complexidade.}
\label{fig:conv2081695b}
\end{figure}

\comofeito{mesma construção da figura anterior (hospital em preto,
medianas Direta e OSS tracejadas, linha vertical na conversão), agora
para o faturamento por saída (reais) e a fração de alta complexidade do
mesmo hospital.}

\comoler{mesma leitura da figura anterior. Na mortalidade, a queda
começa em 2021, antes da linha vertical, e depois da conversão o
hospital passa a andar colado na referência das OSS; no faturamento, a linha
preta desce enquanto as referências sobem.}

\begin{figure}[H]
\centering
\subfig{fig_ae_09_conversor_2078287_mort_all.png}{Mortalidade geral}
\subfig{fig_ae_09_conversor_2078287_tmp.png}{TMP}
\caption{Centro de Referência da Saúde da Mulher, Pérola Byington (CNES
2078287): conversão para OSS em 2023. Mortalidade geral e TMP.}
\label{fig:conv2078287}
\end{figure}

\comofeito{mesma construção da trajetória anterior, agora para o Pérola
Byington: ano no eixo horizontal, mortalidade e TMP no vertical, linha
preta do hospital contra as medianas tracejadas de Direta e OSS, e linha
vertical no ano da conversão (2023).}

\begin{figure}[H]
\centering
\subfig{fig_ae_09_conversor_2078287_custo_saida.png}{Faturamento por saída}
\subfig{fig_ae_09_conversor_2078287_pct_alta_complex.png}{Fração alta
complexidade}
\caption{Centro de Referência da Saúde da Mulher, Pérola Byington (CNES
2078287): conversão para OSS em 2023. Faturamento por saída e fração de alta complexidade.}
\label{fig:conv2078287b}
\end{figure}

\comofeito{mesma construção, para o faturamento por saída e a fração de
alta complexidade do Pérola Byington.}

\comoler{a linha vertical está no último ano do painel: tudo o que se
vê antes dela é o hospital ainda sob administração Direta. A figura
serve de ponto de partida para acompanhamento futuro, não de avaliação
da mudança.}

\begin{figure}[H]
\centering
\subfig{fig_ae_09_conversor_2082225_mort_all.png}{Mortalidade geral}
\subfig{fig_ae_09_conversor_2082225_tmp.png}{TMP}
\caption{Hospital Kátia de Souza Rodrigues, Taipas (CNES 2082225):
conversão para OSS em 2025. Mortalidade geral e TMP.}
\label{fig:conv2082225}
\end{figure}

\comofeito{mesma construção das trajetórias anteriores, para o Hospital
Kátia de Souza Rodrigues (Taipas); mortalidade e TMP no eixo vertical,
ano no horizontal. A linha vertical fica no último ano do painel (2025),
pois a conversão é desse ano: tudo à esquerda é o hospital ainda sob
administração Direta.}

\begin{figure}[H]
\centering
\subfig{fig_ae_09_conversor_2082225_custo_saida.png}{Faturamento por saída}
\subfig{fig_ae_09_conversor_2082225_pct_alta_complex.png}{Fração alta
complexidade}
\caption{Hospital Kátia de Souza Rodrigues, Taipas (CNES 2082225):
conversão para OSS em 2025. Faturamento por saída e fração de alta complexidade.}
\label{fig:conv2082225b}
\end{figure}

\comofeito{mesma construção, para o faturamento por saída e a fração de
alta complexidade do Hospital Kátia de Souza Rodrigues.}

\comoler{mesma situação do hospital anterior: um único ano após a
mudança. Note, ainda assim, a subida de mortalidade e de faturamento ao longo
da década sob gestão Direta, o contexto em que a conversão foi
decidida.}

\begin{figure}[H]
\centering
\subfig{fig_ae_09_conversor_2091755_mort_all.png}{Mortalidade geral}
\subfig{fig_ae_09_conversor_2091755_tmp.png}{TMP}
\caption{Hospital Geral de Vila Penteado (CNES 2091755): conversão para
OSS em 2025. Mortalidade geral e TMP.}
\label{fig:conv2091755}
\end{figure}

\comofeito{mesma construção, para o Hospital Geral de Vila Penteado
(conversão em 2025): mortalidade e TMP no eixo vertical, ano no
horizontal, hospital em preto contra as medianas Direta e OSS.}

\begin{figure}[H]
\centering
\subfig{fig_ae_09_conversor_2091755_custo_saida.png}{Faturamento por saída}
\subfig{fig_ae_09_conversor_2091755_pct_alta_complex.png}{Fração alta
complexidade}
\caption{Hospital Geral de Vila Penteado (CNES 2091755): conversão para
OSS em 2025. Faturamento por saída e fração de alta complexidade.}
\label{fig:conv2091755b}
\end{figure}

\comofeito{mesma construção, para o faturamento por saída e a fração de
alta complexidade do Hospital Geral de Vila Penteado.}

\comoler{este hospital tem perfil de baixo óbito (mortalidade sempre
abaixo de 1\%, típica de maternidade), e por isso sua linha preta corre
bem abaixo das referências. Também aqui há apenas um ano após a
conversão.}

\begin{figure}[H]
\centering
\subfig{fig_ae_09_conversor_2750511_mort_all.png}{Mortalidade geral}
\subfig{fig_ae_09_conversor_2750511_tmp.png}{TMP}
\caption{Hospital Estadual de Presidente Prudente (CNES 2750511):
conversão para OSS em 2025. Mortalidade geral e TMP.}
\label{fig:conv2750511}
\end{figure}

\comofeito{mesma construção, para o Hospital Estadual de Presidente
Prudente (conversão em 2025): mortalidade e TMP no eixo vertical, ano no
horizontal, hospital em preto contra as medianas Direta e OSS.}

\begin{figure}[H]
\centering
\subfig{fig_ae_09_conversor_2750511_custo_saida.png}{Faturamento por saída}
\subfig{fig_ae_09_conversor_2750511_pct_alta_complex.png}{Fração alta
complexidade}
\caption{Hospital Estadual de Presidente Prudente (CNES 2750511):
conversão para OSS em 2025. Faturamento por saída e fração de alta complexidade.}
\label{fig:conv2750511b}
\end{figure}

\comofeito{mesma construção, para o faturamento por saída e a fração de
alta complexidade do Hospital Estadual de Presidente Prudente.}


\clearpage
\section{O que anda junto com o quê}
\label{sec:correlacao}

A Figura~\ref{fig:correlacao} resume, num único mapa, o quanto cada par
de indicadores anda junto. A medida usada (correlação de Spearman) vai
de 0, quando dois indicadores não têm relação, a 1, quando crescem
sempre juntos. A mortalidade ajustada foi retirada do mapa por ser
estatisticamente idêntica à mortalidade geral (correlação de 1,000
após arredondamento entre as duas; todas as demais relações coincidem
casa a casa).

Os pares mais fortes fazem sentido imediato para a gestão:
\textbf{faturamento por saída e fração de alta complexidade (0,77)}, faturamento e
tempo de permanência (0,52), e as duas ocupações entre si (0,54). A
\textbf{complexidade estrutural conversa com tudo}: 0,56 com alta
complexidade, 0,62 com ocupação de internação, 0,58 com faturamento, 0,60 com
TMP e 0,25 com mortalidade, o retrato numérico do papel de grande
divisor que a Seção~\ref{sec:complexidadeporte} descreveu. Já a
mortalidade se associa mais à ocupação de UTI (0,36) do que à de
internação (0,22), na direção esperada: onde a casuística é mais grave,
usa se mais UTI e morre se mais.

\comoler{procure a linha de um indicador e percorra as colunas: cada
célula diz o quanto aquele par anda junto (0 sem relação, 1 relação
perfeita). A diagonal vale sempre 1 porque compara o indicador com ele
mesmo. Repare na coluna da complexidade estrutural, escura quase
inteira: é o confundidor que precisa ser descontado em qualquer
comparação.}

\begin{figure}[H]
\centering
\includegraphics[width=.62\textwidth]{fig_ae_10_correlacao_spearman.png}
\caption{Mapa de correlações entre os indicadores, as ocupações e a
complexidade estrutural. Azul mais escuro indica relação mais forte;
todos os valores estão impressos.}
\label{fig:correlacao}
\end{figure}

\comofeito{mapa de calor. Nos dois eixos aparecem os mesmos itens --- os
sete indicadores, as duas taxas de ocupação e a complexidade estrutural
--- e cada célula traz a \textbf{correlação de Spearman} entre o par
correspondente, medida que vai de 0 (sem relação de ordem) a 1 (crescem
sempre juntos). Azul mais escuro indica relação mais forte; a diagonal
vale sempre 1 porque compara cada item consigo mesmo. Escolheu-se
Spearman por captar relações de ordem sem exigir forma linear.}

A Figura~\ref{fig:variancia} e a Tabela~\ref{tab:variancia} respondem a
outra pergunta essencial: \textbf{quando um indicador varia, essa
variação vem de comparar hospitais diferentes ou do mesmo hospital
mudando ao longo dos anos?} A resposta é contundente: a variação vem
quase toda das diferenças entre hospitais. Na fração de alta complexidade,
\textbf{apenas 4,5\% da variação total acontece dentro do mesmo hospital}
ao longo da década; no faturamento são 8,3\%, no tempo de permanência
13,1\%, na mortalidade 16,0\%, na ocupação de internação 16,7\%, e só a
ocupação de UTI passa de um terço (38,2\%).
Em termos de gestão: \textbf{cada hospital é muito parecido consigo
mesmo, ano após ano, e muito diferente dos vizinhos}. Os indicadores
têm forte inércia institucional, e mudanças reais dentro de um hospital
são eventos raros e valiosos de observar, o que dá ainda mais peso aos
cinco casos de conversão da Seção~\ref{sec:conversores}.

\begin{table}[H]
\centering
\caption{De onde vem a variação de cada indicador: parcela que ocorre
dentro do mesmo hospital ao longo dos anos e parcela que vem das
diferenças entre hospitais.}
\label{tab:variancia}
\small
\begin{tabular}{lrr}
\toprule
Indicador & Dentro do hospital (\%) & Entre hospitais (\%) \\
\midrule
Fração alta complex.   & 4,5  & 95,5 \\
Faturamento por saída        & 8,3  & 91,7 \\
TMP                    & 13,1 & 86,9 \\
Mortalidade geral      & 16,0 & 84,0 \\
Mortalidade ajustada   & 16,0 & 84,0 \\
Ocupação internação    & 16,7 & 83,3 \\
Ocupação UTI           & 38,2 & 61,8 \\
\bottomrule
\end{tabular}
\end{table}

\comofeito{cada linha é um indicador e as duas colunas repartem 100\% da
sua variação total em duas parcelas: a que ocorre \textbf{dentro do mesmo
hospital} ao longo dos anos e a que vem das \textbf{diferenças entre
hospitais}. As parcelas saem da decomposição da variância de cada
indicador nessas duas componentes.}

\comoler{quanto menor a barra, mais estável o indicador dentro de cada
hospital e mais as diferenças observadas são traço permanente da
instituição. A fração de alta complexidade, com 5\%, é quase uma
impressão digital do hospital; a ocupação de UTI, com 38\%, é a que mais
responde ao tempo, em boa parte por causa da pandemia.}

\begin{figure}[H]
\centering
\includegraphics[width=.62\textwidth]{fig_ae_14_variancia_within.png}
\caption{Parcela da variação de cada indicador que acontece dentro do
mesmo hospital ao longo do tempo.}
\label{fig:variancia}
\end{figure}

\comofeito{gráfico de barras horizontais. Cada barra é um indicador
(eixo vertical) e o comprimento mostra a \textbf{parcela da variação
total daquele indicador que ocorre dentro do mesmo hospital ao longo dos
anos} (eixo horizontal, em \%); o restante, até 100\%, vem das
diferenças entre hospitais. A parcela foi obtida decompondo a variância
de cada indicador em uma componente dentro do hospital e outra entre
hospitais.}


\clearpage
\section{Qualidade dos registros e cuidados adotados}
\label{sec:qualidade}

Nenhuma base administrativa é perfeita, e parte do trabalho
exploratório é localizar exatamente onde estão os problemas para que
eles não contaminem as conclusões. A boa notícia: os problemas deste
painel são poucos, identificados e circunscritos.

\textbf{Ocupação acima de 100\%.} Na internação, 15,5\% das observações
(492) ultrapassam 100\%, e na UTI, 6,5\% (207). Em sua maioria isso é
informação verdadeira de sobrecarga (mais diárias que a capacidade
cadastrada comporta) e não deve ser apagado. O problema real está na
cauda extrema da UTI: \textbf{os 15 maiores valores pertencem todos a
2020 e 2021}, liderados pelo \textbf{Hospital Irmãos Penteado e Santa
Casa de Campinas (CNES 2022648)}, que registra 875\% de ocupação de UTI
em 2021: são 6.389 diárias divididas pelos \textbf{2 leitos SUS de UTI}
então cadastrados, pois a razão do SIH é computada sobre os leitos SUS
(o cadastro total do hospital indicava 16 leitos de UTI, que ainda
assim dariam 109\% de ocupação). A
explicação é operacional: na onda da pandemia os hospitais expandiram
leitos de UTI às pressas e o cadastro não acompanhou, produzindo
divisões por um denominador defasado. Parte dos valores acima de 100\%
tem origem em um \textbf{descasamento de calendário}: leitos que
entraram no cadastro apenas no meio do ano (caso típico das UTIs
expandidas na pandemia) convivem com diárias acumuladas em período
maior que o de cadastro, inflando a razão. Esse viés não é corrigível
com a granularidade anual do painel; a winsorização no percentil 99 e o
tratamento da UTI em duas partes o mitigam, e ele deve ser lembrado na
leitura dos anos de 2020 e 2021. O cuidado adotado: analisar a
ocupação como razão positiva, tratar a UTI em duas partes (ter ou não
UTI ativa; e, entre os que têm, a intensidade de uso) e aparar a cauda
extrema (winsorização no percentil 99: 125\% na internação e 171\% na
UTI), com atenção redobrada ao biênio pandêmico.\footnote{Duas
alternativas metodológicas ficam registradas para trabalhos futuros: a
sensibilidade da winsorização a cortes no percentil 95 e no percentil
98, e a modelagem estrutural das diárias como contagem (Binomial
Negativa) com o logaritmo dos leitos como termo de exposição, que
dispensa a razão de ocupação e o seu denominador problemático.}

\textbf{Denominadores frágeis na pandemia.} Ao excluir as internações
COVID, alguns hospitais poderiam ficar com pouquíssimas internações
restantes, tornando suas taxas instáveis. No painel isso ocorre em
\textbf{apenas 2 observações, ambas do Hospital Infante D.\ Henrique
(CNES 2097613, São José do Rio Preto)}: 13 saídas não COVID em 2020 e
somente 4 em 2021, ano em que o TMP calculado degenera para zero (o
único zero de TMP de todo o painel). As duas observações recebem
marcação de denominador frágil e ficam de fora das análises de
resultado.

\textbf{Ruído estatístico dos hospitais pequenos.} A
Figura~\ref{fig:funil} organiza essa discussão: cada ponto é um
hospital num ano, posicionado pelo volume de saídas (eixo horizontal) e
pela mortalidade (eixo vertical); as linhas mostram a faixa em que a
mortalidade deveria variar se as diferenças fossem puro acaso
estatístico, faixa que é larga para hospitais pequenos e estreita para
os grandes. O resultado: \textbf{cerca de 66\% das observações caem
fora até da faixa mais generosa} (1.033 acima e 1.051 abaixo). Isso diz
duas coisas ao mesmo tempo: as diferenças de mortalidade entre
hospitais são reais e estruturais, ligadas ao perfil assistencial, e
não flutuação de sorte; e, abaixo de mil saídas por ano, o acaso passa
a pesar de verdade, o que justifica a cautela extra com hospitais
pequenos.

\textbf{Valores extremos e completude.} Pela regra usual de detecção
(3 intervalos interquartílicos acima do terceiro quartil), há 18
observações extremas de mortalidade (acima de 17,2\%), 87 de faturamento
(acima de R\$~4.223 em valores correntes, ver
Seção~\ref{sec:tendencia}), 178 de alta complexidade e 33 de TMP. Os
extremos de permanência são poucos porque a subpopulação de longa
permanência que os concentraria saiu na exclusão dos hospitais
estruturalmente incomparáveis. Quanto à
completude, \textbf{o painel não tem nenhum valor faltante nos sete
indicadores} nem nas variáveis de estratificação; as únicas lacunas
estão em campos cadastrais (nome fantasia e especialização), sem
consequência para as análises.

\comoler{o funil é estreito à direita (hospitais grandes, onde o acaso
pesa pouco) e largo à esquerda (hospitais pequenos, onde uma morte a
mais ou a menos muda a taxa). Pontos fora das linhas curvas indicam
mortalidade que o acaso não explica: a nuvem inteira transbordando das
linhas mostra que as diferenças entre hospitais são estruturais, e a
região à esquerda de mil saídas mostra onde as taxas devem ser lidas
com reserva.}

\begin{figure}[H]
\centering
\includegraphics[width=.62\textwidth]{fig_ae_16_funnel_mortalidade.png}
\caption{Funil da mortalidade: cada ponto é um hospital num ano; a
linha horizontal é a mortalidade média da rede (5,57\%) e as linhas
curvas delimitam a faixa esperada por puro acaso estatístico (95\% e
99,8\%), conforme o volume de saídas.}
\label{fig:funil}
\end{figure}

\comofeito{gráfico de dispersão (funil). Cada ponto é um hospital em um
ano, posicionado pelo \textbf{volume de saídas} no eixo horizontal (em
escala logarítmica) e pela \textbf{mortalidade geral} no vertical. A
linha horizontal é a mortalidade média da rede (5,57\%) e as duas linhas
curvas são os limites dentro dos quais a mortalidade deveria variar se as
diferenças fossem só acaso estatístico (95\% e 99,8\%), calculados a
partir do volume de cada hospital, por isso o funil é largo à
esquerda (hospitais pequenos) e estreito à direita (grandes).}

\clearpage
\section{Como estimamos os efeitos: método e salvaguardas}
\label{sec:estrategia}

Até aqui o documento descreveu os dados. Nesta segunda parte, entramos
na fase de \textbf{estimação}: usamos modelos estatísticos para medir
quanto de cada diferença entre hospitais pode ser atribuído ao modelo
de gestão, depois de descontar complexidade, porte e ano. A ideia
central é simples: comparar hospitais como se estivessem em igualdade
de condições. Um hospital OSS de grande porte e alta complexidade não
pode ser comparado diretamente com uma Santa Casa pequena; o modelo faz
esse desconto de forma sistemática, para todos os hospitais ao mesmo
tempo. Em todas as equações, a complexidade entra pelo escore Barcelona
\textbf{estrutural} (estrutura física e tecnológica), conforme a
salvaguarda da Seção~\ref{sec:painel}.

Cada indicador recebeu a família de modelo decidida empiricamente na
primeira parte deste documento: a mortalidade usa uma regressão Beta
com inflação em zero (que acomoda tanto a escala de porcentagens quanto
os raros anos sem óbito); a fração de alta complexidade usa o modelo
ZOIB (que trata separadamente os hospitais sem nenhuma produção de alta
complexidade e os raros casos de produção exclusiva); o faturamento real por
saída e o tempo de permanência são modelados no logaritmo, comparando
cada hospital com ele mesmo ao longo do tempo (os chamados efeitos
fixos); a produção usa a Binomial Negativa, adequada para contagens com
grande dispersão; a ocupação de internação é modelada no logaritmo sem
truncar em 100\%; e a ocupação de UTI usa um modelo em duas partes,
primeiro se o hospital tem UTI ativa, depois com que intensidade a usa.
A mistura de duas Betas que a Seção~\ref{sec:distribuicoes} identificou
na mortalidade não precisa de molde próprio aqui: os dois grupos de
hospitais são absorvidos pelas covariáveis de complexidade estrutural e
perfil, e os resíduos quantílicos reunidos no Apêndice~\ref{ap:diagnosticos} (r igual
a 0,976) comprovam que a Beta única condicional basta.

Cinco salvaguardas valem para todos os modelos. Primeira: o faturamento entra
sempre \textbf{em valores reais}, deflacionado pelo IPCA para preços de
2025, porque a primeira parte mostrou que o crescimento nominal do
faturamento era, na íntegra, inflação. Segunda: as duas observações de
denominador frágil do Hospital Infante D Henrique (CNES 2097613, São
José do Rio Preto, 2020 e 2021, com 13 e 4 saídas fora da COVID) ficam
fora de todas as equações. Terceira: o tratamento especial que o
painel anterior dava aos hospitais de longa permanência deixou de ser
necessário --- a ETAPA F removeu essa subpopulação do painel inteiro,
de modo que a antiga dummy de longa permanência fica sem variação e sai
das equações de tempo de permanência onde antes entrava.\footnote{O
teto de 30,5 dias que motivava esse tratamento era uma censura
administrativa do ciclo de faturamento da AIH; com a subpopulação
removida do painel, ele não aparece mais nas distribuições
(Seção~\ref{sec:distribuicoes}) e o ajuste do tempo de permanência é
hoje quase perfeito. A regra permanece definida no código para eventual
reativação, mas não tem efeito prático no painel corrente.} Quarta: as
ocupações extremas são limitadas ao percentil 99 (125,1\% na internação
e 171,2\% na UTI), preservando a informação de sobrecarga real e
removendo o erro de cadastro documentado do biênio pandêmico. Quinta e
mais importante: \textbf{todas as margens de erro são calculadas
agrupando por hospital}, reconhecendo que 11 anos do mesmo CNES não são
11 informações independentes; sem esse agrupamento, os efeitos
pareceriam muito mais precisos do que de fato são.

Um cuidado adicional merece explicação. Como quase toda a variação dos
indicadores está \emph{entre} hospitais e não \emph{dentro} deles ao
longo do tempo (a Seção~\ref{sec:correlacao} documentou parcelas de
variação interna de 4,5\% a 38,2\%), o contraste entre gestões
pode ser lido de duas maneiras: comparando hospitais diferentes
(between) ou acompanhando os cinco hospitais que trocaram de gestão
(within). Na mortalidade usamos a especificação de Mundlak, que separa
as duas leituras dentro do mesmo modelo; no faturamento e no tempo de
permanência, os efeitos fixos de hospital fazem a leitura within pura.
E a ressalva permanente continua de pé: a categoria Privado tem apenas
3 hospitais e entra nos modelos somente para não contaminar as demais
comparações; nenhum número dela deve ser interpretado.

\clearpage
\section{Resultados por indicador}
\label{sec:resultados}

A Figura~\ref{fig:contrasteoss} resume o contraste central da pesquisa:
quanto cada indicador difere entre a gestão por OSS e a administração
Direta, já descontados complexidade, porte e ano, com as margens de
erro agrupadas por hospital.

\textbf{Uma advertência precisa acompanhar todos os números desta seção,
e ela não está nas margens de erro.} Os coeficientes estimados \emph{dentro}
do hospital (o tempo de permanência e o faturamento) são identificados por
\textbf{cinco conversores}, e a Subseção~\ref{sub:oqueforam} documenta que
nenhum dos cinco é um caso limpo: uma conversão foi decidida em resposta a
uma crise de gestão, outra coincidiu com mudança de prédio, dois hospitais
já operavam com serviços terceirizados antes da mudança formal, e os dois
únicos com janela posterior longa foram para a mesma organização social.
Os intervalos de confiança medem incerteza amostral e não sabem nada
disso.

O quadro que emerge é o de uma \textbf{rede que trabalha em ritmo mais
alto}: as OSS \textbf{atendem 43,7\% mais pacientes}, com os leitos de
internação \textbf{22,7\% mais ocupados} e a UTI \textbf{19,3\% mais
usada} entre os hospitais que a possuem, e as três diferenças são grandes
o bastante para não haver dúvida sobre a sua existência.\footnote{A
intensidade de uso de UTI é estimada com a família Gama, que acomoda a
cauda de UTIs de uso quase nulo; ver a Seção~\ref{sec:inferenciafinal}.}
É uma diferença de regime entre redes com missões distintas, e não um
efeito demonstrado da conversão de um hospital individual (ver as
Seções~\ref{sec:efeitogestao} e~\ref{sec:inferenciafinal}).

O \textbf{paciente fica 10,2\% menos tempo internado} nas OSS quando se
compara cada hospital consigo mesmo, e é isso que fecha a conta do ritmo:
\textbf{mais gente passa pelos mesmos leitos porque cada um fica menos
tempo neles}. A margem de erro sustenta a diferença, mas com cinco
hospitais convertidos ela tende a ser otimista, e o teste mais exigente da
Seção~\ref{sec:inferenciafinal} não a confirma nem a contradiz. Daí a
classificação de sugestiva.

O \textbf{faturamento real por saída} aparece 8,3\% menor, e aqui a
resposta honesta é que \textbf{os dados não decidem}: a margem vai de
23,8\% de queda a 10,4\% de alta, ou seja, comporta hospital mais barato e
hospital mais caro. Com cinco conversões, não há como cravar se as OSS
custam mais ou menos por paciente atendido. \textit{\label{ressalva:pagamento}Ressalva de leitura (regime de pagamento): hospitais OSS são custeados por contrato de gestão com orçamento global, não por reembolso de AIH: o valor registrado por saída é, para eles, produção sem consequência financeira direta; a diferença de faturamento pode refletir o próprio regime de pagamento, a variável cujo efeito se busca medir, e não uso de recursos ou qualidade assistencial. Esta ressalva vale para toda leitura de faturamento no documento e é retomada, em forma abreviada, onde o indicador reaparece.} Os efeitos de ano do mesmo modelo confirmam, agora em
ambiente controlado, a queda real do faturamento da década: 2025 fica 8,9\%
abaixo de 2015 em preços comparáveis.

Na \textbf{mortalidade}, o modelo estima uma queda \textbf{de 5,8\% para
4,6\%} depois da conversão, medida dentro do mesmo hospital, num
patamar de significância que já é limítrofe. \textbf{Esse número não deve
ser citado}: ele não sobrevive ao confronto com o hospital gêmeo da
Seção~\ref{sec:inferenciafinal}. Duas leituras acompanham o resultado e
importam para a gestão. Primeiro, \textbf{as OSS não são hospitais de
mortalidade estruturalmente baixa}: o modelo separa o que é diferença
entre hospitais do que é mudança dentro do hospital, e a parcela entre
hospitais não aponta vantagem alguma; toda a queda estimada vem da
trajetória dos cinco conversores. Segundo, \textbf{nenhum hospital OSS ou
Municipal registra um ano sequer sem óbito} no painel, contra
Filantrópicos e Privados que registram: são redes que recebem casos de
outra natureza, e é mais uma razão para não ler mortalidade bruta como
placar de gestão. Na \textbf{fração de alta complexidade} não há diferença
entre OSS e Direta por método nenhum; quem determina essa fração é a
complexidade instalada do hospital e o seu porte, como era de esperar.

A Tabela~\ref{tab:contrasteoss} traz os mesmos contrastes em número. A
figura seguinte é a leitura visual dela; quem quiser conferir os valores
que sustentam cada afirmação desta seção encontra-os aqui.

\begin{table}[H]
\centering
\caption{O que muda quando o hospital é gerido por OSS em vez da
administração Direta, já descontados complexidade, porte e ano. As linhas
estão ordenadas pela \textbf{firmeza da evidência}, e não pelo tamanho do
número: é a ordem em que um gestor deve confiar neles. A coluna
\textbf{Diferença} é a única que precisa ser lida para decidir; as três
últimas são o registro técnico que sustenta cada linha e estão explicadas
logo abaixo da tabela.}
\label{tab:contrasteoss}
\small
\setlength{\tabcolsep}{4.5pt}
\begin{tabular}{llrrrr}
\toprule
Indicador & Diferença & Evidência & Coef. & EP & $z$ \\
\midrule
\multicolumn{6}{l}{\emph{Diferença firme, medida entre as redes}}\\
Produção (saídas)            & $+43{,}7\%$ de pacientes atendidos & firme & $0{,}3629$ & 0,0895 & $4{,}06$ \\
Ocupação de internação       & $+22{,}7\%$ de leito ocupado & firme & $0{,}2044$ & 0,0779 & $2{,}62$ \\
Ocupação de UTI              & $+19{,}3\%$ de uso da UTI & firme & $0{,}1769$ & 0,0603 & $2{,}93$ \\
\midrule
\multicolumn{6}{l}{\emph{Diferença sugestiva, medida dentro do mesmo hospital}}\\
Tempo de permanência         & $-10{,}2\%$ de tempo internado & sugestiva & $-0{,}1080$ & 0,0468 & $-2{,}31$ \\
\midrule
\multicolumn{6}{l}{\emph{Sem diferença que os dados sustentem}}\\
Mortalidade (geral)          & de 5,8\% para 4,6\%, não sustentada & frágil & $-0{,}2237$ & 0,1311 & $-1{,}71$ \\
Faturamento real por saída   & $-8{,}3\%$ por paciente, dentro da margem & nenhuma & $-0{,}0866$ & 0,0948 & $-0{,}91$ \\
Fração alta complexidade     & de 8,9\% para 7,6\%, dentro da margem & nenhuma & $-0{,}2044$ & 0,2348 & $-0{,}87$ \\
\bottomrule
\end{tabular}
\end{table}

\noindent\textbf{Como ler a coluna Evidência.} \emph{Firme} quer dizer que
a diferença aparece em todos os métodos aplicados e não depende de nenhum
hospital em particular; é retrato de como as redes operam hoje.
\emph{Sugestiva} quer dizer que a diferença é medida comparando cada
hospital consigo mesmo, antes e depois da conversão, e que os cinco
conversores disponíveis bastam para apontá-la mas não para prová-la
(Seções~\ref{sec:efeitogestao} e~\ref{sec:inferenciafinal}).
\emph{Frágil} quer dizer que o número existe no modelo mas some quando o
hospital convertido é comparado a um gêmeo construído a partir dos não
convertidos, e por isso não deve ser citado
(Seção~\ref{sec:inferenciafinal}). \emph{Nenhuma} quer dizer que a
diferença medida é menor que a própria margem de erro: os dados não
distinguem as duas redes nesse indicador.

\comofeito{cada linha é um dos sete modelos principais e as colunas trazem
o contraste entre a gestão por OSS e a administração Direta.
\textbf{Diferença} traduz o resultado do modelo para a unidade que o
gestor usa. Nos indicadores de quantidade ela é a variação percentual.
Nos dois que já são percentuais, a mortalidade e a fração de alta
complexidade, o modelo estima o quanto a taxa se desloca, e a coluna traz
esse deslocamento \textbf{escrito como as duas taxas que ele liga}: a
mortalidade sai de 5,8\% e chega a 4,6\%, a fração de alta complexidade sai
de 8,9\% e chega a 7,6\%. A taxa de partida é a média observada da
administração Direta, que é a rede de referência da comparação. Escrever o
deslocamento como se fosse variação percentual seria erro grosseiro: cair
de 5,8\% para 4,6\% é um deslocamento de 1,2 na escala da taxa, mas
equivale a perder um quinto do valor. \textbf{Coef.} é o valor estimado na escala interna de cada
modelo (o \emph{logito} para as duas proporções, o logaritmo para os
demais), nenhuma delas legível diretamente, e por isso a coluna
\textbf{Diferença} existe. \textbf{EP} é a margem de erro calculada
agrupando as observações por hospital, reconhecendo que 11 anos do mesmo
CNES não são 11 informações independentes. \textbf{$z$} é o coeficiente
dividido pela margem de erro, e valores acima de 2 em módulo indicam
diferença estatisticamente firme. Nas duas proporções, $e^{coef}-1$ seria
razão de chances e não variação percentual da média, e é por isso que ali a
coluna Diferença traz o par de taxas em vez de um único número. A linha da
ocupação de UTI já traz a reestimação pela família Gama da
Seção~\ref{sec:inferenciafinal}.}
\newpage
\comoler{cada linha é um indicador; o quadrado verde é a melhor
estimativa do efeito de ser OSS em vez de Direta, e a haste horizontal
é a margem de erro. Quando a haste cruza a linha vertical do zero, como
no faturamento, não se pode afirmar diferença; quando fica inteira de um
lado, como na produção e na ocupação, a diferença é estatisticamente
firme. Valores à direita do zero indicam OSS maior; à esquerda, OSS
menor.}

\begin{figure}[H]
\centering
\includegraphics[width=.92\textwidth]{fig_est_03_contraste_oss_direta.png}
\caption{Contraste OSS vs Direta nos modelos principais, em efeito
percentual com intervalo de 95\%; a mortalidade, que tem escala
própria, aparece anotada no título. Faturamento e TMP: identificação
\emph{within} (efeitos fixos, sustentada pelos cinco conversores);
produção e ocupações: contraste predominantemente \emph{between} entre
as redes OSS e Direta, de natureza associativa.}
\label{fig:contrasteoss}
\end{figure}

\comofeito{gráfico de pontos com margem de erro. Cada linha é um
indicador (eixo vertical) e a posição do quadrado no eixo horizontal é o
\textbf{efeito estimado de ser OSS em vez de Direta}, em variação
percentual, já descontados complexidade, porte e ano; a haste horizontal
é o intervalo de confiança de 95\%, com as margens agrupadas por
hospital. Quando a haste cruza a linha vertical do zero, não se pode
afirmar diferença. A mortalidade, por ter escala própria, é anotada no
título em vez de entrar no eixo.}

\clearpage
\section{O efeito da mudança de gestão: estudo de eventos}
\label{sec:efeitogestao}

A forma mais exigente de medir o efeito da gestão OSS é acompanhar
\textbf{os cinco hospitais que de fato mudaram} (Conjunto Hospitalar de
Sorocaba, CNES 2081695; Pérola Byington, CNES 2078287; Hospital Kátia
de Souza Rodrigues de Taipas, CNES 2082225; Hospital Geral de Vila
Penteado, CNES 2091755; e Hospital Estadual de Presidente Prudente,
CNES 2750511) ano a ano em torno da conversão, comparando cada um com
sua própria história e com a evolução simultânea dos demais 284
hospitais. É o chamado \textbf{estudo de eventos}: o ano anterior à
mudança serve de marco zero, e o modelo estima quanto o hospital se
afastou do próprio padrão em cada ano antes e depois da conversão.

As Figuras~\ref{fig:eventomorttmp} e~\ref{fig:eventocustoocup} mostram
os resultados para quatro indicadores. \textbf{Em termos puramente
descritivos}, depois da conversão a
\textbf{mortalidade cai 1,30 ponto percentual já no ano da mudança},
aprofunda para 2,64 e 2,91 pontos nos dois anos seguintes e estabiliza
em 2,23 pontos na janela de três ou mais anos; a leitura causal desses
números depende do teste de tendências anteriores, examinado logo
adiante e rejeitado justamente para a mortalidade. O \textbf{tempo de
permanência} cai de forma crescente até 24,6\% na janela longa. O
\textbf{faturamento real por saída} não reage no ano da mudança, mas acumula
queda de 13,5\% três ou mais anos depois, resultado que deve ser lido
com cautela: o teste de tendências anteriores do faturamento é rejeitado e
40\% dos conversores placebo da Seção~\ref{sec:inferenciafinal}
produzem efeito igual ou maior. \textit{Vale aqui a ressalva do regime
de pagamento enunciada na Seção~\ref{sec:resultados}: para as OSS, o
faturamento por saída é produção sem consequência financeira direta.} A \textbf{ocupação de
internação} sobe após a conversão, com pico no primeiro ano.

Antes de comemorar, é preciso olhar o lado esquerdo das figuras, os
anos \emph{anteriores} à conversão. Ali mora o teste de credibilidade:
se o hospital já vinha melhorando antes de mudar de gestão, parte da
melhora posterior não pode ser atribuída à OSS. O teste formal (os
pontos à esquerda da linha tracejada serem todos estatisticamente
nulos) \textbf{passa para o tempo de permanência (p igual a 0,254) e
para a fração de alta complexidade (p igual a 0,229), é rejeitado para
a mortalidade e a ocupação (p abaixo de 0,001) e é rejeitado também
para o faturamento (p igual a 0,046), ainda que com folga menor}. A rejeição na mortalidade tem nome e
sobrenome, já antecipados na Seção~\ref{sec:conversores}: em Sorocaba
o pico de mortalidade imediatamente anterior à conversão sugere
reversão à média, e no Pérola Byington a melhora começou antes da
mudança formal.

\textbf{O achado que sobrevive a este teste é o tempo de permanência},
e ele merece destaque porque é o único candidato a efeito causal que
esta pesquisa produz. Só ele passa no teste dos anos anteriores com
folga confortável (p igual a 0,254): antes da conversão os hospitais
não vinham encurtando internações, e depois passaram a encurtar. Isso é
o que o desenho de estudo de eventos foi construído para detectar, e é
o que ele detectou. O efeito é de cerca de 10\% de queda, cresce ao
longo dos anos posteriores e reaparece em todas as especificações que
testamos, inclusive na hierárquica da Seção~\ref{sec:fase2} e no
estimador de Callaway e Sant'Anna da
Seção~\ref{sec:inferenciafinal}. Para uma administradora, é a
afirmação mais forte que este trabalho autoriza sobre o efeito de
mudar o modelo de gestão: \textbf{depois da conversão, os mesmos
hospitais passaram a girar seus leitos mais rápido}.

A ressalva vem colada e não pode ser separada dele. Esse efeito é
identificado por \textbf{cinco hospitais}, dois deles com janela
posterior utilizável, e os três testes desenhados para situações de
poucos tratados, permutação, bootstrap selvagem por cluster (uma
técnica de reamostragem que recalcula a margem de erro sorteando
hospitais inteiros milhares de vezes, apropriada quando há poucas
unidades tratadas) e controle sintético,
não o corroboram (Seção~\ref{sec:inferenciafinal}). Também não o
contradizem: com esse número de unidades, eles não distinguem um efeito
nulo de uma ausência de poder para detectá-lo. O achado é, portanto,
\textbf{sugestivo forte, não demonstrado}, e é assim, com as duas
metades juntas, que deve ser citado.

Para a mortalidade e a ocupação, os números medem uma associação
temporal real, mas não um efeito causal limpo.

\comoler{cada ponto é o desvio do hospital convertido em relação ao seu
próprio padrão, no ano indicado; a haste vertical é a margem de erro e
o losango escuro marca o ano de referência, imediatamente anterior à
mudança. À esquerda da linha tracejada o hospital ainda era Direta; à
direita, já era OSS. Efeito crível pede pontos próximos de zero à
esquerda e afastamento só à direita: é o que se vê no tempo de
permanência, mas não na mortalidade, onde os anos anteriores já
oscilam. Log pontos podem ser lidos, aproximadamente, como variação
percentual dividida por 100.}

\begin{figure}[H]
\centering
\subfig{fig_est_01_evento_mort_all.png}{Mortalidade geral, em pontos
percentuais}
\subfig{fig_est_01_evento_tmp.png}{Tempo de permanência, em log pontos
(um log ponto equivale, aproximadamente, a 1\% de variação)}
\caption{Estudo de eventos da conversão para OSS: mortalidade geral e
tempo de permanência, amostra completa.}
\label{fig:eventomorttmp}
\end{figure}

\comofeito{estudo de eventos. O eixo horizontal traz os \textbf{anos
desde a conversão para OSS} (valores negativos antes, positivos depois,
com o ano $-1$ de referência) e o vertical, o \textbf{efeito estimado} do
indicador em relação a esse ano de referência (mortalidade em pontos
percentuais; TMP em log pontos, que equivalem, aproximadamente, à
variação percentual dividida por 100). Cada ponto é o desvio médio dos
conversores naquele ano em relação ao próprio padrão, controlando ano e
hospital; a haste vertical é a margem de erro.}

\begin{figure}[H]
\centering
\subfig{fig_est_01_evento_custo_real.png}{Faturamento real por saída, em log
pontos}
\subfig{fig_est_01_evento_ocupacao_internacao_w.png}{Ocupação de
internação, em log pontos}
\caption{Estudo de eventos da conversão para OSS: faturamento real por saída
e ocupação de internação, amostra completa.}
\label{fig:eventocustoocup}
\end{figure}

\comofeito{mesma construção da figura anterior (anos desde a conversão no
eixo horizontal, efeito estimado no vertical), agora para o faturamento
real por saída e a ocupação de internação, ambos em log pontos.}

\comoler{mesma leitura das figuras anteriores. Nesta figura, nenhum ponto
posterior à conversão se firma longe do zero: a mudança de gestão não
alterou o perfil de complexidade da produção desses hospitais.}

\begin{figure}[H]
\centering
\includegraphics[width=.62\textwidth]{fig_est_01_evento_pct_alta_complex.png}
\caption{Estudo de eventos da conversão para OSS: fração de alta
complexidade, em pontos percentuais, amostra completa.}
\label{fig:eventopctalta}
\end{figure}

\comofeito{mesma construção do estudo de eventos, para a fração de alta
complexidade em pontos percentuais no eixo vertical, ao longo dos anos
desde a conversão no horizontal.}

Como exercício de robustez, o estudo foi refeito na \textbf{região de
suporte comum} (faixas Barcelona 3 e 4, portes médio e grande), onde
Direta e OSS convivem em número suficiente. Sorocaba (faixa 5) e
Presidente Prudente (faixa 2) saem dessa amostra, que fica com Pérola
Byington, Taipas e Vila Penteado. A Figura~\ref{fig:eventosuporte}
mostra o resultado para mortalidade e tempo de permanência: os efeitos
posteriores ficam ainda maiores (mortalidade 2,05 pontos menor no ano
da conversão, tempo de permanência 18,6\% menor), mas o teste dos anos
anteriores é rejeitado em todos os cinco indicadores, e com folga (p
abaixo de 0,001 em todos, com estatística F entre 17,0 e 556,1), e a janela posterior se
apoia em três hospitais, dois deles com um único ano observado após a
mudança. Esses números confirmam a direção dos achados, sem
constituir prova independente.

\comoler{mesma leitura do estudo de eventos anterior, agora só com
hospitais comparáveis entre si. Note que restam menos anos após a
conversão: dos três conversores presentes, dois mudaram em 2025.}

\begin{figure}[H]
\centering
\subfig{fig_est_02_evento_suporte_mort_all.png}{Mortalidade geral, em
pontos percentuais}
\subfig{fig_est_02_evento_suporte_tmp.png}{Tempo de permanência, em log
pontos}
\caption{Estudo de eventos na região de suporte comum (faixas 3 e 4,
portes médio e grande): mortalidade geral e tempo de permanência.}
\label{fig:eventosuporte}
\end{figure}

\comofeito{mesma construção do estudo de eventos (mortalidade e TMP), mas
restrita à \textbf{região de suporte comum} (faixas Barcelona 3 e 4 e
portes médio e grande), onde Direta e OSS convivem em número
suficiente; ficam de fora Sorocaba (faixa 5) e Presidente Prudente
(faixa 2), restando três conversores.}


\clearpage
\section{Robustez e diagnóstico dos modelos}
\label{sec:robustez}

Duas baterias de verificação fecham a fase de estimação.

\textbf{Sem o biênio pandêmico.} Todos os modelos foram reestimados
excluindo 2020 e 2021, os anos em que a casuística e os denominadores
sofreram os choques documentados na primeira
parte. A Figura~\ref{fig:robustezpandemia} compara os contrastes OSS vs
Direta nas duas amostras: as diferenças são mínimas (o tempo de
permanência passa de 10,2\% para 11,4\% de queda, a produção de 43,7\%
para 44,7\% de alta, e assim por diante). \textbf{Nenhum dos
contrastes principais depende dos anos de pandemia.}

\begin{table}[H]
\centering
\caption{Contraste OSS \emph{versus} Direta estimado no painel completo
e na amostra que exclui 2020 e 2021, nos sete modelos principais. Cada
bloco traz o \textbf{efeito} já traduzido para a unidade de leitura, ao
lado do valor bruto do modelo e da sua margem de erro. A leitura
pretendida é a comparação entre os dois blocos: quanto mais parecidos,
menos o resultado depende dos anos de pandemia.}
\label{tab:sempandemia}
\small
\setlength{\tabcolsep}{4pt}
\begin{tabular}{lrrrrrr}
\toprule
& \multicolumn{3}{c}{Painel completo} & \multicolumn{3}{c}{Sem 2020 e 2021} \\
\cmidrule(lr){2-4}\cmidrule(lr){5-7}
Modelo & Efeito & Coef. & EP & Efeito & Coef. & EP \\
\midrule
Mortalidade (geral)        & 5,8\% $\to$ 4,6\% & $-0{,}2237$ & 0,1311 & 5,6\% $\to$ 4,4\% & $-0{,}2329$ & 0,1157 \\
Fração alta complexidade   & 8,9\% $\to$ 7,6\% & $-0{,}2044$ & 0,2348 & 8,9\% $\to$ 7,7\% & $-0{,}1937$ & 0,2333 \\
Faturamento real por saída & $-8{,}3\%$  & $-0{,}0866$ & 0,0948 & $-8{,}0\%$  & $-0{,}0830$ & 0,0978 \\
Tempo de permanência       & $-10{,}2\%$ & $-0{,}1080$ & 0,0468 & $-11{,}4\%$ & $-0{,}1207$ & 0,0527 \\
Produção (saídas)          & $+43{,}7\%$ & $0{,}3629$  & 0,0895 & $+44{,}7\%$ & $0{,}3694$  & 0,0892 \\
Ocupação de internação     & $+22{,}7\%$ & $0{,}2044$  & 0,0779 & $+22{,}3\%$ & $0{,}2015$  & 0,0753 \\
Ocupação de UTI            & $+19{,}3\%$ & $0{,}1769$  & 0,0603 & $+22{,}1\%$ & $0{,}1995$  & 0,0600 \\
\quad {\footnotesize cálculo alternativo} & $+22{,}8\%$ & $0{,}2058$ & 0,0747 & $+24{,}3\%$ & $0{,}2175$ & 0,0723 \\
\bottomrule
\end{tabular}
\end{table}

\comofeito{cada linha é um dos sete modelos principais, e cada bloco de
três colunas é uma amostra: à esquerda o painel inteiro, à direita a mesma
estimação depois de retirar 2020 e 2021. \textbf{Efeito} é o resultado
traduzido para a unidade que o gestor usa. Nos indicadores de quantidade
ele é a variação percentual; nos dois que já são taxas, a mortalidade e a
fração de alta complexidade, ele é o par de taxas que o deslocamento liga,
e a seta se lê ``sai de'' e ``chega a'', com a taxa de partida sendo a
média observada da administração Direta naquela amostra. \textbf{Coef.} é o
mesmo resultado no valor bruto do modelo, que não se lê diretamente, e é
por isso que a coluna Efeito existe. \textbf{EP} é a margem de erro,
calculada agrupando as observações por hospital, porque onze anos do mesmo
hospital não são onze informações independentes. Nenhum efeito muda de
sinal e nenhum muda de magnitude de forma relevante, que é o que a figura
seguinte mostra em forma gráfica. A ocupação de UTI aparece em
\textbf{duas linhas} porque foi calculada de duas maneiras: a primeira é a
adotada no documento, escolhida depois que os testes de aderência do
Apêndice~\ref{ap:diagnosticos} mostraram que ela descreve melhor estes
dados; a segunda usa a mesma conta das outras seis linhas e fica na tabela
para que a coluna possa ser lida de cima a baixo na mesma régua.}
\newpage
\comoler{cada linha traz duas estimativas do mesmo efeito: o círculo
azul escuro usa todos os anos e o quadrado verde exclui 2020 e 2021.
Quanto mais próximos os dois símbolos, mais robusto o resultado; aqui
eles praticamente se sobrepõem em todas as linhas.}

\begin{figure}[H]
\centering
\includegraphics[width=.82\textwidth]{fig_est_05_robustez_sem_pandemia.png}
\caption{Contraste OSS vs Direta estimado no painel completo e na
amostra sem 2020 e 2021, para os modelos em escala percentual.}
\label{fig:robustezpandemia}
\end{figure}

\comofeito{gráfico de pontos. Cada linha é um indicador (eixo vertical) e
o eixo horizontal traz o \textbf{efeito OSS vs Direta} em variação
percentual. Cada linha traz \textbf{dois símbolos do mesmo efeito}: um
estimado no painel completo e outro na amostra que exclui 2020 e 2021.
Quanto mais próximos os dois símbolos, mais o resultado independe dos
anos de pandemia.}

\textbf{Aderência dos modelos.} A segunda bateria confere se cada modelo
de fato descreve os dados que ele resume. Todos passam, e o material
completo, com o diagnóstico de cada um e os gráficos correspondentes,
está no Apêndice~\ref{ap:diagnosticos}.

% A terceira bateria, "Descasamento de calendário nos denominadores", saiu
% daqui em 04/08/2026, a pedido expresso de João, junto com a sua tabela
% (a antiga Tabela 10, `tab:calendario`) e as duas ressalvas que a fechavam.
%
% O QUE ELA MEDIA, para quem precisar reabrir o assunto: as duas ocupações
% dividem produção acumulada no ano por capacidade cadastrada vezes 365, e
% quando numerador e denominador não cobrem o mesmo período a razão mede
% calendário e não desempenho. A bateria não corrigia o problema — ele não é
% corrigível na granularidade anual —, ela media o TETO dele, reestimando os
% contrastes sem os conjuntos que continham as observações afetadas. O
% resultado era tranquilizador em toda a escada: excluir os 30 anos parciais
% de UTI move o contraste de 19,3% para 19,8% (0,06 erro-padrão), e no degrau
% mais severo os dois indicadores de ocupação se deslocam em DIREÇÕES
% OPOSTAS (+0,31 e −0,38 erro-padrão), assinatura de ruído de amostragem e
% não de viés sistemático, que empurraria ambos para o mesmo lado.
%
% Foi por esse resultado que a bateria saiu: ela custava três páginas e uma
% tabela de dez linhas para concluir que nada do que o documento afirma
% depende do assunto. Toda ela continua reproduzível por
% `analises/verificar_bateria_calendario.py` e `auditar_uti_ano_parcial.py`,
% e o registro da decisão está na seção 21 de PENDENCIAS.md.
Fecham o quadro as ressalvas permanentes de leitura. A identificação
do efeito de gestão dentro de hospital repousa em \textbf{cinco
conversores}, três deles com um único ano após a mudança; as margens
de erro agrupadas com tão poucas unidades tratadas tendem a ser
otimistas, e os efeitos devem ser lidos como ordens de grandeza. A
categoria Privado (3 CNES) permaneceu como controle de absorção em
todos os modelos, sem interpretação. E os resultados de mortalidade e
ocupação do estudo de eventos carregam a contaminação de tendências
anteriores já descrita, que o tempo de permanência não apresenta.

\clearpage
\section{O teste de estresse final: quanta confiança os efeitos merecem}
\label{sec:inferenciafinal}

A Seção~\ref{sec:robustez} encerrou com um aviso: todo o efeito de
gestão medido dentro de hospital repousa em cinco conversores, e
margens de erro calculadas com tão poucas unidades tratadas tendem a
ser otimistas. Esta seção fecha a pesquisa quantificando exatamente
isso, com três instrumentos usados na literatura para situações de
poucos tratados, e resolve a pendência técnica da UTI.

\textbf{Conversores de mentira.} O teste mais intuitivo é o de
permutação: sorteamos 5 hospitais falsos entre os 284 que nunca
mudaram de gestão, fingimos que eles converteram nos mesmos anos dos
conversores reais e medimos o efeito que o modelo atribuiria a essa
conversão inventada. Repetindo o sorteio 1.999 vezes, obtemos a régua
do acaso. O resultado: a queda de faturamento real atribuída à OSS é comum
entre os falsos conversores (40\% deles produzem efeito tão grande
quanto o real), enquanto \textbf{a queda de tempo de permanência de
10,2\% é maior que a de 85\% dos placebos} (p igual a 0,153). Um
segundo método, o bootstrap selvagem por cluster, sabidamente
conservador nesta situação, aponta na mesma direção e com p ainda mais
alto: 0,324 no tempo de permanência e 0,597 no faturamento.

Somando ao controle sintético descrito adiante, são \textbf{três testes
de poder limitado} e vale registrar o que eles têm em comum. Permutação
(p igual a 0,153), bootstrap selvagem (p igual a 0,324) e controle
sintético (com ajuste pré degenerado, ver adiante) convergem em
\textbf{não conseguir distinguir um efeito nulo de uma ausência de
poder para detectá-lo}. Com cinco hospitais tratados, e apenas dois com
janela posterior utilizável, é isso que esses testes conseguem dizer:
não que o efeito não exista, e sim que este desenho não tem como
separar as duas hipóteses. A tradução: o
efeito no tempo de permanência é consistente em direção e tamanho, mas
\textbf{cinco hospitais não bastam para sustentar a afirmação com o
rigor estatístico convencional}; ele deve ser tratado como evidência
sugestiva forte, não como fato demonstrado.

\textbf{Um estimador moderno para conversões em anos diferentes.} Como
as conversões ocorreram em 2019, 2023 e 2025, a comparação clássica
pode misturar hospitais já convertidos no grupo de controle. O
estimador de Callaway e Sant'Anna evita isso comparando cada coorte de
conversão apenas com os hospitais que nunca mudaram. A
Figura~\ref{fig:cstwfe} mostra que, no ano da conversão, os dois
métodos coincidem em todos os indicadores. A figura traz apenas a
mortalidade e o tempo de permanência, mas a coincidência vale para os
cinco, e a Tabela~\ref{tab:cstwfezero} traz os cinco pares para que o
leitor não precise tomar a afirmação como palavra: as diferenças entre
os dois métodos aparecem só na terceira casa decimal. Nas janelas intermediárias
o método clássico exagera o efeito — na mortalidade (2,6 e 2,9 pontos
contra 2,0 e 2,3 do alternativo) e no tempo de permanência (0,25
contra 0,16 em log pontos) — e as janelas longas se aproximam, com
intervalos sobrepostos (na mortalidade, 2,23 pontos negativos no
método clássico contra 2,44 no alternativo; no tempo de permanência,
0,28 contra 0,30 em log pontos). Nenhuma conclusão muda; as
magnitudes das janelas intermediárias — da mortalidade e do tempo de
permanência — eram os pontos inflados.\footnote{A comparação foi feita
para os cinco indicadores, em quatro janelas cada um; o texto e a figura
discutem a mortalidade e o tempo de permanência, que são os dois
achados em disputa. Nos outros três os dois métodos também concordam:
no faturamento real ambos ficam negativos nas quatro janelas, na
ocupação de internação ambos partem de perto de zero e voltam a zero na
janela longa, e na fração de alta complexidade ambos oscilam sem
padrão. A tabela completa está em \texttt{tab\_est\_cs\_evento.csv}.}

\comoler{os círculos escuros repetem o estudo de eventos da
Seção~\ref{sec:efeitogestao}; os quadrados verdes são o estimador
alternativo. Quando os dois símbolos quase se tocam, como no ano zero
e na janela longa, o resultado não depende do método; onde eles se
afastam --- nos anos intermediários, tanto da mortalidade quanto do
tempo de permanência --- é onde o método clássico exagerava.}

\begin{table}[H]
\centering
\caption{Os dois métodos no \textbf{ano da conversão}, nos cinco
indicadores estimados. TWFE é o método clássico do estudo de eventos da
Seção~\ref{sec:efeitogestao}; CS é o estimador de Callaway e
Sant'Anna. A figura desta seção mostra as duas primeiras linhas; as
outras três estão aqui para que a coincidência entre os métodos possa
ser conferida em todos os indicadores, e não apenas nos dois
ilustrados.}
\label{tab:cstwfezero}
\small
\begin{tabular}{llrr}
\toprule
Indicador & Unidade & TWFE & CS \\
\midrule
Mortalidade geral          & pontos percentuais & $-1{,}296$ & $-1{,}293$ \\
Tempo de permanência       & log pontos         & $-0{,}145$ & $-0{,}145$ \\
Faturamento real por saída & log pontos         & $-0{,}054$ & $-0{,}054$ \\
Fração alta complexidade   & pontos percentuais & $-0{,}395$ & $-0{,}382$ \\
Ocupação de internação     & log pontos         & $0{,}110$  & $0{,}110$ \\
\bottomrule
\end{tabular}
\end{table}

\comofeito{cada linha é um indicador e as duas colunas de número trazem
o efeito estimado \textbf{no próprio ano da conversão} (evento zero) por
cada um dos dois métodos, na unidade indicada. Quanto mais próximos os
dois valores, menos o resultado daquele indicador depende da escolha do
estimador.}

\begin{figure}[H]
\centering
\subfig{fig_est_07_cs_twfe_mort.png}{Mortalidade geral, em pontos
percentuais}
\subfig{fig_est_07_cs_twfe_tmp.png}{Tempo de permanência, em log pontos}
\caption{Estudo de eventos estimado pelo método clássico (TWFE, com
intervalo de confiança) e por Callaway e Sant'Anna, que compara cada
coorte de conversão só com hospitais nunca convertidos.}
\label{fig:cstwfe}
\end{figure}

\comofeito{estudo de eventos com \textbf{dois estimadores sobrepostos}. O
eixo horizontal traz os anos desde a conversão e o vertical, o efeito
estimado (mortalidade em pontos percentuais; TMP em log pontos). Os
círculos escuros, com intervalo de confiança, são o método clássico
(TWFE) e os quadrados verdes são o estimador de Callaway e Sant'Anna, que
compara cada coorte de conversão apenas com hospitais que nunca mudaram
de gestão. A proximidade dos dois símbolos mede o quanto o resultado
depende do método.}

\textbf{O gêmeo sintético.} Para Sorocaba e Pérola Byington, os dois
conversores com janela posterior utilizável, construímos um
\emph{controle sintético}: uma combinação ponderada de hospitais nunca
convertidos que reproduz a trajetória do hospital antes da mudança de
gestão. Se a conversão teve efeito, o hospital real deve descolar do
seu gêmeo depois dela. A Figura~\ref{fig:sintetico} traz o resultado
para a mortalidade, e ele pede duas honestidades. A primeira:
\textbf{em Sorocaba não há descolamento} — o gêmeo também cai depois
de 2019 (confirmando que parte da queda era retorno natural ao
patamar) e o hospital real fica, em média, \textbf{0,54 ponto
percentual abaixo do gêmeo}, lacuna que os placebos produzem com
facilidade (p igual a 0,67). A segunda é de robustez, pelo mesmo
princípio da variante sem conversores da Seção~\ref{sec:fase2}: no
painel anterior à ETAPA F (314 hospitais), este mesmo exercício dava
Sorocaba \textbf{1,09 ponto \emph{acima} do gêmeo} — a exclusão dos 25
hospitais estruturalmente incomparáveis trocou os doadores e o
\textbf{sinal da lacuna inverteu}. Não é um caso de ajuste pré-período
pior: o RMSPE pré --- a medida de quanto o gêmeo sintético erra ao
reproduzir o hospital \emph{antes} da conversão, em que zero significa
reprodução exata --- permanece exatamente 0,0 nas duas rodadas,
interpolação
perfeita em ambas --- e é justamente por isso que a inferência aqui usa
o placebo baseado na lacuna pós-conversão, e não a razão entre RMSPE pós
e pré (a estatística usual do controle sintético), que fica indefinida
quando o ajuste pré é perfeito. A causa é a \textbf{degeneração de
pesos}, e vale explicá-la em português comum porque ela reaparece
adiante e é o que sustenta boa parte das ressalvas desta seção: dizemos
que um resultado é \emph{degenerado} quando o método encontra não uma
resposta, mas muitas respostas igualmente boas, e escolhe uma delas por
acaso do cálculo. O número que sai não é errado; ele simplesmente não é
único, e por isso não carrega informação sobre o hospital. Aqui, com
apenas 4 anos
pré-conversão (2015 a 2018) e centenas de doadores potenciais, várias
combinações de pesos produzem ajuste pré-período idêntico e perfeito;
trocar quais doadores estão disponíveis levou o otimizador a uma
combinação igualmente válida no pré-período, mas com leitura oposta no
pós. Um resultado cujo sinal depende da composição do pool de doadores
não sustenta leitura direcional em nenhum dos dois sentidos. A isso
soma-se um problema que não é estatístico e sim de desenho, descrito na
Subseção~\ref{sub:oqueforam}: a conversão de Sorocaba foi decidida
\emph{em resposta} à crise de gestão de 2018, de modo que o período
pré-conversão retrata o hospital no seu pior momento. Um gêmeo sintético
construído para reproduzir esse período reproduz também a crise, e a
melhora posterior que se observa nos dois --- no hospital real e no
gêmeo --- é compatível com reversão à média sem efeito algum da nova
gestão. Já
\textbf{no Pérola Byington o hospital
fica 1,06 ponto abaixo do gêmeo}, na direção esperada, mas dentro do
que sorteios de placebo produzem (p igual a 0,32). No faturamento real, a
lacuna do Pérola é de R\$ 354 a menos por saída (p igual a 0,25).
\textit{Vale aqui a ressalva do regime de pagamento enunciada na
Seção~\ref{sec:resultados}: para as OSS, o faturamento por saída é
produção sem consequência financeira direta.} Em nenhum caso o descolamento é raro o bastante para
valer como prova.

\textbf{O exercício completo, incluindo o tempo de permanência.} O
controle sintético foi rodado para os três indicadores nos dois
conversores, e a Tabela~\ref{tab:sintetico} traz as seis linhas. As
três primeiras são as discutidas acima; as três restantes não haviam
sido apresentadas, e é preciso dizer com cuidado o que elas mostram --- e
o que não mostram. No tempo de permanência, o hospital real fica
0,24 dia abaixo do gêmeo em Sorocaba e 0,13 dia abaixo no Pérola, sobre
permanências medianas de cerca de 7,1 e 3,4 dias, com p de 0,78 e 0,75.
\textbf{Isso não refuta o efeito no tempo de permanência: apenas não o
corrobora.} A diferença importa. Com dois hospitais utilizáveis e um
ajuste pré degenerado, um p alto significa que o teste não distingue
"o efeito é nulo" de "não há poder para detectar um efeito deste
tamanho" --- exatamente a mesma limitação de poder que a permutação e o
bootstrap selvagem enfrentam no começo desta seção. A coluna do RMSPE
pré é a que sustenta a ressalva: as cinco linhas com valor 0,0 têm
interpolação perfeita no período anterior à conversão, a degeneração de
pesos descrita acima, e por isso são pouco informativas em qualquer
direção. Registramos as seis porque o critério tem de valer por igual:
a linha de mortalidade de Sorocaba, que sustenta parte do argumento
desta seção, tem a mesma degeneração das linhas de permanência.

\begin{table}[H]
\centering
\caption{Controle sintético: as seis combinações de hospital e
indicador efetivamente estimadas. A \textbf{lacuna pós} é a distância
média entre o hospital real e o seu gêmeo depois da conversão (negativa
significa hospital abaixo do gêmeo). O \textbf{RMSPE pré} mede o erro de
ajuste no período anterior à conversão: o valor 0,0 indica interpolação
perfeita e, com ela, a degeneração de pesos discutida no texto. O
\textbf{p placebo} é a proporção de hospitais nunca convertidos que
produzem lacuna igual ou maior quando submetidos ao mesmo exercício.}
\label{tab:sintetico}
\small
\begin{tabular}{llrrrr}
\toprule
Hospital & Indicador & Lacuna pós & RMSPE pré & $p$ placebo & Doadores \\
\midrule
Sorocaba & Mortalidade (p.p.)        & $-0{,}54$   & 0,00   & 0,67 & 283 \\
Sorocaba & Permanência (dias)        & $-0{,}24$   & 0,00   & 0,78 & 283 \\
Sorocaba & Faturamento real (R\$)    & $-56{,}38$  & 0,00   & 0,75 & 283 \\
Pérola   & Mortalidade (p.p.)        & $-1{,}06$   & 0,29   & 0,32 & 283 \\
Pérola   & Permanência (dias)        & $-0{,}13$   & 0,00   & 0,75 & 283 \\
Pérola   & Faturamento real (R\$)    & $-353{,}92$ & 0,00   & 0,25 & 283 \\
\bottomrule
\end{tabular}
\end{table}

\comofeito{cada linha é uma combinação de hospital convertido e
indicador. \textbf{Lacuna pós} é a média, ao longo dos anos posteriores
à conversão, da diferença entre o hospital real e o seu gêmeo sintético,
na unidade do indicador. \textbf{RMSPE pré} é a raiz do erro quadrático
médio do ajuste no período \emph{anterior} à conversão --- quanto menor,
melhor o gêmeo reproduz o hospital antes da mudança, e o valor exato de
0,0 indica que a reprodução foi perfeita, situação em que os pesos do
gêmeo deixam de ser únicos. \textbf{$p$ placebo} é a fração dos
hospitais nunca convertidos que, submetidos ao mesmo exercício, produzem
lacuna de tamanho igual ou maior. \textbf{Doadores} é o número de
hospitais elegíveis para compor o gêmeo.}

\comoler{compare a linha preta com a azul tracejada depois da linha
vertical: é essa distância que a Tabela~\ref{tab:sintetico} resume na
coluna da lacuna. Nos dois hospitais as duas linhas seguem praticamente
juntas no período posterior, que é o que os $p$ de 0,78 e 0,75
traduzem em número.}

\begin{figure}[H]
\centering
\subfig{fig_est_06_sintetico_2081695_tmp.png}{Conjunto Hospitalar
de Sorocaba (CNES 2081695), conversão em 2019}
\subfig{fig_est_06_sintetico_2078287_tmp.png}{Pérola Byington
(CNES 2078287), conversão em 2023}
\caption{Tempo de permanência: hospital real contra o seu controle
sintético. O exercício não corrobora o efeito estimado pelos demais
métodos --- e, pelas razões de poder discutidas no texto, também não o
refuta.}
\label{fig:sinteticotmp}
\end{figure}

\comofeito{mesma construção do controle sintético da figura anterior
(hospital real em linha preta, gêmeo sintético em azul tracejado, linha
vertical no ano da conversão), agora com o \textbf{tempo de permanência
em dias} no eixo vertical em lugar da mortalidade.}

\comoler{a linha preta é o hospital; a azul tracejada é o gêmeo
sintético, montado para imitar o hospital antes da linha vertical. O
efeito da conversão é a distância entre as duas linhas depois da linha
vertical. Em Sorocaba as duas linhas caem juntas, com a preta um pouco
abaixo, e terminam praticamente sobrepostas; no Pérola, a preta cai
abaixo da azul, mas com o gêmeo também caindo.}

\begin{figure}[H]
\centering
\subfig{fig_est_06_sintetico_2081695_mort_pp.png}{Conjunto Hospitalar
de Sorocaba (CNES 2081695), conversão em 2019}
\subfig{fig_est_06_sintetico_2078287_mort_pp.png}{Pérola Byington
(CNES 2078287), conversão em 2023}
\caption{Mortalidade geral: hospital real contra o seu controle
sintético, construído só com hospitais nunca convertidos.}
\label{fig:sintetico}
\end{figure}

\comofeito{controle sintético. O eixo horizontal é o ano e o vertical, a
mortalidade geral (em pontos percentuais). A linha preta é o hospital
real; a azul tracejada é o \textbf{gêmeo sintético}, uma combinação
ponderada de hospitais que nunca converteram, montada para reproduzir a
trajetória do hospital \emph{antes} da linha vertical (a conversão).
Depois da linha vertical, a distância entre as duas linhas é a estimativa
do efeito da conversão.}

\textbf{A pendência da UTI, resolvida.} A intensidade de uso de UTI
foi reestimada com a família Gama, adequada à cauda de UTIs de uso
quase nulo que o modelo em logaritmo não acomodava. O diagnóstico de
resíduos salta de 0,814 para \textbf{0,944}
(Figura~\ref{fig:qqgama}, no Apêndice~\ref{ap:diagnosticos}),
e o resultado substantivo permanece: uso de UTI \textbf{19,3\% maior
nas OSS} (p igual a 0,003), com a variante de corte mínimo de
atividade confirmando (23,3\%). Este valor \textbf{substitui} os
22,8\% que a especificação em logaritmo dava na
Seção~\ref{sec:resultados}: a direção e a significância se mantêm, o
nível cai cerca de três pontos e a estimativa citável deste indicador
passa a ser a da reestimação.

\textbf{Balanço final da fase de estimação.} O quadro que resiste a
toda a bateria de testes é: as OSS operam um regime assistencial
claramente mais intenso (mais produção, mais ocupação, mais uso de
UTI), com diferenças firmes \textbf{enquanto retrato transversal das
redes}; a única evidência com pretensão causal que resiste é a do
tempo de permanência, e é preciso ser exato sobre o que isso quer
dizer: ela \textbf{sobrevive a todos os testes que têm poder para
avaliá-la} --- o teste de tendências anteriores (p igual a 0,254), o
estimador de Callaway e Sant'Anna e a reestimação sem o biênio
pandêmico --- e \textbf{fica sem corroboração em todos os que têm o
poder limitado pelo número de tratados}: permutação (p igual a 0,153),
bootstrap selvagem (p igual a 0,324) e controle sintético (degenerado,
p de 0,78 e 0,75). Nenhum destes três a contradiz; nenhum a sustenta.
Ela permanece sugestiva com cinco conversores, não definitiva; o faturamento
real por saída não mostra diferença demonstrável em nenhum método; e
a queda de mortalidade após as conversões, embora visível nos dados
brutos, não produz descolamento raro frente ao gêmeo sintético em
nenhum dos dois conversores testados — e em Sorocaba o próprio sinal
da lacuna inverte com a mudança de painel da ETAPA F, o que rebaixa o
resultado a frágil, não adverso nem favorável. A segunda fase da estimação
(Seção~\ref{sec:fase2}) quantificou esse último ponto da forma mais
direta possível: retirando os cinco conversores da amostra, a
diferença de mortalidade entre as redes OSS e Direta simplesmente
desaparece (de 0,262 negativo o coeficiente vira 0,003 positivo, com
intervalo de 0,414 negativo a 0,420 positivo, estatisticamente nulo),
ou seja, \textbf{a mortalidade menor é um fenômeno da trajetória dos
hospitais que mudaram de gestão, sem contraparte na comparação entre
as redes}. Novos anos de dados, especialmente
para os três hospitais convertidos em 2025, são o caminho natural para
transformar as sugestões em conclusões.

\clearpage
\section{A segunda camada de evidência: hierarquia e eficiência}
\label{sec:fase2}

A segunda fase da estimação atacou a mesma pergunta com três famílias
de métodos ainda não usadas no documento: modelos
\textbf{hierárquicos} (que tratam os 11 anos de cada hospital como
parentes, não como observações soltas), a análise envoltória de dados
(\textbf{DEA}, que constrói uma fronteira de melhores práticas e mede
a distância de cada hospital até ela) e a fronteira estocástica
(\textbf{SFA}, a contraparte paramétrica do DEA). O pacote completo de
resultados, com scripts, diagnósticos e um relatório de execução
auditável, está na pasta resultados\_fase2 do repositório; aqui ficam
os achados que mudam ou reforçam a leitura do documento.

\textbf{O que a camada hierárquica confirma.} Estimados com um
efeito próprio para cada hospital, o tempo de permanência e o faturamento
real reproduzem quase exatamente os resultados da
Seção~\ref{sec:resultados} (queda de 10,2\% no TMP; faturamento sem
diferença), a produção continua maior nas OSS (33\%, com intervalo
inteiro acima do zero) e a ocupação aponta na mesma direção com
magnitude menor (11,5\%), \textbf{mas com intervalo que cruza o zero na
especificação hierárquica}, em ambos os casos porque parte do
contraste entre as redes é absorvida pelo efeito de hospital; e a mortalidade aparece
\textbf{0,262 menor no logito} (intervalo de 0,410 a 0,114 negativos;
na escala da própria taxa, isso equivale a cerca de \textbf{1,3
ponto percentual a menos de mortalidade}, na altura da mediana do
painel, que é a forma interpretável do mesmo número),
com precisão maior que a da
Seção~\ref{sec:resultados}.\footnote{Duas especificações alternativas
foram estimadas para esse mesmo coeficiente e não são reportadas aqui
por apontarem na mesma direção e magnitude: sem o termo de Mundlak,
0,233 negativo; com efeito aleatório de ano, 0,260 negativo. A tabela
completa está em \texttt{tabB\_sensibilidade\_mortalidade.csv}.} O teste
decisivo, porém, é outro: \textbf{retirando os cinco conversores, a
diferença de mortalidade entre as redes desaparece}: o coeficiente vai
a 0,003 positivo, com intervalo de 0,414 negativo a 0,420 positivo,
cruzando o zero com folga, enquanto TMP, produção, ocupação e
eficiência não só sobrevivem como se \textbf{fortalecem}: sem os
conversores o TMP vai de 0,108 para 0,281 negativo, a produção de 0,282
para 0,492, a ocupação de 0,108 para 0,218 (passando a excluir o zero,
o que não fazia no modelo principal) e a eficiência de 0,293 para
0,331 negativo. Os dois coeficientes, lado a lado, são
a evidência direta de que a identificação causal repousa sobre os
conversores (sobretudo Sorocaba e Pérola Byington, os dois com janela
posterior utilizável), e não sobre uma diferença ampla entre as redes:
a mortalidade menor é uma história dos hospitais que mudaram de
gestão, não uma diferença entre as redes OSS e Direta. E essa história
tem explicação nomeada, não apenas a constatação de que o efeito
desaparece sem os conversores: como a Subseção~\ref{sub:oqueforam}
documenta, a conversão de Sorocaba foi decidida em resposta a uma crise
de gestão, o que deprime o período anterior e faz da reversão à média a
explicação mais econômica para a melhora posterior; e no Pérola Byington
a mudança de gestão coincidiu com a mudança para um prédio novo, o que
torna o tratamento confundido. As duas alternativas não são hipóteses de
gabinete: são o que de fato aconteceu nos dois hospitais que sustentam o
resultado.

\comoler{cada linha compara duas estimativas do efeito da mesma
categoria: o círculo escuro trata cada categoria isoladamente; o
quadrado verde deixa as categorias emprestarem informação umas às
outras, na proporção da quantidade de dados de cada uma. Quanto menor
o grupo, maior o recuo do quadrado em direção ao zero: o Privado, com
3 hospitais, quase não sustenta efeito próprio. É a ilustração formal
da regra deste documento de nunca interpretar essa categoria.}

\begin{figure}[H]
\centering
\includegraphics[width=.58\textwidth]{figB02_encolhimento_categorias.png}
\caption{Efeitos de categoria na mortalidade com e sem partial
pooling: com apenas 3 hospitais, o efeito aparente do Privado encolhe
em quase 90\% quando o modelo pondera a evidência pelo tamanho do
grupo.}
\label{fig:encolhimento}
\end{figure}

\comofeito{gráfico de pontos. Cada linha é uma categoria de gestão (eixo
vertical) e o eixo horizontal traz o \textbf{efeito da categoria no
logito da mortalidade} (centrado). Cada linha traz duas estimativas: o
círculo escuro trata cada categoria isoladamente e o quadrado verde
aplica \emph{partial pooling}, em português "agrupamento parcial": o
modelo trata as categorias como aparentadas e permite que uma tome
emprestada a informação das outras, tanto mais quanto menos dados ela
própria tiver. Quanto menor o grupo, mais o quadrado recua em direção
ao zero.}

\textbf{Eficiência: a ressalva vem antes do resultado.} O DEA e a SFA
medem eficiência produtiva: quanto produto (saídas, incluída a alta
complexidade) cada hospital extrai dos seus insumos. Um dos insumos é
o \textbf{valor total real} pago pelo SUS, e ele exige a ressalva de
primeira ordem desta seção: o valor do SIH é reembolso por produção
faturada, não medida direta de recursos consumidos, e os regimes de
pagamento diferem por categoria (as OSS recebem pelo contrato de
gestão, de modo que o faturamento por AIH pode divergir do recurso
efetivo; os filantrópicos convivem com subfinanciamento e subsídio
cruzado). Parte da distância entre categorias pode refletir regime de
pagamento, não gestão. Duas verificações sustentam que a ordenação
não é artefato do insumo financeiro: a \textbf{produção física por
leito}, que não usa dinheiro nenhum, já ordena as redes do mesmo jeito
(55 saídas por leito nas OSS e 51 nos Municipais, contra 36 na Direta
e 35 nos Filantrópicos); e o DEA reestimado \textbf{só com leitos como
insumo} preserva e até amplia a vantagem (medianas corrigidas de viés
de 0,65 nas OSS contra 0,42 na Direta), com a segunda etapa
confirmando a distância
menor das OSS à fronteira (coeficiente de 1,03 negativo no logaritmo
da distância, intervalo bootstrap de 1,67 a 0,50 negativos, inteiro
longe do zero). A ressalva fica assim resolvida \textbf{no nível das
categorias} para a distância entre OSS e as demais: o afastamento da
OSS na frente não depende do insumo financeiro em nenhuma das
verificações. Já a ordem interna das outras três categorias é frágil,
como os números detalhados adiante mostram, e \textbf{os dois métodos
discordam sobre ela}. No DEA, Direta e Público Municipal ficam separados
por 0,002 (medianas corrigidas de viés de 0,634 e 0,632), um empate; na
SFA, quem empata é outro par, Direta e Filantrópico, por 0,004 (0,788
contra 0,784), enquanto o \textbf{Público Municipal aparece em 0,912},
isto é, 0,125 acima da Direta. Os dois métodos põem a OSS na frente e o
Privado no fim, e \textbf{trocam entre si a segunda e a terceira
posições}. Abaixo da OSS, portanto, o que os dados sustentam é um
\textbf{empate entre Direta e Público Municipal no DEA} (que é o
instrumento em que as conclusões de eficiência deste documento se apoiam,
pelo motivo estabelecido adiante), com o
Filantrópico atrás, e não uma hierarquia firme de quatro posições. Ela
\textbf{não} se estende ao hospital individual: a correlação de ordem
entre os escores da versão financeira e da versão só com leitos, ambas
corrigidas de viés, é de 0,51 por hospital-ano (0,48 entre as medianas
por hospital), o ranking hospital a hospital embaralha quando o
dinheiro sai do modelo, e a classificação de eficiência de um hospital
específico só deve ser citada depois de conferida nas duas versões.

\newpage
Com a ressalva posta, o resultado: na eficiência corrigida de viés
(bootstrap de Simar e Wilson, 2.000 réplicas por ano), as medianas do
período são \textbf{OSS 0,733, Direta 0,634, Público Municipal 0,632 e
Filantrópico 0,595}. Direta e Público Municipal praticamente empatam
(diferença de 0,002), sugerindo três patamares em vez de quatro: OSS à
frente, Direta e Municipal próximas no meio, Filantrópico atrás. Na
regressão da distância à fronteira sobre
categoria, complexidade, porte e ano, a vantagem OSS equivale a
\textbf{cerca de 25\% menos distância à fronteira que a Direta}
(exp(-0,293) menos 1), estável retirando o biênio pandêmico (0,272
negativo), os cinco conversores (0,331 negativo) ou restringindo ao
suporte comum (0,319 negativo), as quatro estimativas com intervalo
inteiro abaixo do zero. Ao contrário da mortalidade,
\textbf{a eficiência
maior das OSS é um traço da rede}, presente na comparação pura entre
hospitais distintos. Vale relembrar aqui, porque a frase anterior é o
tipo de afirmação que viaja sozinha, que essa conclusão é firme
\textbf{no nível das categorias} e não desce ao hospital individual: a
correlação de ordem entre a versão financeira e a versão só com leitos
é de apenas 0,51, como o parágrafo anterior registrou. E o receio de os 3 hospitais Privados
distorcerem a fronteira não se confirma: retirados, os escores dos
demais praticamente não se movem (correlação de ordem 1,000), porque
o perfil deles opera longe da fronteira.

\comoler{como nos violinos das seções anteriores: a silhueta engorda
onde há mais hospitais e o número é a mediana do grupo. Os dois painéis
concordam nas pontas, com a OSS na frente e o Privado no fim, e
\textbf{discordam no meio}: no painel (a), do DEA, a Direta (0,634) e o
Município (0,632) empatam; no painel (b), da fronteira estocástica, o
Município (0,912) sobe bem acima da Direta (0,788). O painel (b) é
diagnóstico, não medida: pelo motivo estabelecido adiante, nenhum número
dele é citável. O Privado aparece só como registro.}

\begin{figure}[H]
\centering
\subfig{figE01_violino_te_categoria.png}{Eficiência DEA corrigida de
viés}
\subfig{figS01_te_categoria.png}{Eficiência técnica da fronteira
estocástica}
\caption{Eficiência por categoria nos dois métodos. O painel (a), do DEA,
é a medida em que este documento se apoia. O painel (b), da fronteira
estocástica, é \textbf{diagnóstico e não medida}: a concentração das OSS
junto ao valor 1 é uma degeneração numérica do modelo, e a fronteira que
o gerou tem a monotonicidade violada em 94,9\% das observações, de modo
que \textbf{nenhum valor do painel (b) é citável, em nenhuma categoria}
(o texto explica). Como nos violinos da
Seção~\ref{sec:categorias}, onde a silhueta é mais gorda há mais
hospitais com aquele valor de eficiência; a largura não é percentual nem
contagem, e não se compara entre categorias.}
\label{fig:eficiencia}
\end{figure}

\comofeito{gráficos de violino, um por categoria de gestão no eixo
horizontal. O eixo vertical traz a \textbf{eficiência técnica} (entre 0 e
1, quanto mais perto de 1 mais próximo da fronteira de melhores
práticas): no painel (a), a eficiência DEA corrigida de viés, e no (b), a
eficiência da fronteira estocástica (SFA). Cada violino é a distribuição
da categoria e o número é a mediana. No painel (b) a concentração das OSS
junto a 1 é uma degeneração numérica do modelo, e nenhum valor do painel
é citável, em nenhuma categoria, pelo motivo estabelecido no texto.}

\newpage
\textbf{A contraparte paramétrica, e o que ela não é.} A
fronteira estocástica (Translog (uma forma matemática mais flexível
que a alternativa clássica, Cobb Douglas, por não obrigar os insumos a
se combinarem sempre na mesma proporção), preferida ao Cobb Douglas
com folga larga: a estatística de razão de verossimilhança entre as duas
formas é de 567,0 com 6 graus de liberdade, com p da ordem de
$10^{-119}$; e 95,5\% da variância composta é atribuída à ineficiência)
põe as OSS na mesma ponta de cima em que o DEA as põe, e traz
três honestidades a registrar, que juntas retiram dela o estatuto de
medida:
as OSS ficam tão coladas à fronteira paramétrica que o coeficiente
delas na equação de ineficiência perde identificação (o análogo da
separação perfeita da Seção~\ref{sec:estrategia}), de modo que nem o
nível nem a posição das OSS na SFA são citáveis,
e a não-identificação aparece igual nas quatro amostras e nas
três sensibilidades, com o ponto oscilando entre $-772$ e $-1.971$ e o
erro-padrão entre 400 e 1.896, que é a assinatura do problema e não uma
faixa de estimativas;\footnote{Numa das amostras (sem o biênio
pandêmico) o valor sai com $p$ de $8{,}3\times10^{-7}$. Isso \emph{não}
é precisão: sob separação a verossimilhança é praticamente plana numa
região distante de zero, e o intervalo de Wald calculado ali não tem
interpretação. Um $p$ pequeno sobre um parâmetro degenerado mede o
formato da aproximação, não o dado.} \textbf{o insumo de UTI é mal
comportado na forma Translog}, com a monotonicidade violada em
\textbf{94,9\% das observações} (97,8\% entre as que têm leito de UTI
habilitado, tenham-no usado ou não no ano; 96,9\% se o corte for ter
registrado diária de UTI),
contra 99,7\% de respeito para leitos e 99,6\% para valor: herança
do logaritmo de um mais UTI com 19\% de zeros e da colinearidade alta
com leitos totais e valor, e o motivo pelo qual qualquer leitura fina
de UTI nesta fronteira precisa da sensibilidade que a remove; e, na
outra ponta, Direta (0,788) e Filantrópico (0,784)
praticamente empatam em ponto estimado, a mesma proximidade que o DEA
mostra entre Direta e Município (sem intervalo de confiança calculado
em nenhum dos dois métodos para confirmar ou afastar sobreposição). A
correlação de ordem entre os dois conjuntos de escores é de
\textbf{0,62} (0,63 excluindo as OSS, de modo que a semelhança não
depende do grupo problemático). Ela mede quanto os dois procedimentos se
parecem quando aplicados aos mesmos dados, e \textbf{não é corroboração
de um pelo outro}: os dois consomem a mesma informação, e um deles, pelo
motivo estabelecido adiante, não descreve uma função de produção.
A semelhança não cresce de forma gradual: fica achatada entre 0,50 e
0,55 de 2015 a 2019, salta para a faixa de 0,66 a 0,74 já em 2020 e se
mantém nesse patamar mais alto até 2025, sem tendência adicional, um
degrau que coincide com o início da pandemia, não uma trajetória
ascendente ao longo da década.

\textbf{A equação de ineficiência, e por que nenhum coeficiente dela é
citável.} A fronteira estocástica estima, em paralelo à produção, uma
segunda equação que explica \emph{quanta} ineficiência cada hospital
carrega. O coeficiente das OSS ali é o número degenerado já discutido; a
equação tem \textbf{outros sete}, e \textbf{cinco} deles são fortemente
significativos. Como a variável explicada é a ineficiência, sinal
negativo significa mais eficiente.

A tabela é reproduzida porque omiti-la seria pior: ela contém o único
resultado desta fase que o documento chegou a não mostrar, e resultado
que desaparece depois de encontrado é omissão. Mas ela vem com uma
advertência que vale para cada uma das suas linhas e que precisa ser lida
\emph{antes} dos números. \textbf{A monotonicidade da fronteira que gerou
estes coeficientes está violada em 94,9\% das observações}, como o
parágrafo anterior registrou. Uma fronteira violada em quase toda a
amostra \textbf{não se comporta como função de produção}, e significância
mede a precisão de uma estimativa, não a validade da estrutura de onde
ela sai: coeficiente preciso extraído de fronteira mal-comportada é
preciso sobre um objeto que não é o que o nome diz. Por isso, e não por
conveniência de resultado, \textbf{nenhum coeficiente desta tabela é
interpretável como magnitude de ineficiência, e nenhum entra em conclusão
alguma deste documento}.

A violação também não se concentra em categoria nenhuma, o que impede
tratá-la como problema de uma linha só: ela vai de 93,6\% no Filantrópico
a \textbf{99,3\% na administração Direta}, que é justamente a referência
de todos os contrastes da tabela. Cinco vírgula sete pontos de amplitude
sobre uma base de 93,6\% é defeito geral da forma funcional. A fronteira
inteira está desqualificada como medida, não uma parte dela.

\begin{table}[H]
\centering
\caption{Equação de ineficiência da fronteira estocástica (Battese e
Coelli 1995). Referência: administração Direta, porte médio.
Coeficiente negativo indica \emph{menos} ineficiência.
\textbf{Tabela tarjada: a monotonicidade da fronteira que produziu estes
coeficientes está violada em 94,9\% das observações, de 93,6\% no
Filantrópico a 99,3\% na própria Direta. Nenhum valor desta tabela é
interpretável como magnitude nem entra em conclusão deste documento; ela
é diagnóstico da inadequação da forma Translog.}}
\label{tab:equacaoinef}
\small
\begin{tabular}{lrrr}
\toprule
Termo & Coeficiente & Erro-padrão & $p$ \\
\midrule
Público Municipal        & $-5{,}454$ & 1,020 & $9{,}0\times10^{-8}$ \\
Filantrópico             & $-0{,}244$ & 0,163 & 0,134 \\
Privado                  & $+2{,}474$ & 0,396 & $4{,}2\times10^{-10}$ \\
Complexidade (padronizada) & $-2{,}159$ & 0,325 & $3{,}2\times10^{-11}$ \\
Grande Porte             & $+2{,}470$ & 0,333 & $1{,}3\times10^{-13}$ \\
Especial                 & $+7{,}401$ & 0,917 & $7{,}0\times10^{-16}$ \\
Tendência linear         & $-0{,}010$ & 0,023 & 0,668 \\
\bottomrule
\end{tabular}
\end{table}

\textbf{O que a tabela mostra, lida como diagnóstico e não como medida.}
O coeficiente do Público Municipal ($-5{,}454$, $p$ de
$9{,}0\times10^{-8}$) é o resultado que faltava ao documento, e é também
o que melhor ilustra a advertência acima. Ele é preciso, e aponta o
Público Municipal como menos ineficiente que a administração Direta.
\textbf{O DEA, que é o instrumento adequado, não confirma essa ordem, e
tampouco a refuta.} Nas medianas corrigidas de viés impressas acima a
Direta (0,634) fica à frente do Município (0,632); nas médias o Município
fica à frente (0,644 contra 0,620). Ano a ano, no escore corrigido de
viés, o Município supera a Direta em seis dos onze anos, e os seis são de
2020 a 2025. No estrato de alta complexidade da
Seção~\ref{sec:segmentacao}, que é onde a comparação entre redes é
viável, o Município fica à frente; no estrato intermediário, atrás.
\textbf{A ordem entre estas duas redes depende de qual escore e de qual
estatística se escolhe, e por isso este documento não a afirma em direção
nenhuma.} O que há é um sinal recorrente e sem corroboração por
instrumento adequado, o que é diferente tanto de achado quanto de nada.

As demais linhas têm o mesmo estatuto. O Filantrópico não se distingue da
Direta ($p$ de 0,134). O porte carrega o maior coeficiente da tabela
(Especial, $+7{,}401$, quase três vezes o do Grande Porte) enquanto a
complexidade anda no sentido oposto ($-2{,}159$), e não há tendência
temporal ($p$ de 0,668). Lidos como \emph{ordens de grandeza relativas
dentro de uma fronteira que se sabe mal especificada}, esses sinais
sugerem que a estrutura física do hospital move a ineficiência medida
mais do que a rede que o administra. Lidos como magnitudes, não dizem
nada, pelo motivo já enunciado.

\textbf{A regra, enunciada uma vez para todo o documento.} Nenhum
coeficiente e nenhum nível da fronteira estocástica é citável, em nenhuma
categoria de gestão, nem aqui nem nas seções seguintes. A SFA permanece
no documento como \textbf{diagnóstico da própria inadequação da forma
Translog a esta tecnologia de produção}, e a violação em 94,9\% é ela
mesma um achado: diz que a Translog não descreve bem o modo como estes
hospitais transformam insumo em produto. \textbf{As conclusões de
eficiência deste documento apoiam-se no DEA.}

\comoler{quanto mais a nuvem sobe da esquerda para a direita, mais os
dois procedimentos se parecem ao ordenar os hospitais. A faixa verde
colada no topo é a degeneração das OSS descrita no texto (o modelo
paramétrico as coloca todas na fronteira); mesmo descontando essa faixa,
a semelhança permanece. Semelhança não é confirmação: o eixo vertical vem
de uma fronteira cuja monotonicidade está violada em 94,9\% das
observações, e nenhum valor dele é citável.}

\begin{figure}[H]
\centering
\includegraphics[width=.58\textwidth]{figS02_dispersao_sfa_dea.png}
\caption{Semelhança entre os dois procedimentos de eficiência: cada ponto
é um hospital num ano, posicionado pelo escore DEA (horizontal) e pela
eficiência SFA (vertical). O eixo vertical é diagnóstico, não medida:
nenhum valor de SFA é citável.}
\label{fig:convergencia}
\end{figure}

\comofeito{gráfico de dispersão. Cada ponto é um hospital em um ano,
posicionado pelo \textbf{escore de eficiência DEA} no eixo horizontal e
pela \textbf{eficiência técnica SFA} no vertical (ambos entre 0 e 1).
Quanto mais a nuvem sobe da esquerda para a direita, mais os dois métodos
(de naturezas opostas) concordam sobre quem é eficiente. A faixa
verde colada no topo é a degeneração das OSS descrita no texto.}

\newpage
\subsection*{Como os modelos de ocupação estão especificados}
\label{sub:etapa5}

Os modelos de ocupação desta fase tomam três decisões de especificação que
convém declarar, porque nas três havia alternativa plausível e a escolha
muda o que o número significa.

\textbf{O limiar de aparagem} da cauda superior é o percentil 99 da
subamostra em que a ocupação de UTI é \emph{definida}, e não do painel
completo. Calculado sobre o painel completo, o limiar seria contaminado
pelas observações de hospitais sem UTI, que não são ocupação baixa: são
ausência de ocupação (Seção~\ref{sec:categorias}).

\textbf{A amostra} do modelo de UTI é definida pela capacidade instalada,
isto é, ter leito de UTI habilitado, independentemente de tê-lo usado no
ano. O critério alternativo (restringir às observações com ocupação
positiva) selecionaria a amostra pela própria variável que se quer
explicar.

\textbf{A ocupação de internação não é aparada acima de 100\%}, porque esse
excesso não é ruído de cauda: é a assinatura aritmética da superlotação
descrita na Seção~\ref{sec:segmentacao}. Apará-lo apagaria o sinal.

Definir a amostra pela capacidade mantém dentro dela as observações de
hospitais com UTI habilitada e nenhuma diária no ano. Como a família Gama
não admite valor zero, a especificação principal é o \textbf{modelo em duas
partes com Gama na parte contínua}, que acomoda essa massa em zero num
componente próprio e preserva a função de variância da Gama.

\textbf{O resultado.} As OSS operam a UTI \textbf{30,4\% mais
intensamente} que a administração Direta, com intervalo de 95\% entre
$+5{,}7\%$ e $+60{,}8\%$. A ocupação de internação aponta na mesma direção,
$+11{,}5\%$, mas com o intervalo cruzando o zero, de modo que ali não se
pode afirmar diferença.

\subsubsection*{A escolha de família move mais o número do que os defeitos}

Entre todas as decisões de modelagem, a que mais desloca a estimativa é a
\textbf{família de distribuição}, mais do que qualquer das três decisões
de especificação acima. Por isso as cinco especificações consideradas são
reportadas juntas.

\begin{table}[H]
\centering
\caption{Intensidade de uso da UTI, OSS contra administração Direta, sob
cinco especificações. Efeito e intervalo em percentual.}
\label{tab:familiasuti}
\small
\begin{tabular}{lrrrl}
\toprule
Especificação & $n$ & Efeito & IC 95\% & Zero \\
\midrule
\textbf{Duas partes, Gama (principal)} & 2.489 & $\mathbf{+30{,}4\%}$ & $[+5{,}7\%;\ +60{,}8\%]$ & exclui \\
Gama, apenas os positivos              & 2.482 & $+30{,}4\%$ & $[+5{,}7\%;\ +60{,}8\%]$ & exclui \\
Tweedie sem os zeros                   & 2.482 & $+27{,}0\%$ & $[+7{,}6\%;\ +49{,}8\%]$ & exclui \\
Tweedie com os zeros                   & 2.489 & $+26{,}7\%$ & $[+6{,}7\%;\ +50{,}4\%]$ & exclui \\
\textbf{Lognormal sem os zeros}        & 2.482 & $+25{,}9\%$ & $[-0{,}3\%;\ +59{,}0\%]$ & \textbf{cruza} \\
\bottomrule
\end{tabular}
\end{table}

A amplitude entre as cinco é de \textbf{4,5 pontos percentuais}, menor que a
largura do intervalo de confiança de qualquer uma delas. A direção é a mesma
nas cinco, e \textbf{quatro excluem o zero}.

\textbf{A quinta é a lognormal, e ela precisa ser nomeada}, porque
``significativo em quatro de cinco'' sem dizer qual é a exceção não é
robustez reportada. \textbf{E ela não discorda.} O ponto dela está a
\textbf{0,0087} em pontos de logaritmo do Tweedie estimado sobre exatamente
a mesma amostra, distância menor que a que separa quaisquer duas das
outras famílias. O que muda é a \textbf{precisão}: o erro-padrão da
lognormal é 41\% maior que o do Tweedie e 11\% maior que o da Gama.
\textbf{Ela perde a significância por variância, não por deslocamento}, e
por margem mínima: o limite inferior fica a três milésimos do zero.

A razão é que a lognormal \textbf{estima outra quantidade}. As demais
modelam a média condicional; a lognormal modela a mediana condicional. Sob
assimetria à direita, que é a forma desta distribuição, o efeito sobre a
mediana é sistematicamente mais atenuado e menos preciso que o efeito sobre
a média. É \emph{estimar diferente}, não \emph{estimativa discordante}, e
isso fortalece a leitura de robustez em vez de qualificá-la: a única
especificação que não atinge significância é também a única que não está
medindo a mesma coisa que as outras, e ainda assim entrega ponto
equivalente.

\comofeito{um alerta sobre a coluna de ajuste que não aparece nesta tabela:
o critério de informação da lognormal (2.135) é numericamente muito menor
que o das demais (entre 21.608 e 22.695), e \textbf{isso não significa
ajuste melhor}. O da lognormal é calculado sobre o logaritmo da resposta e o
das outras sobre a resposta, escalas diferentes. O critério é comparável
\emph{dentro} de cada escala, nunca entre elas.}

\newpage
\textbf{O que fica para a leitura do documento.} A segunda fase não
altera nenhuma conclusão; ela redistribui os pesos. O \textbf{regime
operacional mais intenso e a eficiência produtiva maior das OSS são
traços robustos da rede}, mas é preciso dizer \emph{qual} componente do
regime operacional resiste a quê. \textbf{Produção e intensidade de uso
da UTI excluem o zero nas quatro amostras}, inclusive sem os
conversores (produção de $+0{,}26$ a $+0{,}49$ em pontos de log;
intensidade de UTI de $+0{,}19$ a $+0{,}50$). \textbf{A ocupação de
internação, não:} o intervalo cruza o zero em duas das quatro, na
especificação principal ($-0{,}041$ a $+0{,}277$) e sem o biênio
pandêmico ($-0{,}077$ a $+0{,}260$), e só o exclui sem os
conversores e no suporte comum. Escrever que o regime operacional
inteiro é robusto ``em todas as amostras'' seria estender à ocupação
uma robustez que ela não tem, e que as duas passagens anteriores desta
mesma Seção já haviam registrado; o \textbf{TMP menor} reaparece
com a mesma magnitude na especificação hierárquica, o que confirma que
o resultado não dependia da forma do modelo, mas não altera a
limitação que vem dos cinco conversores, e ele segue sem corroboração
nos testes de poder limitado (permutação p igual a 0,153, bootstrap
selvagem 0,324, controle sintético degenerado); o \textbf{faturamento} segue sem
diferença em qualquer método; e a \textbf{mortalidade menor}
concentra-se inteiramente na trajetória dos cinco conversores, o
que a aproxima das ressalvas causais das
Seções~\ref{sec:efeitogestao} e~\ref{sec:inferenciafinal} e a afasta
de uma leitura de superioridade da rede OSS nesse desfecho.

\clearpage
\section{Fase qualitativa: seleção e roteiro de entrevistas}
\label{sec:qualitativa}

% comandos locais desta seção. Eram cabeçalhos em negrito soltos, o que
% deixava 17 páginas sem uma única entrada no sumário nem bookmark de
% PDF; agora são unidades de seccionamento de verdade. A aparência é
% praticamente a mesma --- \subsection já é \large\bfseries e
% \subsubsection já é \bfseries no corpo do texto ---, com o ganho da
% numeração e da navegabilidade. (xcolor já carregado no preâmbulo do
% principal.)
\newcommand{\parte}[1]{\subsection{#1}}
\newcommand{\bloco}[1]{\subsubsection{#1}}
\newcommand{\achado}[1]{%
  \par\smallskip
  \noindent\begingroup\small\itshape
  \textcolor{black!70}{Achado que a pergunta testa: #1}%
  \endgroup\par\smallskip}

\noindent\textbf{Uma nota sobre o estatuto desta seção.} O que se segue é
uma \textbf{sugestão}, construída por critérios puramente estatísticos:
os hospitais foram escolhidos pela sua posição nas distribuições de
eficiência, permanência e produção do painel, e as perguntas foram
derivadas dos achados quantitativos da fase de estimação. Nada disso
substitui o julgamento de campo. A decisão sobre \textbf{quais hospitais
de fato importam e quais perguntas fazer} cabe à experiência e ao
conhecimento institucional da pesquisadora principal, Priscilla Reinisch
Perdicaris, cuja leitura da rede hospitalar paulista é o fator decisivo
desta fase. As listas e os roteiros abaixo são apenas um ponto de partida
informado pelos dados, a ser confirmado, reordenado, enxugado ou ampliado
conforme esse julgamento.

\parte{Apresentação e método}

Esta seção operacionaliza a fase qualitativa da pesquisa sobre o efeito
do modelo institucional de gestão no desempenho dos hospitais SUS de São
Paulo, sobre o mesmo painel balanceado de \textbf{289 hospitais} de 2015
a 2025 (painel pós-ETAPA F, 15/07/2026). A fase de estimação, já
concluída, deixou quatro achados que as entrevistas devem explicar por
dentro. O primeiro é que o regime operacional mais intenso e a
eficiência técnica maior das OSS são traços robustos da rede, presentes
em todos os métodos. O segundo é que o tempo médio de permanência cerca
de 10\% menor nas OSS é evidência sugestiva de variação dentro do
hospital, apoiada em apenas cinco conversores. O terceiro é que a
mortalidade menor concentra-se na trajetória dos conversores e não
sobrevive ao controle sintético, o que recomenda cautela e investigação
de mecanismo (o exercício em Sorocaba, em particular, tem o sinal
instável entre o painel de 314 e o de 289 --- ver
Seção~\ref{sec:inferenciafinal} --- então não deve ser tomado como
achado firme em nenhuma direção). O quarto é que o faturamento real por
saída não difere entre categorias: a diferença de gestão, se existe, não
aparece no faturamento. \textit{Ressalva de leitura (regime de
pagamento): hospitais OSS são custeados por contrato de gestão com
orçamento global, não por reembolso de AIH --- o valor registrado por
saída é, para eles, produção sem consequência financeira direta; a
diferença de faturamento pode refletir o próprio regime de pagamento, a
variável cujo efeito se busca medir, e não uso de recursos ou qualidade
assistencial.} Toda referência a complexidade nesta fase usa o escore
estrutural, como no restante do documento (Seção~\ref{sec:painel}).

A seleção de entrevistados separa um \textbf{núcleo de doze entrevistas
previstas} de uma \textbf{reserva de dez entrevistas sugeridas} e
\textbf{duas entrevistas curtas de validação de registro}. Nenhum bloco
é obrigatório à exceção dos conversores; os cinco conversores da
administração Direta para OSS são os \textbf{prioritários}, por
ancorarem toda a leitura causal do documento. O núcleo, descrito a
seguir, reúne os conversores e os desviantes positivos e negativos de
eficiência persistente no suporte comum de porte e complexidade. A
reserva sugerida, cujas listas estão no Apêndice~\ref{ap:selecao},
reúne os pares casados de contraste com
estrutura parecida e eficiência divergente e os casos de maior
discordância entre DEA e SFA. As duas entrevistas curtas de validação de
registro checam observações com valores extremos ou frágeis no painel,
não hipóteses de gestão. O grupo Privado, com apenas três hospitais,
fica fora de todos os blocos comparativos pelo tamanho; o painel de 289
\textbf{não tem mais hospitais de longa permanência} --- a ETAPA F
removeu os dezoito que existiam no painel de 314, então essa exclusão,
que valia antes, deixou de ser necessária. As entrevistas são
semiestruturadas: o roteiro-base comum vale para todos os entrevistados
dos blocos de gestão (a--e), e os módulos específicos por perfil
acrescentam perguntas próprias de cada grupo; o módulo de validação de
registro (bloco f) não usa o roteiro-base. Cada pergunta traz, em
itálico, o achado quantitativo que ela testa; o entrevistador deve
conhecer esses achados, mas não deve revelá-los ao entrevistado antes da
resposta.

\parte{O núcleo: hospitais previstos para entrevista}

\bloco{Bloco a: conversores da Direta para OSS (prioritários)}
\begin{center}\small
\begin{tabular}{llllrr}
\toprule
CNES & Hospital & Município & Categoria & DEA & SFA \\
\midrule
2081695 & Conjunto Hospitalar Sorocaba & Sorocaba & OSS (conversor) & 0,563 & 0,999 \\
2078287 & Centro De Referencia Da Saude Da Mulhe & Sao Paulo & OSS (conversor) & 0,827 & 0,805 \\
2082225 & Hospital Katia De Souza Rodrigues Taip & Sao Paulo & OSS (conversor) & 0,818 & 0,912 \\
2091755 & Hospital Geral De Vila Penteado Dr Jos & Sao Paulo & OSS (conversor) & 0,582 & 0,893 \\
2750511 & Hospital Estadual Dr Odilo Antunes De  & Presidente Prudente & OSS (conversor) & 0,698 & 0,900 \\
\bottomrule
\end{tabular}
\end{center}
\comofeito{as seis tabelas de seleção têm as mesmas seis colunas.
\textbf{CNES} é o código do estabelecimento; \textbf{Hospital} e
\textbf{Município} identificam-no; \textbf{Categoria} é o modelo de
gestão; e \textbf{DEA} e \textbf{SFA} são os dois escores de eficiência
técnica do hospital (entre 0 e 1), tomados pela \textbf{mediana dos anos
de 2015 a 2025} --- o DEA vem da análise envoltória de dados e o SFA, da
fronteira estocástica (Seção~\ref{sec:fase2}). \textbf{Ressalva que vale
para todas as seis tabelas:} nas linhas de categoria OSS, o valor SFA de
0,999 é a \textbf{degeneração numérica} descrita na
Seção~\ref{sec:fase2} --- a eficiência técnica das OSS na fronteira
estocástica é praticamente uma constante (0,9993 a 0,9994 em todo o
painel), não uma medida do hospital. Para os hospitais OSS, portanto,
vale apenas a coluna DEA; o número da coluna SFA não deve ser citado nem
comparado com o de outras categorias.}

Os cinco conversores são o material mais valioso da pesquisa, pois
permitem observar o mesmo hospital antes e depois da mudança de gestão.
As \emph{figuras} das suas trajetórias nos quatro indicadores já constam
da Seção~\ref{sec:conversores} (Figuras~\ref{fig:conv2081695}
a~\ref{fig:conv2750511b}) e por isso não se repetem aqui; as fichas
abaixo trazem as medianas antes e depois, e é aqui --- e só aqui --- que
o \textbf{faturamento real por saída (R\$ de dez/2025)} dos cinco
hospitais aparece.

\noindent\textbf{2081695, Conjunto Hospitalar Sorocaba} (Sorocaba, interior; Grande Porte, faixa Barcelona 5, complexidade estrutural 675,0). Conversor da Direta para OSS em 2019. Antes e depois (medianas): TMP de 8,29 para 7,13 dias (queda de 14,0\%); mortalidade geral de 9,34\% para 7,95\% (queda de 14,8\%); faturamento real por saída (R\$ de dez/2025) de R\$ 2.734,63 para R\$ 2.401,47 (queda de 12,2\%); saídas hospitalares de 13.159,50 para 16.243,00 (alta de 23,4\%).\par\medskip
\noindent\textbf{2078287, Centro de Referência da Saúde da Mulher, Pérola Byington} (São Paulo, capital; Médio Porte, faixa Barcelona 4, complexidade estrutural 266,0). Conversor da Direta para OSS em 2023. Antes e depois (medianas): TMP de 4,12 para 3,40 dias (queda de 17,5\%); mortalidade geral de 7,55\% para 4,57\% (queda de 39,5\%); faturamento real por saída (R\$ de dez/2025) de R\$ 1.989,40 para R\$ 1.423,76 (queda de 28,4\%); saídas hospitalares de 5.583,00 para 14.036,00 (alta de 151,4\%).\par\medskip
\noindent\textbf{2082225, Hospital Kátia de Souza Rodrigues, Taipas} (São Paulo, capital; Grande Porte, faixa Barcelona 4, complexidade estrutural 380,0). Conversor da Direta para OSS em 2025, com apenas um ano posterior e sem leitura de tendência. Antes e depois (medianas): TMP de 5,41 para 5,35 dias (queda de 1,1\%); mortalidade geral de 7,90\% para 6,47\% (queda de 18,1\%); faturamento real por saída (R\$ de dez/2025) de R\$ 869,18 para R\$ 977,07 (alta de 12,4\%); saídas hospitalares de 9.804,00 para 10.839,00 (alta de 10,6\%).\par\medskip
\noindent\textbf{2091755, Hospital Geral de Vila Penteado} (São Paulo, capital; Grande Porte, faixa Barcelona 4, complexidade estrutural 350,0). Conversor da Direta para OSS em 2025, com apenas um ano posterior e sem leitura de tendência. Antes e depois (medianas): TMP de 5,03 para 5,50 dias (alta de 9,3\%); mortalidade geral de 7,48\% para 8,31\% (alta de 11,1\%); faturamento real por saída (R\$ de dez/2025) de R\$ 940,08 para R\$ 1.246,13 (alta de 32,6\%); saídas hospitalares de 8.350,50 para 7.783,00 (queda de 6,8\%).\par\medskip
\noindent\textbf{2750511, Hospital Estadual de Presidente Prudente} (Presidente Prudente, interior; Médio Porte, faixa Barcelona 2, complexidade estrutural 94,0). Conversor da Direta para OSS em 2025, com apenas um ano posterior e sem leitura de tendência. Antes e depois (medianas): TMP de 3,72 para 3,91 dias (alta de 5,4\%); mortalidade geral de 0,60\% para 0,68\% (alta de 13,4\%); faturamento real por saída (R\$ de dez/2025) de R\$ 1.132,81 para R\$ 1.082,27 (queda de 4,5\%); saídas hospitalares de 3.675,00 para 3.513,00 (queda de 4,4\%).\par\medskip

\bloco{Bloco b: desviantes positivos (quartil superior persistente)}
\begin{center}\small
\begin{tabular}{llllrr}
\toprule
CNES & Hospital & Município & Categoria & DEA & SFA \\
\midrule
7573162 & Hospital Regional De Jundiai & Jundiai & OSS & 0,905 & 0,999 \\
6164366 & Hospital Estadual Americo Brasiliense & Americo Brasiliense & Direta & 0,810 & 0,908 \\
5718368 & Hosp Mun M Boi Mirim & Sao Paulo & Público Municipal & 0,870 & 0,962 \\
2076896 & Hospital Sao Luiz Gonzaga & Sao Paulo & Filantrópico & 0,922 & 0,928 \\
\bottomrule
\end{tabular}
\end{center}
\comofeito{mesmas seis colunas da tabela do bloco a (CNES, hospital,
município, categoria e os escores medianos DEA e SFA de 2015 a 2025),
com a mesma ressalva: o SFA de 0,999 da linha OSS é degeneração
numérica, não medida, e não deve ser citado nem comparado.}

Um desviante positivo por categoria, todos no suporte comum de porte e
complexidade, escolhidos por manterem eficiência alta ao longo dos onze
anos. A pergunta que eles respondem é quais práticas concretas sustentam
esse giro alto de forma persistente. Em cada ficha abaixo, o
\textbf{marcador financeiro$\times$físico} compara o rank do hospital no
DEA com insumo financeiro (leitos e valor real) contra o rank no DEA só
com leitos: divergência grande é candidata a artefato do regime de
pagamento, não de gestão real, e fica registrada como pergunta em aberto
para a entrevista, não como conclusão.

\noindent\textbf{7573162, Hospital Regional de Jundiaí} (Jundiaí, interior; Médio Porte, faixa Barcelona 3, complexidade estrutural 246,0). Desviante positivo da categoria OSS (quartil superior persistente de eficiência técnica em DEA e SFA): mediana de 0,905 no DEA entre 2015 e 2025, no suporte comum de faixa 3 e Médio Porte; a mediana SFA de 0,999 não é citável (\emph{degeneração OSS}), de modo que a seleção deste hospital se apoia no DEA. Marcador financeiro$\times$físico: 4\textordmasculine{} no rank DEA com insumo financeiro contra 18\textordmasculine{} só com leitos (diferença de 14 posições).\par\medskip
\noindent\textbf{6164366, Hospital Estadual Américo Brasiliense} (Américo Brasiliense, interior; Médio Porte, faixa Barcelona 3, complexidade estrutural 187,0). Desviante positivo da categoria Direta, por critério relaxado (quartil superior de DEA e acima da mediana de SFA): medianas de 0,810 no DEA e 0,908 na SFA entre 2015 e 2025, no suporte comum de faixa 3 e Médio Porte. Marcador financeiro$\times$físico: 27\textordmasculine{} no rank DEA com insumo financeiro contra 75\textordmasculine{} só com leitos (diferença de 48 posições).\par\medskip
\noindent\textbf{5718368, Hospital Municipal M'Boi Mirim} (São Paulo, capital; Grande Porte, faixa Barcelona 4, complexidade estrutural 469,0). Desviante positivo da categoria Público Municipal (quartil superior persistente de eficiência técnica em DEA e SFA): medianas de 0,870 no DEA e 0,962 na SFA entre 2015 e 2025, no suporte comum de faixa 4 e Grande Porte. Marcador financeiro$\times$físico: 12\textordmasculine{} no rank DEA com insumo financeiro contra 4\textordmasculine{} só com leitos (diferença de $-$8 posições --- ranqueia melhor sem o dinheiro, na direção oposta dos outros três hospitais deste bloco).\par\medskip
\noindent\textbf{2076896, Hospital São Luiz Gonzaga} (São Paulo, capital; Grande Porte, faixa Barcelona 3, complexidade estrutural 241,0). Desviante positivo da categoria Filantrópico (quartil superior persistente de eficiência técnica em DEA e SFA): medianas de 0,922 no DEA e 0,928 na SFA entre 2015 e 2025, no suporte comum de faixa 3 e Grande Porte. Marcador financeiro$\times$físico: 1\textordmasculine{} no rank DEA com insumo financeiro contra 66\textordmasculine{} só com leitos (diferença de 65 posições --- a maior do bloco; candidato prioritário para a pergunta de mecanismo).\par\medskip

\noindent A trajetória de cada um destes quatro hospitais nos quatro
indicadores, contra a mediana da sua categoria, está no
Apêndice~\ref{ap:selecao}.

\bloco{Bloco c: desviantes negativos (quartil inferior persistente)}
\begin{center}\small
\begin{tabular}{llllrr}
\toprule
CNES & Hospital & Município & Categoria & DEA & SFA \\
\midrule
2028840 & Instituto De Infectologia Emilio Ribas & Sao Paulo & Direta & 0,447 & 0,409 \\
2698471 & Secao Hospital Municipal Dr Arthur Dom & Santos & Público Municipal & 0,378 & 0,717 \\
2080354 & Hospital Santo Antonio Santos & Santos & Filantrópico & 0,351 & 0,399 \\
\bottomrule
\end{tabular}
\end{center}
\comofeito{mesmas seis colunas da tabela do bloco a (CNES, hospital,
município, categoria e os escores medianos DEA e SFA de 2015 a 2025).}

O espelho do bloco anterior: um hospital por categoria no quartil
inferior persistente de eficiência, para ouvir quais restrições
impedem o giro, na visão de quem opera, antes de confrontá-las com os
dados do painel.

\noindent\textbf{2028840, Instituto de Infectologia Emílio Ribas} (São Paulo, capital; Grande Porte, faixa Barcelona 4, complexidade estrutural 459,0). Desviante negativo da categoria Direta (quartil inferior persistente de eficiência técnica em DEA e SFA): medianas de 0,447 no DEA e 0,409 na SFA entre 2015 e 2025, no suporte comum de faixa 4 e Grande Porte.\par\medskip
\noindent\textbf{2698471, Hospital Municipal Dr. Arthur Domingues Pinto} (Santos, interior; Médio Porte, faixa Barcelona 3, complexidade estrutural 122,0). Desviante negativo da categoria Público Municipal (quartil inferior persistente de eficiência técnica em DEA e SFA): medianas de 0,378 no DEA e 0,717 na SFA entre 2015 e 2025, no suporte comum de faixa 3 e Médio Porte.\par\medskip
\noindent\textbf{2080354, Hospital Santo Antônio} (Santos, interior; Grande Porte, faixa Barcelona 3, complexidade estrutural 126,0). Desviante negativo da categoria Filantrópico (quartil inferior persistente de eficiência técnica em DEA e SFA): medianas de 0,351 no DEA e 0,399 na SFA entre 2015 e 2025, no suporte comum de faixa 3 e Grande Porte.\par\medskip

\noindent As trajetórias destes três hospitais também estão no
Apêndice~\ref{ap:selecao}.

\parte{A reserva sugerida: entrevistas adicionais}

Além do núcleo acima, dois blocos de \textbf{sugestões} foram montados
para o caso de o tempo de campo permitir: os \textbf{pares casados de
contraste} (bloco d), que aprofundam a comparação principal do
documento pondo lado a lado hospitais de estrutura quase idêntica e
eficiência divergente, e os \textbf{casos de discordância entre DEA e
SFA} (bloco e), que ajudam a entender o que os modelos de eficiência
não capturam. São dez hospitais, e como podem não ser entrevistados,
suas tabelas, fichas e figuras ficam no
Apêndice~\ref{ap:selecao}. Os módulos de roteiro correspondentes
(D e E) continuam nesta seção, adiante.

\parte{Validação de registro: as duas entrevistas curtas}

\bloco{Bloco f: validação de registro (entrevistas curtas)}
\begin{center}\small
\begin{tabular}{llllrr}
\toprule
CNES & Hospital & Município & Categoria & DEA & SFA \\
\midrule
2022648 & Hospital Irmaos Penteado E Santa Casa  & Campinas & Filantrópico & 0,354 & 0,664 \\
2097613 & Hospital Infante D Henrique & Sao Jose Do Rio Preto & Filantrópico & 0,758 & 0,264 \\
\bottomrule
\end{tabular}
\end{center}
\comofeito{mesmas seis colunas da tabela do bloco a (CNES, hospital,
município, categoria e os escores medianos DEA e SFA de 2015 a 2025);
estas duas linhas foram escolhidas por problema de registro, não por
eficiência.}

Duas entrevistas curtas (vinte a trinta minutos), só com o responsável
pelo cadastro CNES e faturamento SIH, restritas ao registro --- não são
entrevistas de gestão e não usam o roteiro-base.

\noindent\textbf{2022648, Hospital Irmãos Penteado e Santa Casa de Campinas} (Campinas, interior; Médio Porte, faixa Barcelona 3, complexidade estrutural 156,0). Ocupação de UTI registrada em 2021 chega a 875,2\% (2020: 561,6\%; leitos UTI cadastrados: 16). Verificar leitos CNES cadastrados contra leitos operacionais na pandemia.\par\medskip
\noindent\textbf{2097613, Hospital Infante D. Henrique} (São José do Rio Preto, interior; Grande Porte, faixa Barcelona 2, complexidade estrutural 39,0). Denominadores frágeis em 2020--21 (saídas registradas: 13 em 2020, 4 em 2021; TMP zero em 2021). Verificar se o hospital fechou, mudou de papel ou parou de faturar pelo SIH nesse período; as duas observações já são tratadas como denominador frágil e excluídas dos desfechos em toda a estimação.\par\medskip

\parte{Roteiro-base comum aos blocos de gestão (a--e)}

O roteiro-base cobre seis domínios e dura de cinquenta a sessenta
minutos. Vale para os entrevistados dos blocos a, b e c do núcleo e, se
realizados, dos pares casados sugeridos e dos casos de divergência. As
perguntas são abertas; os itens entre colchetes são estímulos de
aprofundamento para o entrevistador.

\bloco{Governança e desenho institucional}
\begin{enumerate}[itemsep=2pt]
  \item Como se organiza a direção do hospital: quem decide o quê entre a diretoria local, a mantenedora ou OSS e a Secretaria? [organograma real e formal; conselhos; tempo de mandato da direção]
  \par\nopagebreak\achado{a eficiência técnica maior das OSS é traço robusto da rede, com DEA e SFA concordando; a hipótese é que o desenho de governança seja o canal.}
  \item Que decisões precisam de autorização externa ao hospital e quanto tempo levam, do pedido à resposta? [contratação, compra acima de um teto, abertura de leito]
  \par\nopagebreak\achado{o regime operacional mais intenso das OSS sugere menor fricção decisória; a pergunta mede a fricção diretamente.}
  \item Com que frequência a direção examina indicadores assistenciais e o que acontece quando um indicador piora? [rotina de reunião; quem responde pelo indicador]
  \par\nopagebreak\achado{desviantes positivos mantêm eficiência alta por onze anos; persistência sugere rotina gerencial, não sorte.}
\end{enumerate}

\bloco{Autonomia de recursos humanos e compras}
\begin{enumerate}[itemsep=2pt]
  \item Como o hospital contrata e desliga pessoal assistencial, e quanto tempo passa entre a necessidade identificada e o profissional trabalhando? [concurso e CLT; terceirização; banco de horas]
  \par\nopagebreak\achado{a produção das OSS é cerca de 43,7\% maior a estrutura comparável (leitura transversal); a elasticidade de RH é o mecanismo candidato.}
  \item Como funcionam as compras e a manutenção de equipamentos: prazos típicos, o que mais trava? [licitação e compra direta; estoque; equipamento parado]
  \par\nopagebreak\achado{o faturamento real por saída não difere entre categorias; se as OSS produzem mais com faturamento igual, a diferença deve estar no uso dos insumos, não no preço.}
  \item Qual é a rotatividade das equipes médica e de enfermagem, e o que a explica? [vínculos múltiplos; plantões; salário relativo]
  \par\nopagebreak\achado{a estabilidade de equipe é explicação concorrente à gestão para a eficiência persistente dos desviantes positivos.}
\end{enumerate}

\bloco{Contrato de gestão e metas}
\begin{enumerate}[itemsep=2pt]
  \item Existe contrato, convênio ou instrumento com metas quantitativas? Quais indicadores, quem apura, e o que acontece quando a meta não é cumprida? [glosa; aditivo; renegociação]
  \par\nopagebreak\achado{o contrato de gestão é a diferença institucional formal entre OSS e administração direta; a queda de 10,2\% no TMP (within) e a alta de 43,7\% na produção (transversal) são os efeitos que ele teria de explicar.}
  \item As metas chegam ao cotidiano clínico? Dê um exemplo de decisão de ponta (alta, agenda cirúrgica, escala) que mudou por causa de meta. [quem traduz a meta para a equipe]
  \par\nopagebreak\achado{o TMP menor nas OSS é evidência dentro do hospital, apoiada em cinco conversores; a pergunta busca o elo micro entre contrato e permanência.}
\end{enumerate}

\bloco{Relação com a Secretaria e financiamento}
\begin{enumerate}[itemsep=2pt]
  \item Como é a interação com a Secretaria, estadual ou municipal: cobrança, apoio técnico, regulação de fila? Com que frequência? [câmaras técnicas; visitas; sistemas]
  \par\nopagebreak\achado{categorias públicas e OSS respondem à mesma Secretaria com resultados diferentes; a pergunta separa o vínculo formal da relação praticada.}
  \item Como o dinheiro chega e o que acontece numa crise de caixa? [teto MAC; incentivos; atraso de repasse; quem socorre]
  \par\nopagebreak\achado{o faturamento real por saída não difere entre categorias: se o financiamento é semelhante, a diferença de desempenho não é explicada por mais dinheiro.}
\end{enumerate}

\bloco{Gestão de leitos e fluxo do paciente}
\begin{enumerate}[itemsep=2pt]
  \item Descreva a rotina de altas: horário, critérios, quem decide, existe reunião diária de leitos? [alta assistida; previsão de alta na admissão]
  \par\nopagebreak\achado{o TMP cerca de 10\% menor nas OSS torna a rotina de gestão de altas o mecanismo mais direto que a entrevista pode verificar.}
  \item Quem controla a fila de internação e a ocupação do leito (NIR, regulação própria, CROSS)? Como se decide quem entra? [leito de retaguarda; internação social]
  \par\nopagebreak\achado{o giro alto dos desviantes positivos exige controle fino de entrada e saída; a pergunta identifica o instrumento.}
  \item Que parte da cirurgia migrou para o regime ambulatorial ou hospital-dia, e o que impediu de migrar mais? [centro cirúrgico; anestesia; retorno]
  \par\nopagebreak\achado{produção maior com TMP menor é compatível com substituição de internação por ambulatório; a pergunta testa essa composição.}
\end{enumerate}

\bloco{Resposta à pandemia (2020 e 2021)}
\begin{enumerate}[itemsep=2pt]
  \item O que mudou em leitos, equipes e fluxo em 2020 e 2021, e quanto da expansão foi física de fato e quanto foi apenas cadastro? [leitos de campanha; UTI improvisada; registro no CNES]
  \par\nopagebreak\achado{as séries de 2020 e 2021 têm extremos de ocupação de UTI e denominadores frágeis; a distinção entre leito real e leito cadastrado é decisiva para a leitura desses anos.}
  \item O que a pandemia deixou de permanente na operação do hospital? [telessaúde; protocolos; equipes]
  \par\nopagebreak\achado{a estimação de robustez sem 2020 e 2021 preserva os achados principais; a pergunta explora se houve mudança estrutural que os dados posteriores já refletem.}
\end{enumerate}

\parte{Módulos específicos por perfil}

\bloco{Módulo A: conversores da Direta para OSS (bloco a)}

Reconstituir a linha do tempo da transição: quando a conversão começou a
ser discutida, quando o contrato foi assinado, quando a OSS assumiu de
fato. Atenção metodológica: no Pérola Byington e em Sorocaba a melhora
dos indicadores começou \emph{antes} da conversão formal, e é
indispensável perguntar o que já estava em curso no período anterior
(intervenção da Secretaria, troca de direção, obra, mudança de perfil),
pois a resposta separa o efeito da gestão OSS de uma tendência
pré-existente.

\begin{enumerate}[itemsep=2pt]
  \item Conte a linha do tempo da transição para OSS: primeira discussão, decisão, assinatura, assunção. O que aconteceu no hospital em cada etapa? [greves; judicialização; transição de direção]
  \par\nopagebreak\achado{as datas de conversão adotadas (2019, 2023 e três em 2025) vêm do cadastro; a linha do tempo real pode antecedê-las e reinterpretar o estudo de eventos.}
  \item O que mudou concretamente no primeiro ano de OSS: rotinas, chefias, sistemas, escalas? O que não mudou? [primeiros noventa dias; o que a OSS trouxe pronto]
  \par\nopagebreak\achado{no estudo de eventos, o TMP cai já nos primeiros anos posteriores; a pergunta identifica qual prática explica a queda rápida.}
  \item Quanto da equipe assistencial e das chefias permaneceu após a conversão? [estatutários cedidos; recontratação CLT]
  \par\nopagebreak\achado{a continuidade de equipe distingue efeito de gestão de efeito de composição: se a equipe é a mesma, a mudança é de método, não de pessoas.}
  \item Quais metas o contrato de gestão fixou no início e como elas mudaram? Alguma meta era vista como inatingível? [metas de produção e de qualidade]
  \par\nopagebreak\achado{a alta de 43,7\% na produção (transversal) e a queda de 10,2\% no TMP (within) fazem das metas contratuais a explicação institucional candidata.}
  \item (Somente Pérola Byington e Sorocaba) Os indicadores já vinham melhorando antes da conversão formal. O que estava acontecendo no hospital nos dois ou três anos anteriores? [gestor interino; plano de saneamento; antecipação da OSS]
  \par\nopagebreak\achado{a melhora pré-conversão observada nesses dois hospitais ameaça a leitura causal do estudo de eventos; a resposta decide entre antecipação, tendência ou seleção.}
\end{enumerate}

\bloco{Módulo B: desviantes positivos (bloco b)}

O objetivo é extrair práticas concretas e transferíveis por trás do giro
alto e da eficiência persistente, com ceticismo ativo: pedir sempre o
exemplo, o documento, a rotina, sem aceitar declaração de princípios.

\begin{enumerate}[itemsep=2pt]
  \item O hospital está no quartil superior de eficiência da rede há onze anos. Na sua leitura, quais três práticas explicam isso? Para cada uma: desde quando, quem implantou, o que custou? [pedir exemplos datados]
  \par\nopagebreak\achado{a persistência no quartil superior em DEA e SFA simultaneamente afasta artefato de método; deve haver prática real por trás.}
  \item Como funciona, passo a passo, a gestão de altas num dia típico? [reunião de leitos; previsão de alta; alta aos sábados]
  \par\nopagebreak\achado{giro alto com TMP baixo é a assinatura operacional dos desviantes positivos.}
  \item Que procedimentos o hospital faz em regime ambulatorial que vizinhos comparáveis internam? [cirurgia ambulatorial; hospital-dia; protocolos]
  \par\nopagebreak\achado{a produção alta com leitos comparáveis sugere substituição ambulatorial; a pergunta lista os procedimentos.}
  \item Como o hospital se relaciona com a regulação de vagas: aceita tudo que a central manda? Recusa o quê, e com que argumento? [seleção de casos; perfil de demanda]
  \par\nopagebreak\achado{a eficiência alta pode ser gestão ou seleção de pacientes mais leves; a pergunta confronta a hipótese de seleção, a ler junto com a complexidade estrutural do hospital e com o marcador financeiro$\times$físico da ficha de cada hospital.}
\end{enumerate}

\bloco{Módulo C: desviantes negativos (bloco c)}

Simétrico ao módulo B, mas centrado nas restrições percebidas. O tom não
é de auditoria: o objetivo é entender o que impede o giro, na visão de
quem opera, e confrontar depois com as restrições objetivas do painel.

\begin{enumerate}[itemsep=2pt]
  \item Quais são hoje os três maiores obstáculos para o hospital girar mais os leitos? [financiamento; RH; perfil de demanda; física do prédio]
  \par\nopagebreak\achado{o hospital está no quartil inferior persistente de DEA e SFA; a pergunta colhe a explicação nativa antes de confrontá-la com os dados.}
  \item Descreva o paciente típico que fica internado mais tempo do que precisaria. Por que ele não sai? [internação social; retaguarda; exame que não sai]
  \par\nopagebreak\achado{TMP alto com ocupação alta indica leito ocupado por paciente que não precisaria dele; a pergunta identifica o gargalo.}
  \item O financiamento cobre o custo do que o hospital produz? Onde aperta primeiro? [tabela SUS; folha; incentivos]
  \par\nopagebreak\achado{o faturamento real por saída não difere entre categorias na rede; se a restrição percebida é financeira, ela precisa ser conciliada com esse achado.}
  \item Se o hospital pudesse mudar uma única regra (de RH, compra ou contrato), qual mudaria e o que aconteceria? [pergunta projetiva]
  \par\nopagebreak\achado{a resposta revela qual margem de autonomia o gestor acredita que falta, comparável em espelho com o que as OSS dizem ter.}
\end{enumerate}

\bloco{Módulo D: pares casados (bloco d, sugeridos)}

As perguntas dos domínios de recursos humanos e compras, contrato e
metas, e leitos e fluxo do roteiro-base devem ser feitas de forma
espelhada e com a mesma ordem nos dois hospitais do par, para permitir
comparação direta das respostas. Acrescentar:

\begin{enumerate}[itemsep=2pt]
  \item Num dia típico, quem decide a abertura de um leito extra, a contratação de um plantonista e a compra urgente de um insumo, e em quanto tempo? [responder para os três casos]
  \par\nopagebreak\achado{o par tem faixa Barcelona, porte e leitos parecidos com eficiência divergente; a velocidade decisória é a diferença de gestão mais plausível com estrutura igual.}
  \item Como é montada a escala médica do fim de semana e o que acontece com as altas de sexta a domingo? [alta no fim de semana; plantão de sexta]
  \par\nopagebreak\achado{diferenças de TMP entre pares estruturalmente parecidos costumam se concentrar no fim de semana; pergunta espelhada nos dois hospitais do par.}
\end{enumerate}

\bloco{Módulo E: divergência entre DEA e SFA (bloco e, sugeridos)}

A divergência entre os dois métodos sugere algo que os modelos não
capturam: composição de casos atípica, insumo mal medido, produção
heterogênea. A entrevista busca esse fator omitido.

\begin{enumerate}[itemsep=2pt]
  \item O que este hospital faz que hospitais parecidos não fazem: serviços, perfis de paciente, papéis na rede que não aparecem em leitos e saídas? [ensino; referência regional; porta aberta e referenciada]
  \par\nopagebreak\achado{DEA e SFA discordam fortemente sobre este hospital; a explicação provável é dimensão de produto ou insumo que nenhum dos dois modelos mede.}
  \item Os números oficiais de leitos, produção e permanência do hospital refletem a operação real? Onde o registro distorce? [leitos bloqueados; produção não faturada]
  \par\nopagebreak\achado{a divergência entre métodos também pode ser artefato de medida; a pergunta verifica a qualidade do dado antes de interpretar o resto da entrevista.}
\end{enumerate}

\bloco{Módulo F: validação de registro (bloco f)}

Entrevistas curtas, restritas a quem responde pelo cadastro CNES e pelo
faturamento SIH --- não usam o roteiro-base nem os módulos de gestão.

\begin{enumerate}[itemsep=2pt]
  \item (CNES 2022648) Em 2021 a ocupação de UTI calculada chega a 875\%. Quantos leitos de UTI operavam de fato em 2020 e 2021, e quantos estavam cadastrados no CNES? Houve atraso ou represamento de atualização cadastral? [leitos de campanha; UTI COVID]
  \par\nopagebreak\achado{ocupação acima de 100\% em escala implica denominador (leitos CNES) menor que a operação real; a resposta define como tratar as observações extremas de 2020 e 2021.}
  \item (CNES 2097613) O registro mostra menos de 30 saídas sem COVID no biênio 2020--2021 e TMP zero em 2021. O hospital fechou, mudou de papel ou parou de faturar pelo SIH nesse período? [interdição; reforma; migração de convênio]
  \par\nopagebreak\achado{as duas observações desse CNES já são tratadas como denominador frágil e excluídas dos desfechos; a resposta confirma ou corrige esse tratamento.}
\end{enumerate}

\parte{Logística e esforço}

A seleção soma \textbf{24 hospitais únicos}: \textbf{22 entrevistas
plenas} (roteiro-base de cinquenta a sessenta minutos mais o módulo
específico de vinte a trinta minutos, blocos a a e) e \textbf{2
entrevistas curtas} de validação de registro (bloco f, vinte a trinta
minutos). O núcleo previsto (blocos a, b e c) soma doze entrevistas
plenas; a reserva sugerida (blocos d e e, no
Apêndice~\ref{ap:selecao}) acrescenta até dez. O teto
recomendado de campo é de 15 a 18 entrevistas. Havendo necessidade de
priorizar ou cortar, a ordem recomendada é: bloco a (obrigatório), bloco
f (curtas, custo baixo), blocos b e c, cujo contraste entre quartil
superior e inferior é o coração da fase, bloco d, mantendo ao menos um
par completo por âncora e preferindo o par Direta, que espelha a
comparação principal da estimação, e por último bloco e. Recomenda-se
gravar com consentimento, transcrever integralmente e codificar as
respostas contra a lista de achados deste roteiro, para que a fase
qualitativa feche o ciclo: cada mecanismo declarado nas entrevistas deve
ser reexaminável no painel. Toda a lista, no entanto, permanece
subordinada ao ponto de partida desta seção: a escolha final dos
hospitais e das perguntas é da pesquisadora principal.

\clearpage
\section{Síntese e Conclusão Geral}
\label{sec:conclusao}

Este documento perseguiu uma única pergunta ao longo de todas as suas
seções: entre os cinco arranjos institucionais que administram os
hospitais do SUS no estado de São Paulo (\textbf{Filantrópico, Público
Municipal, OSS, Administração Direta e Privado}), algum entrega
desempenho sistematicamente melhor nos cinco indicadores oficiais do
estudo? A resposta foi construída sobre o painel definitivo descrito na
Seção~\ref{sec:painel}: \textbf{289 hospitais acompanhados ano a ano de
2015 a 2025, 3.179 observações balanceadas}, destilado por um funil
auditável a partir de 830 estabelecimentos da base bruta do SIH/SUS. A
conclusão honesta é que a resposta depende radicalmente de qual pergunta
se faz, se quem opera mais intensamente, quem é mais eficiente ou quem
\emph{causa} melhor desfecho ao mudar de gestão, e que a evidência é
forte para umas e frágil para outras. Esta seção amarra as duas metades
do trabalho: a leitura descritiva do painel (Seções~\ref{sec:painel}
a~\ref{sec:qualidade}) e a leitura inferencial e de eficiência
(Seções~\ref{sec:estrategia} a~\ref{sec:fase2}).

\textbf{O que a fase descritiva estabeleceu.} A metade exploratória do
documento mapeou o terreno antes de qualquer modelo. Fixou os moldes
estatísticos de cada indicador (Seção~\ref{sec:distribuicoes}),
retratou o perfil bruto de cada categoria de gestão em medianas e
violinos (Seção~\ref{sec:categorias}) e mostrou que as redes diferem
tanto em desempenho quanto naquilo que as torna comparáveis, a
complexidade e o porte (Seção~\ref{sec:complexidadeporte}). Duas
ressalvas atravessam toda essa leitura e valem para o documento
inteiro. A primeira é que a categoria de gestão de cada hospital segue
a classificação assistencial do SIH adotada operacionalmente pela
equipe (a variável \texttt{modelo\_gestao\_proxy}), e \textbf{não um
cruzamento institucional depurado} caso a caso: é uma definição de
trabalho, não um cadastro jurídico. A segunda é que todo contraste
descritivo entre categorias reflete o \textbf{desempenho médio de quem
está em cada rede}, e não o efeito causal limpo do modelo de gestão,
pois hospitais não são sorteados para uma administração ou outra, e as
redes diferem de partida em porte, complexidade e vocação assistencial.
Os números descritivos situam; não julgam.

\textbf{O que a fase inferencial sustenta, por ordem de robustez.} A
metade causal e de eficiência do documento pode ser resumida não pela
ordem em que os métodos foram aplicados, mas pela \textbf{firmeza da
evidência que cada achado resistiu a produzir}. O que sobrevive a todos
os testes é o \textbf{regime assistencial mais intenso das OSS}: mais
produção, maior ocupação e maior uso de UTI (19,3\% acima da Direta,
p~igual a 0,003, na Seção~\ref{sec:inferenciafinal}), um contraste
transversal firme entre as redes, visível na produção física por leito
que não usa dinheiro nenhum (55 saídas por leito nas OSS contra 36 na
Direta) e replicado em todos os métodos, inclusive sem os conversores.
A esse traço soma-se a \textbf{maior eficiência produtiva das OSS}
(Seção~\ref{sec:fase2}), cerca de 25\% menos distância à fronteira que
a Direta, concordante entre DEA e SFA e igualmente estável à retirada
dos conversores (o coeficiente vai de 0,293 para 0,331 negativo, com
intervalo inteiro abaixo do zero): como a intensidade operacional, é uma característica
\emph{da rede}, medida na comparação entre hospitais distintos, e não um
efeito de conversão.

\textbf{O que é sugestivo, mas não definitivo}, é a queda de cerca de
\textbf{10\% no tempo médio de permanência} associada à gestão por OSS.
Diferentemente dos traços anteriores, esse é um efeito identificado
\emph{dentro} do hospital, pelos cinco conversores que trocaram a
Direta pela OSS (Seções~\ref{sec:efeitogestao}
e~\ref{sec:inferenciafinal}). Ele é consistente em direção e tamanho e
\textbf{sobrevive a todos os testes que têm poder para avaliá-lo}: o
teste de tendências anteriores (p~igual a 0,254), o estimador de
Callaway e Sant'Anna e a reestimação sem o biênio pandêmico. Nos três
testes cujo poder é limitado pelo número de hospitais tratados ele
\textbf{fica sem corroboração}: permutação (p~igual a 0,153), bootstrap
selvagem por cluster (p~igual a 0,324) e controle sintético, cuja
aplicação a este indicador é degenerada (p~de 0,78 e 0,75). Nenhum dos
três o contradiz; nenhum o sustenta --- com cinco hospitais tratados,
eles não distinguem um efeito nulo de uma ausência de poder para
detectá-lo. É evidência sugestiva forte, não fato demonstrado. \textbf{O que
não se sustenta} é a queda de mortalidade após a conversão para OSS.
Visível nos dados brutos dos conversores, ela \textbf{não encontra
apoio no controle sintético}: em Sorocaba a lacuna frente ao gêmeo é
nula (0,54 ponto abaixo, p~igual a 0,67) e o próprio sinal inverte
conforme o painel usado no pool de doadores (era 1,09 acima antes da
ETAPA F) — fragilidade, não efeito; no Pérola Byington a queda existe mas
fica dentro do que sorteios de placebo produzem (p~igual a 0,32). A
Seção~\ref{sec:fase2} deu a esse ponto a demonstração mais direta:
retirando os cinco conversores da amostra, a diferença de mortalidade
entre as redes OSS e Direta simplesmente desaparece. A mortalidade
menor é um \textbf{fenômeno da trajetória dos hospitais que mudaram de
gestão, não uma diferença estrutural entre as redes}. \textbf{O que não
mostrou diferença em método nenhum} foi o faturamento real por saída, igual
entre as redes na leitura transversal (Seção~\ref{sec:resultados}) e na
hierárquica (Seção~\ref{sec:fase2}), e comum entre os placebos no teste
de permutação. \textit{Ressalva de leitura (regime de pagamento): hospitais OSS são custeados por contrato de gestão com orçamento global, não por reembolso de AIH --- o valor registrado por saída é, para eles, produção sem consequência financeira direta; a diferença de faturamento pode refletir o próprio regime de pagamento, a variável cujo efeito se busca medir, e não uso de recursos ou qualidade assistencial.} Sobre o grupo \textbf{Privado} (n~igual a 3), a regra do
documento foi mantida do início ao fim: entrou em todos os modelos
apenas como \textbf{controle de absorção}, para não contaminar as
demais categorias, e nunca recebeu interpretação substantiva, pois a
Seção~\ref{sec:fase2} mostrou seu efeito aparente encolher quase 90\%
assim que o modelo pondera a evidência pelo tamanho do grupo.

\textbf{Limitações do projeto como um todo.} Para além das ressalvas
técnicas discutidas seção a seção, quatro limitações operam no nível do
projeto e condicionam toda leitura dos resultados. A primeira, e a mais
importante, é que \textbf{toda a identificação causal repousa em apenas
cinco conversores}, três deles com um único ano de dados após a mudança
de gestão, uma base estreita para qualquer afirmação de efeito, por
mais cuidadosa que seja a metodologia. A segunda é a \textbf{ausência de
DRG no SIH}: a gravidade do caso, que idealmente ajustaria a comparação,
foi suprida por uma métrica própria de complexidade (de inspiração
Barcelona), cuja versão ponderada por mortalidade
(\texttt{complexidade\_pond\_mort}) permanece provisória. A terceira é
que o painel, por construção, retém apenas hospitais com mediana de
leitos superior a 50, o que \textbf{limita a generalização para
hospitais menores}, ausentes do estudo. A quarta é a \textbf{janela
temporal ainda curta} para os hospitais convertidos em 2025, cujo
pós-tratamento se resume a um único ano e não permite ver a trajetória
se estabilizar.

\textbf{Contribuição.} Com essas fronteiras claras, o que o projeto
entrega é, ao que se sabe, o \textbf{primeiro comparativo institucional
abrangente do estado de São Paulo}, e não apenas da Região
Metropolitana, a reunir, sobre o mesmo painel, três camadas
metodológicas que raramente convivem: a \textbf{estimação causal}
(efeitos aleatórios correlacionados de Mundlak, estudo de eventos,
o estimador de Callaway e Sant'Anna e controle sintético), a
\textbf{análise de eficiência} (DEA com correção de viés e fronteira
estocástica) e a \textbf{validação por modelos bayesianos
hierárquicos}. Mais do que o arsenal, a contribuição está no cuidado
metodológico explícito de \textbf{não confundir efeito de gestão com
efeito de trajetória de conversão}, distinção que, aplicada à
mortalidade, inverteu a leitura ingênua dos dados brutos e é
provavelmente o achado mais consequente do trabalho.

\textbf{Direções futuras.} Três frentes dão continuidade natural à
pesquisa. A primeira é a \textbf{fase qualitativa de entrevistas}, cuja
seleção de entrevistados e roteiro completo já constam da
Seção~\ref{sec:qualitativa}, com o material de apoio no
Apêndice~\ref{ap:selecao}: as entrevistas semiestruturadas ali
desenhadas são dirigidas a investigar o mecanismo por trás do efeito no
tempo de permanência, pois, se ele é real, algo na gestão por OSS
encurta a internação, e entender o quê é o passo que os dados
administrativos não alcançam. A segunda é o acúmulo de \textbf{novos
anos de dados}, sobretudo para os três hospitais convertidos em 2025, o
único caminho para transformar a evidência dentro do hospital, hoje
sugestiva, em conclusão. A terceira é a \textbf{ratificação da forma
funcional da complexidade}, a escolha entre \texttt{complexidade\_pond\_mort}
e o fator de segurança \(s(M)\), ainda pendente de decisão junto à
pesquisadora principal, e que afeta o ajuste de gravidade em todas as
análises de mortalidade.

\textbf{Uma palavra final de cautela.} Os resultados aqui reunidos
foram construídos para serem lidos em conjunto e com seu contexto
metodológico completo. Isolado desse contexto, qualquer número deste
documento é vulnerável ao uso indevido, e dois pontos exigem alerta
expresso. Primeiro, \textbf{nada aqui autoriza a afirmação de que a
gestão por OSS reduz a mortalidade}: o controle sintético não encontra
descolamento significativo em nenhum dos dois conversores testados, a
comparação entre redes sem os conversores é nula, e citar a queda
bruta seria trair o próprio achado central do trabalho. Segundo, toda a
identificação causal apoia-se em \textbf{cinco hospitais convertidos},
uma base estatística frágil por natureza. Por isso, os achados deste
documento \textbf{não devem ser usados isoladamente por gestores
públicos para justificar decisões de modelo de gestão} sem o contexto
metodológico que os cerca. O que a pesquisa oferece com segurança é um
retrato: as OSS operam um regime mais intenso e mais eficiente; o que
ela expressamente não oferece é um veredito causal sobre qual modelo
salva mais vidas.


\clearpage
\appendix
\section{Diagnósticos de ajuste: escolha dos moldes, gráficos QQ e
resíduos quantílicos}
\label{ap:diagnosticos}

\setcounter{figure}{0}
\renewcommand{\thefigure}{A.\arabic{figure}}
\setcounter{table}{0}
\renewcommand{\thetable}{A.\arabic{table}}

Este apêndice reúne todo o material estatístico de ajuste do documento
--- a escolha da família de distribuição de cada indicador, as notas de
qualidade de ajuste, os gráficos QQ e os resíduos quantílicos dos
modelos ---, para que o corpo do texto flua sem interrupções técnicas.
Ele é dirigido a quem quiser conferir a base metodológica; a leitura
substantiva não depende dele. Três noções bastam para acompanhar tudo.
\textbf{Nota AIC}: uma medida de qualidade de ajuste usada para comparar
moldes de distribuição entre si; quanto menor o AIC, melhor o molde
descreve os dados, e o valor isolado não tem interpretação --- só serve
para comparar. \textbf{Gráfico QQ}: compara os dados reais com um molde
matemático de distribuição; se os pontos seguem a reta, o molde descreve
bem os dados, e desvios sistemáticos da reta indicam onde o molde falha.
O número r que acompanha cada QQ resume essa aderência: quanto mais
próximo de 1, melhor. \textbf{Resíduo quantílico}: é a versão do QQ
para um modelo já ajustado com covariáveis; se o modelo descreve bem
os dados, esses resíduos se comportam como uma Normal padrão e os
pontos deitam na diagonal.

\subsection*{Escolha do molde de cada indicador}

Faturamento por saída e TMP são valores positivos com cauda à direita, e
\textbf{cinco moldes clássicos} de distribuição foram postos à prova
para cada um, comparados pela nota AIC e pelo gráfico QQ. O resultado é
consistente (Tabela~\ref{tab:familias}): o molde \textbf{log logístico
(também chamado de Fisk)} vence nas duas amostras, com a \textbf{LogNormal}
logo atrás, e a Gama, molde usado por tradição em análises de custo,
fica sempre entre as piores. No TMP o ajuste do log logístico é
praticamente perfeito (r de 0,9974 na Figura~\ref{fig:qqtmpsemlp}),
porque a consolidação do painel já removeu os hospitais de longa
permanência que, na versão anterior, criavam um teto de 30,5 dias que
nenhum molde contínuo reproduzia. Tanto no faturamento (r de 0,9700,
Figura~\ref{fig:qqcusto}) quanto no TMP (Figura~\ref{fig:qqtmp}) o molde
log logístico acompanha os dados por quase toda a extensão.

\begin{table}[H]
\centering
\caption{Comparação dos cinco moldes de distribuição para faturamento e TMP
pela nota AIC e pelo r do gráfico QQ.}
\label{tab:familias}
\small
\begin{tabular}{llrrr}
\toprule
Amostra & Molde & $n$ & AIC & r do QQ \\
\midrule
Faturamento por saída & Log logística (Fisk) & 3.179 & 50.785 & 0,9700 \\
                & LogNormal            & 3.179 & 50.883 & 0,9079 \\
                & Gaussiana inversa    & 3.179 & 50.952 & 0,8943 \\
                & Gama                 & 3.179 & 51.867 & 0,8310 \\
                & Weibull              & 3.179 & 52.320 & 0,8343 \\
\midrule
TMP & Log logística (Fisk) & 3.178 & 11.878 & 0,9974 \\
                     & LogNormal            & 3.178 & 11.956 & 0,9810 \\
                     & Gaussiana inversa    & 3.178 & 11.989 & 0,9786 \\
                     & Gama                 & 3.178 & 12.150 & 0,9631 \\
                     & Weibull              & 3.178 & 12.920 & 0,9398 \\
\bottomrule
\end{tabular}
\end{table}

\comofeito{cada bloco de linhas é uma amostra (faturamento por saída e
TMP) e cada linha, um molde de distribuição ajustado a ela. A coluna
\textbf{$n$} é o número de observações usadas; \textbf{AIC} é a nota de
ajuste obtida do ajuste de máxima verossimilhança de cada molde (quanto
menor, melhor); e \textbf{r do QQ} é a correlação entre os pontos do
gráfico QQ e a reta ideal (quanto mais perto de 1, melhor). Em cada
amostra os moldes aparecem ordenados do melhor para o pior AIC.}

Para a fração de alta complexidade, a Figura~\ref{fig:betapct} mostra o
ajuste Beta do miolo (os hospitais que realizam alguma alta
complexidade); os dois acúmulos nas pontas --- zero e um --- são
tratados separadamente pelo modelo ZOIB citado na
Seção~\ref{sec:distribuicoes}.

\comoler{a grande maioria dos hospitais que fazem alta complexidade faz
pouca (a concentração à esquerda; metade fica abaixo de 3,7\%), e uma
minoria é altamente especializada (as barras à direita, perto de
100\%). A curva escura é o molde Beta, que acompanha essa queda
contínua.}

\begin{figure}[H]
\centering
\includegraphics[width=.62\textwidth]{fig_ae_02_beta_pct_alta_complex.png}
\caption{Fração de alta complexidade entre os hospitais que realizam
alguma alta complexidade: histograma e molde Beta ajustado. O eixo
vertical está em escala logarítmica para a cauda inteira ficar
visível.}
\label{fig:betapct}
\end{figure}

\comofeito{o eixo horizontal é a fração de alta complexidade, restrita
aos hospitais que realizam alguma (entre 0 e 1), e o vertical acompanha
a quantidade de observações por faixa, em escala logarítmica (para a
cauda inteira ficar visível).
Sobre o histograma foi desenhada a curva do molde Beta ajustado a esses
dados.}

\subsection*{Gráficos QQ da fase exploratória}

As Figuras~\ref{fig:qqcusto} a~\ref{fig:qqtmpsemlp} avaliam os moldes
marginais de faturamento e TMP da Seção~\ref{sec:distribuicoes}.

\comoler{quanto mais os pontos abraçam a reta, melhor o molde. No
painel (b), o molde log logístico segue os dados por quase toda a
extensão, descolando apenas nos pouquíssimos hospitais mais caros; por
isso ele recebe a melhor nota da Tabela~\ref{tab:familias}.}

\begin{figure}[H]
\centering
\subfig{fig_ae_03_qq_lognormal_custo_saida.png}{Molde LogNormal}
\subfig{fig_ae_03_qq_fisk_custo_saida.png}{Molde log logístico (Fisk)}
\caption{Gráficos QQ do faturamento por saída: os pontos são os dados e a
reta é o comportamento esperado pelo molde.}
\label{fig:qqcusto}
\end{figure}

\comofeito{gráfico QQ. O eixo horizontal traz os quantis teóricos do
molde de distribuição e o vertical, os valores do faturamento por saída
ordenados do menor ao maior; a reta é o comportamento esperado se o
molde descrevesse os dados perfeitamente. Cada painel testa um molde
(LogNormal e log logístico/Fisk), e o r no título resume a aderência.}

\comoler{nos dois painéis os pontos acompanham a reta em quase toda a
extensão e descolam apenas nos pouquíssimos hospitais de permanência
mais longa. O painel corrente já não contém a subpopulação de longa
permanência que, na versão anterior desta análise, criava um teto de
30,5 dias e o característico achatamento no alto: a ETAPA F a removeu, e
o ajuste do TMP passou a ser quase perfeito.}

\begin{figure}[H]
\centering
\subfig{fig_ae_03_qq_lognormal_tmp.png}{Molde LogNormal}
\subfig{fig_ae_03_qq_fisk_tmp.png}{Molde log logístico (Fisk)}
\caption{Gráficos QQ do TMP no painel completo.}
\label{fig:qqtmp}
\end{figure}

\comofeito{mesma construção da figura anterior (quantis teóricos no eixo
horizontal, valores ordenados no vertical, reta de referência), agora
para o TMP no painel completo, comparando os moldes LogNormal e log
logístico.}

\comoler{este é o gráfico QQ do TMP no painel atual: os pontos seguem a
reta e o molde descreve os dados quase perfeitamente. Não há mais o
achatamento que a versão anterior do painel exibia, porque os hospitais
de longa permanência que o causavam saíram na consolidação.}

\begin{figure}[H]
\centering
\includegraphics[width=.62\textwidth]{fig_ae_03_qq_fisk_tmp_sem_longa_perm.png}
\caption{Gráfico QQ do TMP no painel de 289 hospitais (log logística):
os pontos ficam praticamente sobre a reta (r de 0,9974).}
\label{fig:qqtmpsemlp}
\end{figure}

\comofeito{gráfico QQ do TMP no painel de 289 hospitais, molde log
logístico: quantis teóricos no eixo horizontal, valores ordenados do TMP
no vertical, reta de referência. Quanto mais os pontos ficam sobre a
reta, melhor o ajuste (aqui r de 0,9974).}

\subsection*{Os moldes melhoram quando se compara o comparável}

Uma pergunta natural: se os QQ do painel completo não são perfeitos, o
molde está errado? As Figuras~\ref{fig:qqcondA} a~\ref{fig:qqcondC}
respondem que não: o desajuste vem de misturar 289 hospitais muito
diferentes num único gráfico. Quando se compara cada hospital consigo mesmo (controlando o
efeito do hospital e do ano), os desvios praticamente desaparecem: no
faturamento, a assimetria restante é quase nula; na mortalidade, a forte
assimetria da distribuição bruta na escala \textbf{logito} (a escala
interna em que o modelo trabalha proporções), de 1,97 negativa, cai
para 0,27 negativa depois do controle; na escala original da
Tabela~\ref{tab:univariada} a assimetria é outra (0,80 positiva), e os
dois números não se contradizem. Em resumo, \textbf{os moldes escolhidos descrevem bem os
dados quando se compara o comparável}.

\comoler{compare o gráfico de cima com o de baixo: em cima, todos os
hospitais juntos, os pontos serpenteiam em volta da reta; embaixo,
comparando cada hospital consigo mesmo, os pontos praticamente deitam
na reta.}

\begin{figure}[H]
\centering
\subfig{fig_ae_11_qq_custo_marginal.png}{Faturamento: todos os hospitais
juntos}
\subfig{fig_ae_11_qq_custo_residuo_ef.png}{Faturamento: cada hospital
comparado consigo mesmo}
\caption{Faturamento por saída: gráfico QQ antes e depois de controlar as
diferenças fixas entre hospitais e entre anos.}
\label{fig:qqcondA}
\end{figure}

\comofeito{dois gráficos QQ do faturamento (quantis teóricos no eixo
horizontal, valores ordenados no vertical). No painel de cima, todos os
289 hospitais entram juntos; no de baixo, controla-se o efeito fixo de
hospital e de ano, isto é, cada hospital é comparado consigo mesmo ao
longo do tempo. Comparar os dois isola se o desajuste vem do molde ou da
mistura de hospitais diferentes.}

\comoler{no painel anterior este par contrastava o desenho em degrau
(com os hospitais de longa permanência) com a reta quase perfeita (sem
eles). No painel corrente a ETAPA F já removeu essa subpopulação, então
as duas versões coincidem e ambas mostram a reta --- a confirmação de
que o desvio era a subpopulação, não o molde.}

\begin{figure}[H]
\centering
\subfig{fig_ae_11_qq_tmp_marginal.png}{TMP: todos os hospitais juntos}
\subfig{fig_ae_11_qq_tmp_marginal_sem_longa_perm.png}{TMP: sem a antiga
subpopulação de longa permanência}
\caption{TMP: gráfico QQ do painel completo e da versão sem a antiga
subpopulação de longa permanência --- no painel corrente as duas
coincidem, porque a ETAPA F já removeu essa subpopulação.}
\label{fig:qqcondB}
\end{figure}

\comofeito{mesma construção da figura anterior, para o TMP: o painel de
cima traz o painel completo e o de baixo, a versão sem a antiga
subpopulação de longa permanência (no painel atual as duas coincidem,
pois a ETAPA F já removeu essa subpopulação).}

\comoler{mesmo padrão do faturamento: a curva em S do gráfico de cima, com
todos os hospitais misturados, vira uma reta bem comportada quando cada
hospital é comparado consigo mesmo. A heterogeneidade entre hospitais,
e não o molde, era a fonte do desajuste.}

\begin{figure}[H]
\centering
\subfig{fig_ae_11_qq_mort_logito_marginal.png}{Mortalidade: todos os
hospitais juntos}
\subfig{fig_ae_11_qq_mort_logito_residuo_ef.png}{Mortalidade: cada
hospital comparado consigo mesmo}
\caption{Mortalidade: gráfico QQ antes e depois de controlar as
diferenças fixas entre hospitais e entre anos.}
\label{fig:qqcondC}
\end{figure}

\comofeito{mesma construção (QQ com todos os hospitais juntos em cima e
cada hospital comparado consigo mesmo embaixo), agora para a mortalidade
na escala logito, a escala interna em que o modelo trabalha proporções.}

\subsection*{Resíduos quantílicos dos modelos estimados}

A fase de estimação prometeu que a escolha das famílias de distribuição
seria validada nos modelos ajustados, e é o que as
Figuras~\ref{fig:qqprimeira} e~\ref{fig:qqsegunda} fazem, com a
Figura~\ref{fig:qqgama} documentando a reestimação da UTI comentada na
Seção~\ref{sec:inferenciafinal}. O resíduo quantílico é um teste de
qualidade: se o modelo descreve bem os dados, esses resíduos se alinham
sobre a diagonal do gráfico. O alinhamento é muito bom na mortalidade
(correlação de 0,976), na fração de alta complexidade (0,991) e na
produção (0,967; o piso achatado na ponta esquerda do QQ, com
observações empilhadas perto de menos 5, corresponde a hospitais
produzindo muito abaixo do que porte e perfil preveriam, desvio
confinado à cauda extrema que não compromete o contraste entre
categorias), e bom no faturamento (0,970) e no tempo de permanência
(0,952), cujas caudas pesadas já estavam diagnosticadas e são absorvidas
pelas margens de erro agrupadas. A exceção honesta é a
\textbf{intensidade de ocupação da UTI} (0,814): a cauda de UTIs com uso
quase nulo escapa do modelo em logaritmo, e esse resultado específico
deve ser lido com reserva; fica registrado para a próxima iteração o uso
de uma família com massa perto de zero para esse indicador --- reestimada
com a família Gama na Figura~\ref{fig:qqgama}.

\comoler{quanto mais os pontos azuis acompanham a linha tracejada,
melhor o modelo descreve os dados; o valor de r no canto é essa
aderência resumida em um número (1 seria perfeito). Compare o
alinhamento quase exato da fração de alta complexidade com o
descolamento da cauda esquerda na UTI, que motivou a reestimação da
figura seguinte.}

\begin{figure}[H]
\centering
\subfig{fig_est_04_qq_mort_zib.png}{Mortalidade geral, modelo Beta com
inflação em zero}
\subfig{fig_est_04_qq_pct_alta_zoib.png}{Fração de alta complexidade,
modelo ZOIB}
\caption{Resíduos quantílicos dos modelos de taxa.}
\label{fig:qqprimeira}
\end{figure}

\comofeito{gráficos de resíduos quantílicos de dois modelos já
estimados (mortalidade, Beta com inflação em zero; alta complexidade,
ZOIB). O eixo horizontal traz os quantis teóricos da Normal e o
vertical, os resíduos quantílicos ordenados; se o modelo descreve bem os
dados, os pontos deitam sobre a diagonal. O r no canto resume a
aderência.}

\begin{figure}[H]
\centering
\subfig{fig_est_04_qq_producao_bn.png}{Produção, Binomial Negativa}
\subfig{fig_est_04_qq_ocup_uti_pos.png}{Ocupação de UTI entre os
positivos, o único ajuste insatisfatório}
\caption{Resíduos quantílicos da produção e da parte positiva da UTI.}
\label{fig:qqsegunda}
\end{figure}

\comofeito{mesma construção da figura anterior (quantis teóricos da
Normal no eixo horizontal, resíduos quantílicos ordenados no vertical,
diagonal de referência), para os modelos de produção (Binomial Negativa)
e da parte positiva da ocupação de UTI.}

\comoler{estes são os quatro modelos restantes, cujos valores de
aderência o texto acima cita: faturamento e tempo de permanência na
especificação de efeitos fixos, ocupação de internação e a versão
ajustada da mortalidade. Em todos, os pontos acompanham a diagonal na
maior parte do trajeto, descolando apenas nas caudas.}

\begin{figure}[H]
\centering
\subfig{fig_est_04_qq_custo_fe.png}{Faturamento real por saída, efeitos
fixos}
\subfig{fig_est_04_qq_tmp_fe.png}{Tempo de permanência, efeitos fixos}
\caption{Resíduos quantílicos dos dois modelos em logaritmo com efeitos
fixos de hospital.}
\label{fig:qqterceira}
\end{figure}

\comofeito{mesma construção dos resíduos quantílicos anteriores (quantis
teóricos da Normal no eixo horizontal, resíduos ordenados no vertical,
diagonal de referência), para os modelos de faturamento real por saída e
de tempo de permanência, ambos estimados com efeitos fixos de hospital.}

\begin{figure}[H]
\centering
\subfig{fig_est_04_qq_ocup_internacao.png}{Ocupação de internação}
\subfig{fig_est_04_qq_mort_ajustada_zib.png}{Mortalidade ajustada,
modelo Beta com inflação em zero}
\caption{Resíduos quantílicos da ocupação de internação e da versão
ajustada da mortalidade.}
\label{fig:qqquarta}
\end{figure}

\comofeito{mesma construção, para o modelo de ocupação de internação e
para a versão da mortalidade que exclui os óbitos fetais, maternos e
neonatais --- esta última é a que sustenta a verificação de robustez da
Seção~\ref{sec:robustez}.}

\comoler{compare com o painel (b) da figura anterior: os pontos que
antes despencavam da linha tracejada na ponta esquerda agora a
acompanham em quase todo o trajeto.}

\begin{figure}[H]
\centering
\includegraphics[width=.5\textwidth]{fig_est_08_qq_uti_gama.png}
\caption{Resíduos quantílicos da intensidade de UTI reestimada com a
família Gama: a aderência sobe de 0,814 para 0,944.}
\label{fig:qqgama}
\end{figure}

\comofeito{mesma construção do resíduo quantílico (quantis teóricos da
Normal no eixo horizontal, resíduos ordenados no vertical), para a
intensidade de ocupação de UTI reestimada com a família Gama; a aderência
sobe de 0,814 para 0,944 em relação ao modelo em logaritmo.}



\end{document}
